% concatenated from main.tex, format/defs.tex, 0-abstract.tex, 1-intro.tex, 2-related-work.tex, 3-main-result.tex, 4-setting.tex, 5-analysis.tex, 6-discussion.tex, 7-experiments.tex, 8-conclusions.tex, A1-SF-proof.tex, A2-lower-bound.tex, A3-bounding-qstar.tex, A4-other-guarantees.tex, A6-higherdy.tex, A8-symmetry-discretized.tex, A9-nonlinear-with-error.tex
\pdfoutput=1
\documentclass[11pt]{article}
\usepackage{microtype}
\usepackage{graphicx}
\usepackage{natbib}
% \usepackage{subfigure}
\usepackage{booktabs}
\usepackage{hyperref}
\usepackage{xcolor}
\usepackage[letterpaper, margin=1in]{geometry}
\usepackage{float}
\linespread{1.0}
\usepackage{booktabs}
\usepackage{colortbl}
\usepackage{xcolor}
\usepackage{adjustbox} % For resizing the table if needed
\usepackage{enumitem}  
\newcommand{\theHalgorithm}{\arabic{algorithm}}
\usepackage{amssymb,amsmath,amsthm, bm, amsfonts}
\usepackage{empheq}
\usepackage{algorithm,algorithmic}
\usepackage{subcaption}
\usepackage{longtable}
\usepackage{tikz}
\usepackage{fancyvrb}
\VerbatimFootnotes
\usetikzlibrary{shapes.geometric, arrows.meta, positioning, shadows}
\usepackage{times}
\usepackage{float}


\newcommand{\ABound}{\kappa_A}
\newcommand{\BBound}{\kappa_B}
\newcommand{\PertBound}{W}
\newcommand{\Diameter}{D}
\newcommand{\CostBound}{\beta}
\newcommand{\CostGradBound}{G}
\newcommand{\hor}{H}
\newcommand{\hortwo}{H_2}
\newcommand{\xidealbound}{D_X}
\newcommand{\specradbound}{1 - \gamma}
\newcommand{\affop}{\Psi}
\newcommand{\BigBound}{D}
\newcommand{\dimn}{d}
\renewcommand{\subset}{\subseteq}

\tikzstyle{box} = [rectangle, draw, fill=blue!10, text width=8em, text centered, rounded corners, minimum height=4em]
\tikzstyle{koopman-space} = [ellipse, draw, fill=red!10, text width=10em, text centered, minimum height=8em, dashed]
\tikzstyle{arrow} = [thick,->,>=Stealth]
\tikzstyle{dashed-arrow} = [thick,->,>=Stealth, dashed]


\newif\ifanonymous
%\anonymoustrue % Set to \anonymousfalse %for the camera-ready version


%\def\pinv{\dag}
%\def\ltwo#1{\left\|{#1}\right\|_2}
%\def\sqrtm#1{{#1}^{1/2}}
\def\half{\frac{1}{2}}
%\def\outerproduct#1#2{{#1}{#2}^{\!\!\!\top}}
%\def\xdomain{\mathcal{X}}
\DeclareMathOperator{\Tr}{\mathrm{Tr}}
\DeclareMathOperator*{\argmin}{argmin}
%%%%%% Author Notes %%%%%%%%
\newcommand{\eh}[1]{$\ll$\textsf{\color{blue} EH : #1}$\gg$}
\newcommand{\anand}[1]{$\ll$\textsf{\color{green} Anand : #1}$\gg$}

\newcommand{\evannote}[1]{\textcolor{red}{\{evan: #1\}}}


\def\A{{\mathcal A}}
\def\M{{\mathcal M}}
\def\X{{\mathcal X}}
\def\H{{\mathcal H}}
\def\Y{{\mathcal Y}}
\def\D{{\mathcal D}}
\def\K{{\mathcal K}}
\def\L{{\mathcal L}}
\def\F{{\mathcal F}}
\def\mP{{\mathcal P}}
\def\reals{{\mathbb R}}
\def\R{{\mathcal R}}
\def\fhat{{\ensuremath{\hat{f}}}}
\def\mH{\mathcal{H}}
\newcommand{\at}{\makeatletter @\makeatother}

\DeclareMathOperator{\clip}{clip}
% \def\norm#1{\mathopen\| #1 \mathclose\|}% use instead of $\|x\|$
\newcommand{\norm}[1]{\left\lVert#1\right\rVert}

\newcommand{\proj}{\mathop{\prod}}

\newcommand{\err}{\mathop{\mbox{\rm error}}}
\newcommand{\rank}{\mathop{\mbox{\rm rank}}}
\newcommand{\ellipsoid}{{\mathcal E}}
\newcommand{\sround}{\mathop{\mbox{\rm SmartRound}}}
\newcommand{\round}{\mathop{\mbox{\rm round}}}
\newcommand{\myspan}{\mathop{\mbox{\rm span}}}
\newcommand{\edges}{\mathop{\mbox{\rm edges}}}
\newcommand{\conv}{\mathop{\mbox{\rm conv}}}
\newcommand{\sign}{\mathop{\mbox{\rm sign}}}

\newcommand{\ignore}[1]{}
\newcommand{\equaldef}{\stackrel{\text{\tiny def}}{=}}
\newcommand{\equaltri}{\triangleq}
\newcommand{\tr}{\ensuremath{{\scriptscriptstyle\mathsf{T}}}}
\def\trace{{\bf Tr}}
\def\dist{{\bf P}}
\def\mydist{{\bf Q}}
\def\reals{{\mathbb R}}
\newcommand{\E}{\mathop{\mbox{\bf E}}}
\newcommand\ball{\mathbb{B}}
\newcommand\var{\mbox{VAR}}
\def\bzero{\mathbf{0}}

\def\Exp{\mathbb{\mathbb{E}}}
\def\mP{{\mathcal P}}
\def\risk{{\mbox{Risk}}}
\def\ent{{\mbox{ent}}}
\def\mG{{\mathcal G}}
\def\mA{{\mathcal A}}
\def\V{{\mathcal V}}
\def\lhat{\hat{\ell}}
\def\Lhat{\hat{L}}
\def\ltil{\tilde{\ell}}
\def\Ltil{\tilde{L}}

\def\Ahat{\hat{\mathbf{A}}}

\def\bold0{\mathbf{0}}
\def\boldone{\mathbf{1}}
\def\boldp{\mathbf{p}}
\def\boldq{\mathbf{q}}
\def\boldr{\mathbf{r}}
\def\boldA{\mathbf{A}}
\def\boldI{\mathbf{I}}
\newcommand\mycases[4] {{
\left\{
\begin{array}{ll}
    {#1} & {#2} \\\\
    {#3} & {#4}
\end{array}
\right. }}
\newcommand\mythreecases[6] {{
\left\{
\begin{array}{ll}
    {#1} & {#2} \\\\
    {#3} & {#4} \\\\
    {#5} & {#6}
\end{array}
\right. }}
\def\ba{\mathbf{a}}
\def\bb{\mathbf{b}}
\def\bc{\mathbf{c}}
\def\bB{\mathbf{B}}
\def\bC{\mathbf{C}}
\def\be{\mathbf{e}}
\def\etil{\tilde{\mathbf{e}}}
\def\bh{\mathbf{h}}
\def\bk{\mathbf{k}}
\def\bu{\mathbf{u}}
\def\bx{\mathbf{x}}
\def\w{\mathbf{w}}
\def\by{\mathbf{y}}
\def\bz{\mathbf{z}}
\def\bp{\mathbf{p}}
\def\bq{\mathbf{q}}
\def\br{\mathbf{r}}
\def\bu{\mathbf{u}}
\def\bv{\mathbf{v}}

\def\bw{\mathbf{w}}
\def\ba{\mathbf{a}}
\def\bm{\mathbf{m}}
\def\bA{\mathbf{A}}
\def\bM{\mathbf{M}}
\def\bS{\mathbf{S}}
\def\bG{\mathbf{G}}
\def\bK{\mathbf{K}}
\def\bL{\mathbf{L}}
\def\bH{\mathbf{H}}
\def\bI{\mathbf{I}}
\def\bJ{\mathbf{J}}
\def\bP{\mathbf{P}}
\def\bQ{\mathbf{Q}}
\def\bV{\mathbf{V}}
\def\bone{\mathbf{1}}


\def\xhat{\hat{\mathbf{x}}}
\def\xbar{\bar{\mathbf{x}}}
\def\vol{\mbox{vol}}
\def\trace{{\bf Tr}}
\def\Exp{\mathbb{\mathbb{E}}}
\newcommand{\val}{\mathop{\mbox{\rm val}}}
\renewcommand{\deg}{\mathop{\mbox{\rm deg}}}
\newcommand{\eps}{\varepsilon}

% Macros for this chapter
\def\bone{\mathbf{1}}
\newcommand{\sphere}{\ensuremath{\mathbb {S}}}
\newcommand{\simplex}{\ensuremath{\Delta}}
\newcommand{\diag}{\mbox{diag}}
%\newcommand{\K}{\ensuremath{\mathcal K}}
\newcommand{\Ktil}{\ensuremath{\mathcal{\tilde{K}}}}
\def\mA{{\mathcal A}}
\newcommand{\x}{\ensuremath{\mathbf x}}
\newcommand{\vn}{\ensuremath{\mathbf n}}
\newcommand{\rv}[1][t]{\ensuremath{\mathbf r_{#1}}}
\newcommand{\y}{\ensuremath{\mathbf y}}
\newcommand{\z}{\ensuremath{\mathbf z}}
\newcommand{\h}{\ensuremath{\mathbf h}}
\newcommand{\xtil}[1][t]{\ensuremath{\mathbf {\tilde{x}}_{#1}}}
\newcommand{\xv}[1][t]{\ensuremath{\mathbf x_{#1}}}
\newcommand{\xvh}[1][t]{\ensuremath{\mathbf {\hat x}}_{#1}}
\newcommand{\ftildeh}[1][t]{\ensuremath{\mathbf {\hat f}}_{#1}}
\newcommand{\xvbar}[1][t]{\ensuremath{\bar{\mathbf x}_{#1}}}
\newcommand{\xstar}{\ensuremath{\mathbf x^{*}}}
\newcommand{\xstarr}{\ensuremath{\mathbf x^{\star}}}
\newcommand{\yv}[1][t]{\ensuremath{\mathbf y_{#1}}}
\newcommand{\zv}[1][t]{\ensuremath{\mathbf z_{#1}}}
\newcommand{\fv}[1][t]{\ensuremath{\mathbf f_{#1}}}
\newcommand{\gv}[1][t]{\ensuremath{\mathbf g_{#1}}}
\newcommand{\qv}[1][t]{\ensuremath{\mathbf q_{#1}}}
\newcommand{\pv}[1][t]{\ensuremath{\mathbf p_{#1}}}
\newcommand{\ftilde}[1][t]{\ensuremath{\mathbf {\tilde f}_{#1}}}
\newcommand{\gtilde}[1][t]{\ensuremath{\mathbf {\tilde g}_{#1}}}
\newcommand{\ev}[1][i]{\ensuremath{\mathbf e_{#1}}}
\newcommand{\uv}{\ensuremath{\mathbf u}}
\newcommand{\vv}{\ensuremath{\mathbf v}}
%\newcommand{\A}[1][t]{\ensuremath{\mathbf A_{#1}}}
\newcommand{\Ftilde}[1][t]{\ensuremath{\mathbf {\tilde F}_{#1}}}
\newcommand{\Gtilde}[1][t]{\ensuremath{\mathbf {\tilde G}_{#1}}}
\newcommand{\hessum}[1][t]{\ensuremath{\boldsymbol{H}_{1:#1}}}
\newcommand{\reg}[1][t]{\ensuremath{\lambda_{#1}}}
\newcommand{\regsum}[1][t]{\ensuremath{\boldsymbol{\lambda}_{1:#1}}}
%\newcommand{\F}{\ensuremath{\mathcal F}}
\newcommand{\ntil}{\tilde{\nabla}}
\newcommand{\Ftil}{\tilde{F}}
\newcommand{\ftil}{\tilde{f}}
\newcommand{\rhotil}{\tilde{\rho}}
\newcommand{\gtil}{\tilde{g}}
\newcommand{\mutil}{\tilde{\mu}}
\newcommand{\Atil}{\tilde{A}}
\newcommand{\Qtil}{\tilde{Q}}
\newcommand{\vtil}{\tilde{v}}
\newcommand\super[1] {{^{(#1)}}}
\newcommand\pr{\mbox{\bf Pr}}
\newcommand\av{\mbox{\bf E}}
\def\regret{\ensuremath{\mbox{Regret}}}
\def\bb{\mathbf{b}}
%\def\bB{\mathbf{B}}
%\def\bB{{B}}
%\def\bC{\mathbf{C}}
\def\bC{{C}}
\def\be{\mathbf{e}}
\def\bu{\mathbf{u}}
\def\bx{\mathbf{x}}
\def\by{\mathbf{y}}
\def\bz{\mathbf{z}}
\def\bp{\mathbf{p}}
\def\bdelta{\mathbf{\delta}}
\def\bq{\mathbf{q}}
\def\br{\mathbf{r}}
\def\bv{\mathbf{v}}
\def\ba{\mathbf{a}}
\def\bg{\mathbf{g}}
\def\bs{\mathbf{s}}
\newcommand{\dm}{\delta_{\times}}
\newcommand{\cA}{\mathcal{A}}
%\def\bA{\mathbf{A}}
%\def\bA{{A}}
\def\bI{\mathbf{I}}
\def\bJ{\mathbf{J}}
\def\bP{\mathbf{P}}
\def\bV{\mathbf{V}}
\def\sgn{\text{sgn}}
\def\eps{\varepsilon}
\def\epsilon{\varepsilon}
\def\veps{\varepsilon}
\newcommand{\etal}{ \textit{et al.}\xspace}
\def\tsum{{\textstyle \sum}}
\def\universal{{\sc Universal}\xspace}
\def\ogd{{\sc Online Gradient Descent}\xspace}
\def\R{\ensuremath{\mathcal R}}
\def\lmid{\lambda_{\text{mid}}}

\def\hestinv{\tilde{\nabla}^{-2}}
\def\hest{\tilde{\nabla}^{2}}
\def\hessinv{\nabla^{-2}}
\def\hess{\nabla^2}
\def\grad{\nabla}
\newcommand{\defeq}{\triangleq}

\newcommand*\xor{\mathbin{\oplus}}
\newcommand{\innerprod}[2]{\langle #1, #2 \rangle}
\newcommand{\matnorm}[2]{\| #1 \|_{#2}}
\newcommand{\invmatnorm}[2]{\| #1 \|_{#2^{-1}}}


%specific to non convex second order paper
\newcommand{\hlinverse}{(\bH + \lambda \bI)^{-1}}
\newcommand{\hlpinverse}{(\bH + \lambda \bI)^{+}}
\newcommand{\hlinversea}[1]{(\bH + #1 \bI)^{-1}}
\newcommand{\hlpinversea}[1]{(\bH + #1 \bI)^{+}}
\newcommand{\ceil}[1]{\left\lceil #1 \right\rceil}

\newcommand{\Tg}{\mathbb{T}_g}
\newcommand{\Th}{\mathbb{T}_h}
\newcommand{\Thh}{\mathbb{T}_{h,1}}
\newcommand{\Time}{\mathbb{T}_{\mathsf{inverse}}}

\newtheorem{innercustomthm}{Theorem}
\newenvironment{stheorem}[1]
  {\renewcommand\theinnercustomthm{#1}\innercustomthm}
  {\endinnercustomthm}
\newtheorem{innercustomlemma}{Lemma}
\newenvironment{slemma}[1]
  {\renewcommand\theinnercustomlemma{#1}\innercustomlemma}
  {\endinnercustomlemma}
%
\newtheorem{ctheorem}{Theorem}
\newtheorem{clemma}{Main Lemma}

\newtheorem{theorem}{Theorem}[section]
%\newtheorem*{theorem*}{Theorem}

\newtheorem{claim}[theorem]{Claim}
\newtheorem{subclaim}{Claim}[theorem]
\newtheorem{proposition}[theorem]{Proposition}
\newtheorem{lemma}[theorem]{Lemma}
\newtheorem{corollary}[theorem]{Corollary}
\newtheorem{conjecture}[theorem]{Conjecture}
\newtheorem{observation}[theorem]{Observation}
\newtheorem{fact}[theorem]{Fact}


\newtheorem{definition}[theorem]{Definition}
\newtheorem{construction}[theorem]{Construction}
\newtheorem{example}[theorem]{Example}
\newtheorem{algorithm1}[theorem]{Algorithm}
\newtheorem{protocol}[theorem]{Protocol}
\newtheorem{assumption}[theorem]{Assumption}

%\newtheorem*{rep@theorem}{\rep@title}
\newcommand{\newreptheorem}[2]{%
\newenvironment{rep#1}[1]{%
 \def\rep@title{#2 \ref{##1}}%
 \begin{rep@theorem}}%
 {\end{rep@theorem}}}


\newreptheorem{theorem}{Theorem}
\newreptheorem{lemma}{Lemma}
\newreptheorem{proposition}{Proposition}
\newreptheorem{claim}{Claim}
\newreptheorem{corollary}{Corollary}
\newreptheorem{mainlemma}{Main Lemma}


\newtheorem{remark}[theorem]{Remark}

%%%%%%%%%%%%%%%%%%%%%%%%%%%%%%%%%%%%%%%%%%%%%%%%%%%%%%%%%%%%%%%%%%%%%%%%%%%%%%%%
% REFERNCE MACROS

\newcommand{\namedref}[2]{\mbox{\hyperref[#2]{#1~\ref*{#2}}}}

\newcommand{\chapterref}[1]{\namedref{Chapter}{#1}}
\newcommand{\sectionref}[1]{\namedref{Section}{#1}}
\newcommand{\appendixref}[1]{\namedref{Appendix}{#1}}
\newcommand{\theoremref}[1]{\namedref{Theorem}{#1}}
\newcommand{\factref}[1]{\namedref{Fact}{#1}}
\newcommand{\remarkref}[1]{\namedref{Remark}{#1}}
\newcommand{\definitionref}[1]{\namedref{Definition}{#1}}
\newcommand{\figureref}[1]{\namedref{Figure}{#1}}
\newcommand{\figurerefb}[2]{\mbox{\hyperref[#1]{Figure~\ref*{#1}#2}}}
\newcommand{\tableref}[1]{\namedref{Table}{#1}}
\newcommand{\lemmaref}[1]{\namedref{Lemma}{#1}}
\newcommand{\mainlemmaref}[1]{\namedref{Main Lemma}{#1}}
\newcommand{\claimref}[1]{\namedref{Claim}{#1}}
\newcommand{\propositionref}[1]{\namedref{Proposition}{#1}}
\newcommand{\corollaryref}[1]{\namedref{Corollary}{#1}}
\newcommand{\constructionref}[1]{\namedref{Construction}{#1}}
\newcommand{\itemref}[1]{\namedref{Item}{#1}}
\newcommand{\propertyref}[1]{\namedref{Property}{#1}}
\newcommand{\protocolref}[1]{\namedref{Protocol}{#1}}
\newcommand{\algorithmref}[1]{\namedref{Algorithm}{#1}}
\newcommand{\assumptionref}[1]{\namedref{Assumption}{#1}}
\newcommand{\stepref}[1]{\namedref{Step}{#1}}
\newcommand{\paramref}[1]{\namedref{Parameter}{#1}}
\newcommand{\footnoteref}[1]{\namedref{Footnote}{#1}}
\newcommand{\equationref}[1]{\mbox{\hyperref[#1]{(\ref*{#1})}}}
\renewcommand{\eqref}{\equationref}
%\newcommand{\equationref}[1]{\namedref{Equation}{#1}}
\newcommand{\exampleref}[1]{\namedref{Example}{#1}}
%\newcommand{\lineref}[1]{\namedref{Line}{#1}}
%%%%%%%%%%%%%%%%%%%%%%%%%%%%%%%%%%%%%%%%%%%%%%%%%%%%%%%%%%%%%%%%%%%%%%%%%%%%%%%%

\numberwithin{equation}{section}


\newcommand{\da}{\text{\ding{172}}\xspace}
\newcommand{\db}{\text{\ding{173}}\xspace}
\newcommand{\dc}{\text{\ding{174}}\xspace}
\newcommand{\dd}{\text{\ding{175}}\xspace}
\newcommand{\de}{\text{\ding{176}}\xspace}
\newcommand{\df}{\text{\ding{177}}\xspace}
\newcommand{\dg}{\text{\ding{178}}\xspace}

% brackets
\newcommand{\braces}[1]{\left\{#1\right\}}
\newcommand{\pa}[1]{\left(#1\right)}
\newcommand{\ang}[1]{\left<#1\right>}
\newcommand{\bra}[1]{\left[#1\right]}
\newcommand{\abs}[1]{\left|#1\right|}

% Brian macros %
\def\memdiam{D}
\def\memgradbound{G_f}

\newcommand{\Return}{\textbf{return} }
\begin{document}
\title{Universal Learning of Nonlinear Dynamics}


%\ifanonymous
%\author{Anonymous Authors}
%\elseif
\author{
  Evan Dogariu\thanks{Courant Institute of Mathematical Sciences, New York University} \and
  Anand Brahmbhatt\thanks{Princeton University } \and
  Elad Hazan\footnotemark[2]
  %\thanks{Princeton University }
}
%\fi
\maketitle


\begin{abstract}\label{sec:0-abstract}

\ignore{
We introduce a new approach to learning dynamical systems based on \emph{improper learning} and spectral filtering. 
Instead of identifying the true dynamics, our method competes with a tractable comparator class: the best high-dimensional linear observer for the observed data. 
This shift avoids the nonconvexity and ill-conditioning of system identification while preserving finite-time performance guarantees, even for marginally stable, noisy, and nonlinear systems.

Our main algorithm, \emph{Online Spectral Filtering} (OSF), is an efficient online convex optimization method that provably learns any marginally stable, observable nonlinear dynamical system with regret bounded by $Q^\star \sqrt{T}$, where $Q^\star$ is a control-theoretic condition number defined through a Luenberger observer program. 
At its core is a new spectral filtering technique for linear dynamical systems that, for the first time, handles general asymmetric dynamics with adversarial process noise and without dependence on hidden state dimension. 
We also present a constructive discretization-based lifting from bounded Lipschitz nonlinear systems to high-dimensional linear systems, providing a principled way to extend OSF beyond the linear setting.

Experiments on synthetic systems, the Lorenz attractor, and Brax physics environments confirm the theoretical predictions, showing that OSF achieves high accuracy, robustness to noise, and the expected scaling with $Q^\star$, and outperforms strong baselines including eDMD and direct observer learning. 
This work bridges tools from control theory and online learning, illustrating how spectral filtering over a well-chosen comparator class can yield algorithms that are both provable and practical for complex dynamical systems.
}


We study the fundamental problem of learning a marginally stable unknown nonlinear dynamical system. We describe an algorithm for this problem, based on the technique of spectral filtering, which learns a mapping from past observations to the next based on a spectral representation of the system. Using techniques from online convex optimization, we prove vanishing prediction error for any nonlinear dynamical system that has finitely many marginally stable modes, with rates governed by a novel quantitative control-theoretic notion of learnability. The main technical component of our method is a new spectral filtering algorithm for linear dynamical systems, which incorporates past observations and applies to general noisy and marginally stable systems. This significantly generalizes the original spectral filtering algorithm to both asymmetric dynamics as well as incorporating noise correction, and is of independent interest.


%Spectral methods have been shown to perform well in predicting sequences generated by linear dynamical systems. In this work, we extend these results to sequences generated by non-linear dynamical systems through the lens of Koopman operator theory. We introduce a novel \emph{Spectral Transformer} unit, which we theoretically analyze under specific assumptions and demonstrate to be effective in practice through empirical evaluation.
\end{abstract}

% \tableofcontents

\section{Introduction}\label{sec:1-introduction}


\renewcommand{\R}{\mathbb{R}}

The problem of learning a dynamical system from observations is a cornerstone of scientific and engineering inquiry. Across disciplines ranging from econometrics and meteorology to robotics and neuroscience, a central task is to predict the future evolution of a system given only a sequence of its past measurements. Formally, given a sequence of observations  $(y_1, ...,y_t)$ generated by an unknown, underlying
dynamical system\footnote{For ease of presentation, we keep the discussion deterministic in the main paper and discuss an extension to stochastic systems in Appendix \ref{sec:nl-w-noise}.} of the form 
\begin{equation}
    \label{eqn:NLDS}
    x_{t+1} = f(x_t), \quad y_t = h(x_t) ,
\end{equation}
the objective is to produce an accurate forecast $\hat{y}_{t+1}$. The challenge lies in the fact that the state $x_t$ is hidden and may be high-dimensional, and the dynamics function $f: \mathcal{X} \to \mathcal{X}$ and observation function $h: \mathcal{X} \to \mathcal{Y}$ are unknown.

Historically, two dominant paradigms have emerged to address this challenge. The first, rooted in control theory and system identification, aims to construct an explicit model of the underlying dynamics. This model-based approach seeks to learn the function $f$ or, more commonly, a linear representation of it. A powerful theoretical tool in this domain is the Koopman operator, which recasts the nonlinear dynamics of states into an infinite-dimensional, yet linear, evolution of observable functions. Data-driven methods such as Dynamic Mode Decomposition (DMD) attempt to find finite-dimensional approximations of this linear operator, yielding a model that can be used for prediction and control \cite{brunton2021modernkoopmantheorydynamical}. While elegant, these methods often rely on strong assumptions to guarantee the convergence of their approximations and can be computationally demanding. 

The second paradigm, driven by recent advances in machine learning, employs powerful, general-purpose black-box architectures to learn the mapping from past to future observations directly. Architectures such as Transformers \cite{vaswani2017attention} and more recent state-space models like Mamba \cite{gu2023mamba} and convolutional models like Hyena \cite{poli2023hyena} have demonstrated remarkable success in sequence modeling tasks. These methods are highly expressive and often achieve state-of-the-art performance, but their operation can be opaque, and they typically lack the formal performance guarantees characteristic of control-theoretic approaches, especially in the presence of adversarial disturbances or for systems operating at the edge of stability.

\paragraph{Provable improper learning for dynamical systems.}
This paper advocates for a third paradigm: improper learning for dynamical systems. The philosophy of this approach is that for the specific goal of prediction, it is not necessary to recover the true, and often intractably complex, underlying dynamics. Instead, one can design an efficient and provably correct algorithm by framing the learning problem as a competition. The algorithm does not attempt to learn a model from the same class as the true system (hence, "improper"); rather, it competes in hindsight against a class of idealized, well-behaved, and computationally tractable predictors. 

This reframing of the problem is a key conceptual leap. It shifts the focus from the difficult, nonconvex problem of system identification to the more manageable, often convex, problem of regret minimization against a carefully chosen comparator class. The central idea of this work is to leverage the rich mathematical structure of control theory not for designing a controller, but for defining this very comparator class. Instead of attempting to learn the true nonlinear system, we design an algorithm that is guaranteed to perform as well as the best possible high-dimensional linear observer system for the given observation sequence. This approach allows us to sidestep the complexities of nonlinear system identification while retaining strong theoretical guarantees. Furthermore, our analysis yields a natural quantitative description of the complexity of learning from observations with improper algorithms.

\paragraph{Our contributions.}
This work makes several contributions to the theory and practice of learning dynamical systems, bridging concepts from classical control theory and modern machine learning.
\begin{enumerate}
    \item 
We introduce \textbf{Observation Spectral Filtering} (OSF), an efficient algorithm based on online convex optimization that is guaranteed to learn any observable nonlinear dynamical system with finitely many marginally stable modes. The algorithm is "improper" in that it does not explicitly estimate the system's hidden states or dynamics; instead, it directly constructs a mapping from past observations to future predictions. 

\item 
\textbf{An LDS guarantee:} The core technical engine of OSF is a learning guarantee for linear dynamical systems (LDS). This method applies to asymmetric dynamics under adversarial noise with performance guarantees that are independent of the system's hidden dimension. This significantly generalizes prior spectral filtering methods, which were largely restricted to symmetric systems, and aligns our work with recent efforts to extend these powerful techniques to broader system classes. 

\item
\textbf{A control-theoretic analysis:} We provide a novel analytical framework that connects the learnability of a nonlinear system to the properties of its best possible high-dimensional linear observer. This connection is formalized through a "Luenberger program," an optimization problem whose solution quantifies the inherent difficulty of learning the system from its outputs. The optimal value of this program, a condition number we denote $Q_\star$, appears directly in our data-dependent learning bounds. 
This framework establishes a quantitative link between the control-theoretic concept of observability and the machine learning concept of learnability, and provides an elegant analytic toolkit with which to derive learning guarantees in nonlinear settings using linear methods. We discuss further implications in Section \ref{sec:discussion}.

\item
\textbf{A simple global linearization:} We introduce a simple construction to approximate any bounded, Lipschitz nonlinear system with a high-dimensional LDS. This technique, based on state-space discretization, provides a rigorous foundation for applying our linear systems algorithm to the nonlinear case. It trades an increase in dimensionality—to which our learning algorithm is immune—for guaranteed linearity and approximation accuracy, and it uses the improper nature of our analysis to circumvent the strong spectral assumptions and convergence difficulties associated with many contemporary data-driven Koopman operator methods. 
\end{enumerate}

Ultimately, this work presents a synthesis of ideas from two distinct fields. It repurposes tools from control theory (such as the Luenberger observer, pole placement, and global linearization) not for their traditional purpose of observer/controller design, but as analytical instruments to define a robust comparator class for a learning algorithm.  By then applying the machinery of online convex optimization, we develop an algorithm whose performance guarantees are directly informed by the structural properties of the underlying prediction problem.  This combination of a constructive global linearization and  spectral filtering yields a universal learning algorithm for any observable nonlinear dynamical system with finitely many marginally stable modes, with regret scaling determined only by the system’s robust observability constant $Q_\star$.


\ignore{
\newpage

The fundamental learning problem in the field of control for dynamical systems is learning a nonlinear dynamical system from observations. Let  $\mathcal{X} \subseteq \R^{d_\mathcal{X}}$ be a state space, and $\mathcal{Y} \subseteq \R^{d_{\mathcal{Y}}}$ be an observation space. A (deterministic or stochastic) dynamical system is described by the equations 
\begin{align}\label{eqn:NLDS}
x_{t+1} = f(x_t) ,\quad  y_{t+1} = h(x_{t+1}) , \end{align}
where the dynamics function is $f: \mathcal{X} \to \mathcal{X}$, the observation function is $h: \mathcal{X} \to \mathcal{Y}$, and $x_0 \in \mathcal{X}$ is the initial state of the system. For simplicity of presentation, we restrict the introduction to deterministic dynamics, but our results generalize to stochastic systems and are detailed henceforth. 
The goal is to accurately predict $y_{t+1}$ given the sequence of observations $(y_0, \ldots, y_t)$,  without knowing the dynamics. For example, if we chose the mean square error as a metric, our goal is to predict a sequence of $\hat{y}_t$, such that as $T \to \infty$ the average loss decays:
$$  \frac{1}{T} \sum_{t =0}^T \norm{\hat{y}_t - y_t}^2 \to  0 .$$
An algorithm that achieves this is said to learn the dynamical system. 
Numerous methods have been proposed for learning dynamical systems, either by attempting to learn the dynamics function itself or by learning a high-dimensional linear representation of the dynamics, which may be constructed using the Koopman operator of the system. 

In contrast, our main result is an efficient algorithm which is {\it improper}, i.e. it does not try to recover what the dynamics function $f$ or the hidden states $x_t$ are. Instead, it constructs a mapping from past observations to future predictions while guaranteeing learning the system. An informal statement of our main result is as follows:

\begin{theorem}[Informal main theorem, see Theorem \ref{theorem:main_nonlinear}]
Let $(f,h, x_0)$ be a nonlinear dynamical system which has finitely-many marginally stable modes and is observable. Then, 
spectral filtering over past observations $y_1, \ldots, y_t$ produces a sequence of predictions $\hat{y}_{t+1}$ such that $  \frac{1}{T} \sum_{t =0}^T \norm{\hat{y}_t - y_t}^2 \to 0 $.
\end{theorem}
The main new technical tool used in our result is a new variant of spectral filtering for linear dynamical systems, stated in Theorem \ref{theorem:main_linear}. 
}






\section{Related Work}

In this section, we survey the past and current methods for sequence prediction in dynamical systems, highlighting their relationships, strengths, and limitations.
While it is impossible to provide an exhaustive account, we highlight the works most relevant to our study.

\subsection{Methods for Learning in Nonlinear Systems}

Learning in nonlinear dynamical systems of the form (\ref{eqn:NLDS}), also called time series prediction in statistics, spans several fields of study. We can roughly divide the existing methods as follows:

\paragraph{Statistical and online learning methods.}

The classical text of Box and Jenkins \cite{box1976time} introduces ARMA (autoregressive moving average) statistical estimation techniques and their extensions, notably ARIMA. These methods are ubiquitous in applications and were extended to the adversarial online learning domain in \cite{anava2013online,kuznetsov2018theory} via online convex optimization \cite{hazan2016introduction}. Later extensions include nonlinear autoregressive methods augmented with kernels, e.g. \cite{cuturi2011autoregressive}.  
Also prominent in the statistical literature are nonparameteric time series methods such as Gaussian processes \cite{rasmussen2003gaussian,murphy2012machine}.


\paragraph{Deep learning methods.}

The most successful methods for learning dynamical systems are based on deep learning. Transformers \cite{vaswani2017attention} are the most widely used architectures that form the basis for large language models. More recent deep architectures are based on convolutional models such as Hyena \cite{poli2023hyena} and spectral architectures such as FlashSTU \cite{liu2024flash}.
Recurrent neural networks and state space models such as Mamba \cite{gu2023mamba} exhibit faster inference and require lower memory, albeit being harder to train and potentially less expressive. 

\paragraph{Koopman operator methods.}

The most relevant technique for learning nonlinear dynamical systems is via the Koopman operator. This methodology relies on the mathematical fact that any (partially observed) dynamical system can be lifted to a linear representation in an infinite-dimensional space. 
The basic idea is to represent the evolution of the system by describing the evolution of all possible system observables via the Koopman operator, which becomes linear and allows for functional-analytic and spectral methods. Detailed theoretical study of this operator can require complicated analysis and strong assumptions; a simple version of a related linearization can be achieved via discretization, which we present in Section \ref{sec:koopman-discrete}. 

Data-driven methods for Koopman operator learning are based on learning an approximate representation of the dynamics in high-dimensional linearly-evolving coordinates. Knowing the top $N$ eigenvalues and eigenfunctions of the Koopman operator allows one to approximately simulate, predict, and control the system using $N$-dimensional linear methods. These eigenfunctions may be learned from data via the use of SVD on the trajectories, a predetermined dictionary of reference functions, or a neural network, which inspires the classical DMD and eDMD algorithms. The linear coordinates are often learned jointly with the dynamics in this linear representation via regression. In short, this pipeline learns a lifting to a linear system on which system identification is performed.
See the survey \cite{brunton2021modernkoopmantheorydynamical} for more details about this approach. 


In contrast, our method is {\it improper}, i.e. it does not learn the mapping of the nonlinear dynamical system to higher dimensions explicitly. We rely on the spectral filtering method to automatically learn the best predictor that competes with a linear predictor with very high-dimensional hidden state. 



\subsection{Methods for Learning in Linear Systems}


The simplest and most commonly studied dynamical systems in the sciences and engineering are linear. 
Given input vectors $u_1, \dots, u_T \in \R^{d_{\textrm{in}}}$, the system generates a sequence of output vectors $y_1, \dots y_T \in \R^{d_{\textrm{out}}}$ according to the following equations 
\begin{align}
\label{eqn:LDS}
    {x}_{t+1} & = A {x}_t + B u_t + w_t \\
    y_t & = C {x}_t + {D} u_t + \xi_t, \notag
\end{align}
where ${x}_0, \dots, {x}_T \in \R^{d_h}$ is a sequence of hidden states, $(A, B, C, {D})$ are matrices which parameterize the LDS, and $w_t, \xi_t$ are perturbation terms. We assume w.l.o.g. that ${D},\xi_t = 0$, as these terms can be folded into the input of the previous iteration. The problem of {\it learning} in such systems refers to predicting the next observation $y_t$ given all previous inputs $u_{1:t}$ and past observations $y_{1:t-1}$. 


The fundamental problem of learning in linear dynamical systems has been studied for many decades. In complex systems the hidden dimension can be very large, and it is thus crucial to avoid explicit complexity in this parameter. Other sources of difficulty come from the spectral properties of the dynamics matrix $A$ such as its spectral radius, which determines how close it is to marginal stability. Asymmetry of $A$ and the presence of process noise $w$ can further complicate the situation. 
We briefly survey the extensive and decades-spanning literature on this topic, focusing on the main approaches or techniques. 

\paragraph{System identification.}
System identification refers to the method of recovering $A,B,C$ from the data. This is a hard nonconvex problem that is known to fail if the system is marginally stable, i.e. has eigenvalues that are close to one. Also, recovering the system means time and memory proportional to the hidden dimension, which can be extremely large. \cite{hardt2018gradient} shows that the matrices $A,B,C$ of an unknown and partially observed stable linear dynamical system can be learned using first-order optimization methods. 

In the statistical noise setting, a natural approach for learning is to perform system identification and then use the identified system for prediction. This approach was taken by \cite{simchowitz2018learning,pmlr-v119-simchowitz20a,pmlr-v97-sarkar19a}. The work of \cite{ghai2020no} extends identification-based techniques to adversarial noise and marginally stable systems. 
The work of \cite{simchowitz2019learning} presented a regression procedure for identifying Markov operators under nonstochastic (or even adversarial) noise. This approach has roots in the classical Ho-Kalman system identification procedure \cite{ho1966effective}, for which the first non-asymptotic sample complexity, although for stochastic noise, was furnished in \cite{oymak2019non}.  Nonstochastic control of unknown, unstable systems without any access to prior information, also called black-box control, was addressed in \cite{chen2021black}.





\paragraph{Regression a.k.a. convex relaxation. } The (auto) regression method predicts according to $\hat{y}_{t} = \sum_{i=0}^m M_i u_{t-i}$. The coefficients $M_i$ can be learned using convex regression, and without dependence on the hidden dimension. The downside of this approach is that if the spectral radius of $A$ is $1-\delta$, it can be seen that $m \approx \frac{1}{\delta}$ terms are needed. 
The regression method can be further enhanced with ``closed loop" components, that regress on prior observations $y_{t-1:1}$. It can be shown using the Cayley-Hamilton theorem that using this method, $d_h$ components are needed, see e.g.  \cite{kozdoba2019line,hazan2017learning}.


 
\paragraph{Method of moments.}

\cite{hazan2019nonstochastic} give a method-of-moments system identification procedure as well as the first end-to-end regret result for nonstochastic control starting with unknown system matrices.  Extension of these methods to tensors were used in \cite{bakshi2023new}, and in \cite{bakshi2023tensor} to learn a mixture of linear dynamical systems. 


\paragraph{Kalman filtering.}
Kalman Filtering \cite{kalman1960new,anderson1991kalman} involves recovering the state $x_t$ from observations. While Kalman filtering is optimal under specific noise conditions, it generally fails in the presence of marginal stability and adversarial noise. In addition, its complexity depends on the hidden dimension of the system.    



\paragraph{Spectral filtering.}
The technique of spectral filtering \cite{hazan2017learning} combines the advantages of many previous methods. It is an efficient method whose complexity does {\bf not} depend on the hidden dimension, and works even for marginally stable linear dynamical systems.  Various extensions of the basic technique were proposed for learning and control \cite{hazan2018spectral,arora2018towards,brahmbhatt2025newapproachcontrollinglinear}, however these do not apply in full generality. The most significant advancement was the recent result of   \cite{marsden2025universalsequencepreconditioning}, who gave an extension to certain asymmetric LDS, those whose imaginary eigenvalues are bounded from above. 
\newline \newline
The main technical contribution of this paper, from which we derive the guarantees for learning nonlinear dynamics, is a spectral filtering method which circumvents previous difficulties and obtains the best of all worlds -- an efficient method whose parametrization does not depend on the hidden dimension, and allows for marginal stability and asymmetry.  Furthermore, our method naturally allows handling of adversarial perturbations in the regret minimization sense. 
Table \ref{tab:my_label} below summarizes the state of the art in terms of the fundamental problem of learning linear dynamical systems, and how our result advances it. 


\begin{table}[h!]
    \centering
    \begin{tabular}{|c|c|c|c|c|}
    \hline
        Method & Marginally stable  & $d_{\text{hidden}}$-free  & asymmetric & adversarial noise \\
    \hline
    \hline
        Sys-ID  & $\times$ & $\times$ & $\checkmark$ & $\times $  \\
    \hline
        Regression, open-loop & $\times$ & $\checkmark$  &  $\checkmark$ & $\checkmark$  \\
    \hline
        Regression, closed-loop  & $\checkmark$ & $\times$ & $\checkmark$  & $\checkmark $ \\
    \hline
        Spectral filtering  \cite{hazan2017learning}& $\checkmark$ & $\checkmark$ & $\times$  & $\times $  \\
    \hline
        USP \cite{marsden2025universalsequencepreconditioning}   & $\checkmark$ & $\checkmark$ & $\sim \checkmark$  & $\times $  \\
    \hline
        OSF (ours)  & $\checkmark$ & $\checkmark$ & $\checkmark$  & $\checkmark $  \\
    \hline
    \end{tabular}
    \caption{Methods for learning LDS.}
    \label{tab:my_label}
\end{table}




\section{Algorithm and Main Result}
Our algorithm leverages spectral filtering of past outputs to build a predictor. By filtering over past observations $y_{t-1}, y_{t-2}, \ldots$, the algorithm can handle nonlinear dynamics, asymmetric linear dynamics, and achieves robustness to adversarial process noise; this is implemented\footnote{To avoid restating many variations of the same algorithm, we state Algorithm \ref{alg:clsf_regressionmain} over both open-loop controls $u_t$ and observations $y_t$. It is understood that when applying our results to a nonlinear system such as (\ref{eqn:NLDS}), since there are no $u_t$'s we take $J^t, M^t \equiv 0$.} in Algorithm \ref{alg:clsf_regressionmain}.  
Our proof techniques also yield learning guarantees for other improper learning algorithms, which we describe in Appendix \ref{appendix:otherguarantees} for completeness. 

 Henceforth, for simplicity of presentation we assume that the dimension of the observation is one, i.e. $d_\mathcal{Y}=d_{\mathrm{obs}}=1$; this is the hardest case to learn. The details that change when the observation space is larger are noted in Appendix \ref{sec:higher-d-y}. 


\begin{algorithm}[htb]
\caption{Observation Spectral Filtering + Regression}
\label{alg:clsf_regressionmain}
\begin{algorithmic}[1]
\STATE \textbf{Input:} Horizon \(T\), number of filters \(h\), number of autoregressive components $m$, step sizes $\eta_t$, convex constraint set \(\K\), prediction bound $R$.
\STATE Compute $\{(\sigma_j, {\phi}_j)\}_{j=1}^k$ the top eigenpairs of the $(T-1)$-dimensional Hankel matrix of \cite{hazan2017learning}.
\STATE Initialize $J_j^0, M_i^0  \in \R^{d_{out} \times d_{in}}$ and $P_j^0, N_i^0  \in \R^{d_{out} \times d_{out}}$ for $j \in \{1, \ldots, m\}$ and $i\in\{1,\dots,h\}$. Let $\Theta^t = (J^t, M^t, P^t, N^t)$ for notation.
\FOR{$t = 0, \ldots, T-1$}
\STATE Compute \begin{align*}\hat{y}_t(\Theta^t) &= \sum_{j=1}^{m-1}J_j^tu_{t-j} +  \sum_{i=1}^h \sigma_i^{1/4}M_i^t \langle \phi_{i} ,  u_{t-2:t-T} \rangle \\ &\;+ \sum_{j=1}^{m}P_j^ty_{t-j} + \sum_{i=1}^h \sigma_i^{1/4}N_i^t \langle \phi_{j} ,  y_{t-1:t-T-1} \rangle\,, \end{align*}
and if necessary project $\hat{y}_t(\Theta^t)$ to the ball of radius $R$ in $\mathcal{Y}$.
\STATE Compute loss $\ell_t(\Theta^t) = \norm{\hat{y}_t(\Theta^t) - y_t}$.
\STATE Update $\Theta^{t+1} \leftarrow \Pi_\K\left[\Theta^t - \eta_t  \nabla_\Theta \ell_t(\Theta^t)\right]$. 
\hfill \textcolor[HTML]{408080}{\# can be replaced with other optimizer}
\ENDFOR
\end{algorithmic}
\end{algorithm}

Theorem \ref{theorem:main_nonlinear} states the learning guarantee this algorithm achieves for nonlinear dynamical systems. 
We will use the construction of the Luenberger program, defined precisely in Definition \ref{def:nonlinearluenbergerprogram}, with an optimal value $Q_\star$ which depends on the dynamical system.

\begin{theorem}[Nonlinear]
    \label{theorem:main_nonlinear}
    Let $y_1, \ldots, y_T$ be the observations of a nonlinear dynamical system satisfying Assumption \ref{assumption:nonlinear}. Fix the number of algorithm parameters $h = \Theta(\log^2T)$ and $m \geq 1$, and for notation define $\gamma = \frac{\log(mT)}{m}$. 
    
  Let $Q_\star = Q_\star\left(f,h,\frac{1}{T^{3/2}}, \left[0, 1 - \frac{1}{\sqrt{T}}\right] \cup \{1\} \cup \mathbb{D}_{1-\gamma}\right)$ as specified in Definition \ref{def:nonlinearluenbergerprogram}.
    Then, running Algorithm \ref{alg:clsf_regressionmain} with $J^t, M^t \equiv 0$, $\mathcal{K} = R_D$ for $D = \Theta((m+h)Q_\star^2)$, and  $\eta_t = \Theta\left(\frac{Q_\star^2}{R\sqrt{t}\log T}\right)$ produces a sequence of predictions $\hat{y}_1, \ldots, \hat{y}_T$ for which
    $$\sum_{t=1}^T \|\hat{y}_t - y_t\| \leq O\left(Q_\star^2 \log(Q_\star) \cdot d_\mathcal{X} R \log(RT) m^2 \log^7(mT)  \cdot \sqrt{T}\right)  . $$
   % If $N$ denotes the number of eigenvalues $\lambda$ of the Koopman operator with $|\lambda| \geq 1-\frac{1}{m}$, then $Q_\star \leq 2^{N \log N}$.
\end{theorem}
\begin{remark}
% For some intuition, if $K: \mathcal{H} \to \mathcal{H}$ is the Koopman operator of the dynamics $f$ on some RKHS $\mathcal{H}$ and $$
% $N(\mathcal{H})$ denotes
    For some intuition, $Q_\star$ is upper bounded by $2^{\widetilde{O}(N)}$, where $N$ is the number of eigenvalues of (a truncation of) the Koopman operator outside the region $\left[0, 1 - \frac{1}{\sqrt{T}}\right] \cup \{1\} \cup \mathbb{D}_{1-\gamma}$. See Remark \ref{remark:perronfrobenius} for more detail about the relationship with the Koopman operator. In other words, Algorithm \ref{alg:clsf_regressionmain} gets a sublinear regret bound on systems which have a $\frac{1}{\sqrt{T}}$ spectral gap and Koopman eigenvalues that are either real or $\gamma$-stable, and it is able to adapt to $o(\log T)$ many eigenvalues that do not meet these criteria. See Section \ref{sec:discussion} for more discussion. For many physical systems of interest the Koopman eigenvalues are real (see Sections \ref{subsection:langevin} and \ref{sec:numerics}), which implies strong regret.
\end{remark}

While the dependence on $T$ is optimal in the theorem above, the constant $Q_\star$ could be very large. 
In Theorem \ref{thm:computational-lb} we provide a computational lower bound which exhibits a marginally stable and bounded dynamical system for which the number of steps needed to learn is exponential in the hidden dimension. 


Next, we arrive to the main technical tool for proving our results: an algorithm for learning marginally stable and asymmetric linear dynamical systems with controls and noise.


\begin{theorem}[Linear]
\label{theorem:main_linear}
Let $y_1, \ldots, y_T$ be the observations of an LDS driven by inputs $u_1, \ldots, u_T$ and disturbances $w_1, \ldots, w_T$ for which Assumption \ref{definition:linear_realizable} is satisfied. Fix the number of algorithm parameters $h = \Theta(\log^2T)$ and $m \geq 1$, and for notation define $\gamma = \frac{\log(mT)}{m}$ and $Q = R_BR_CQ_\star + Q_\star^2$. Set $\mathcal{K} = R_D$ for $D = \Theta\left((m+h) Q\right)$ and $\eta_t = \Theta\left(\frac{Q}{R \sqrt{t}\log T }\right)$, and let $\hat{y}_1, \ldots, \hat{y}_T$ be the sequence of predictions that Algorithm \ref{alg:clsf_regressionmain} produces. 
\begin{enumerate}
\item[(a).] In the noiseless setting $w_t \equiv 0$ let $Q_\star = Q_\star\left(A, C, [0, 1] \cup \mathbb{D}_{1-\gamma}\right)$. Then,
    \begin{align*}
        \sum_{t=1}^T \norm{\hat{y}_t - y_t} \leq O\left(Q R m^2 \log^7(mT) \cdot \sqrt{T} + R_C Q_\star R_{x_0 } \cdot (m + \log T)\right).
    \end{align*}
    
    \item[(b).] In the noiseless setting $w_t \equiv 0$ we may also define $Q_\star = Q_\star\left(A, C, \left[0, 1-\frac{1}{\sqrt{T}}\right] \cup \{1\} \cup \mathbb{D}_{1-\gamma} \right)$, in which case
        \begin{align*}
        \sum_{t=1}^T \norm{\hat{y}_t - y_t} \leq O\left(Q R m^2 \log_+(R_C Q_\star R_{x_0} / R)\log^7(mT) \cdot \sqrt{T}\right).
    \end{align*}

    \item[(c).] For any $\rho \in \left(0, \gamma\right]$, let $Q_\star = Q_\star(A, C, [0, 1-\rho] \cup \mathbb{D}_{1-\gamma})$. Then, with $R' = R + \frac{W}{\rho
    }$,
    \begin{align*}
        \sum_{t=1}^T \norm{\hat{y}_t - y_t} \leq O\left(Q \cdot R' m^2 \log^7(mT) \cdot \sqrt{T} + \frac{R' \log_+(R_C Q_\star R_{x_0} / R') + R_C Q_\star \sum_{t=1}^{T-1}\|w_t\|}{\rho}\right).
    \end{align*}
    
\end{enumerate}
In all cases, if $N$ denotes the number of eigenvalues outside of the desired region then $Q_\star \leq 2^{\widetilde{O}(N)}$. 
\end{theorem}
The above theorem provides an adaptive guarantee in the form of a data-dependent regret bound for Algorithm \ref{alg:clsf_regressionmain} run with a universal choice of parameters\footnote{The only way that any information about the data factors into Algorithm \ref{alg:clsf_regressionmain} is through the optimization, since $Q_\star$ determines the diameter $D$ and learning rates $\eta_t$. If one were to replace OGD with an adaptive/parameter-free optimizer, the algorithm would be fully data-independent.} $m, h$: for example, in result (a) the value of $Q_\star$ measures how easy the signal is to represent via a marginally stable observer system whose eigenvalues are either real or $\widetilde{O}\left(\frac{1}{m}\right)$-stable. If the system originally had a desirable spectral structure then $Q_\star = O(1)$, in which case the regret guarantees are strong. By contrast, if the system has very undesirable spectrum then these guarantees are weaker (since $Q_\star$ may be exponential in the number of undesirable eigenvalues). In general, as long as there are $o( \log T)$ many complex eigenvalues near the unit circle then we have sublinear regret. Furthermore, for the noisy guarantee (c) the predictions made compete with the best choice of $\rho$ in hindsight; if there is significant noise one must take larger $\rho$, whereas if there is almost no noise we can take $\rho$ very small, on the order of $\frac{1}{\sqrt{T}}$. Without knowing any of these considerations, the algorithm does as well as the best possible arrangement.

% We also provide a lower bound (Theorem \ref{theorem:lowerbound}) which states that in this setting the $\sum_{t \leq T} \|w_t\|$ term is optimal, up to constant factors.

To summarize, our results for linear dynamical systems exhibit the following important properties: 
\begin{enumerate}
    \item {\bf Asymmetric systems.} The spectral filtering guarantee is extended to capture asymmetric systems, without requiring small or clustered phases of complex eigenvalues as in previous work. 
    
    \item {\bf No hidden dimension dependence.} There is no explicit\footnote{The hidden dimension does not affect the runtime, number of parameters, or convexity of the optimization in any way. It indirectly affects $Q_\star$ since an LDS with larger $d_h$ has more eigenvalues, and $Q_\star$ scales with the number of eigenvalues that must be moved via pole placement.} dependence on the hidden dimension $d_h$.
    
    \item {\bf Noise robustness.}  Previous results in spectral filtering starting from \cite{hazan2017learning} exhibit compounding loss from noise (or have dependence on hidden dimension, such as Theorem 2 of \cite{hazan2018spectral}). This means that in the presence of adversarial noise, the instantaneous prediction error can be as large as $\|\hat{y}_T - y_T\| \approx T \|w\|$. In contrast, Algorithm \ref{alg:clsf_regressionmain} has  $\|\hat{y}_T - y_T\| \lesssim \|w\| +  \frac{1}{\sqrt{T} }$. Note that there is a lower bound of $\|\hat{y}_T - y_T\| \gtrsim \|w\|$, see Theorem \ref{theorem:lowerbound}.
    

\end{enumerate}
Our main result follows from combining the above theorem with a discretization-based reduction from nonlinear systems to high-dimensional linear systems given in Lemma \ref{lemma:discretization}.






% % \newpage

% % \section{OLD -- Algorithm and Main Result}
% % Our algorithm leverages spectral filtering of past outputs to build predictors for marginally stable, potentially asymmetric linear dynamical systems. By incorporating past observations $y_{t-1}, y_{t-2}, \ldots$ into the predictor via spectral filtering, the algorithm can handle asymmetric linear dynamics and achieves robustness to adversarial process noise. Algorithm~\ref{alg:clsf} formalizes this procedure.




% % \begin{algorithm}[htb]
% % \caption{Closed Loop Spectral Filtering}
% % \label{alg:clsf}
% % \begin{algorithmic}[1]
% % \STATE \textbf{Input:} Horizon \(T\), number of filters \(h\), step size $\eta$, convex constraint set \(\K\).
% % \STATE Compute $\{(\sigma_j, {\phi}_j)\}_{j=1}^k$ the top eigenpairs of a $T$-dimensional Hilbert matrix.
% % \STATE Initialize $M_i^0  \in \reals^{d_{out} \times d_{in}}, N_i^0  \in \reals^{d_{out} \times d_{out}}$ for all $i\in\{0,\dots,h\}$.
% % \FOR{$t = 0, \ldots, T-1$}
% % \STATE Compute $$\widehat{y}_t(M^t, N^t) = y_{t-1} + M_0^t u_{t-1} + \sum_{i=1}^h \sigma_i^{1/4}M_i^t \langle \phi_{j} ,  u_{t-1:t-T} \rangle + N_0^t y_{t-1} +  \sum_{i=1}^h \sigma_i^{1/4}N_i^t \langle \phi_{j} ,  y_{t-1:t-T} \rangle\,. $$
% % \STATE Compute loss $\ell_t(M^t, N^t) = \norm{\widehat{y}_t(M^t, N^t) - y_t}^2$.
% % \STATE Update $(M^{t+1}, N^{t+1}) \leftarrow \Pi_\K\left[(M^{t}, N^{t}) - \eta (\nabla_M\ell_t(M^{t}, N^t), \nabla_N\ell_t(M^t, N^{t}))\right]$. 
% % \ENDFOR
% % \end{algorithmic}
% % \end{algorithm}

% % Theorem~\ref{theorem:realizable} states the formal learning guarantee our algorithm achieves for linear dynamical systems.

% % \begin{theorem}[Linear]
% % \label{theorem:realizable}

% % Let $y_1, \ldots, y_T$ be the observations of an LDS driven by inputs $u_1, \ldots, u_T$ and disturbances $w_1, \ldots, w_T$ for which Assumption \ref{assumption:realizable} is satisfied. Then, with \evannote{mention alg params $h, \eta, \K$}, Algorithm \ref{alg:clsf} produces a sequence of predictions $\hat{y}_1, \ldots, \hat{y}_T$ such that
% % \begin{align}
% % \label{eqn:realizable}
% %     \sum_{t = 1}^T \| \widehat{y}_t - y_t\| \leq O\left(Q_\star\max\left\{R_\Theta, \frac{1+Q_\star}{\sigma_{\min}(C)}\right\}^2R\sqrt{Td_{out}}\log T\right)
% %     + \frac{\norm{C} Q_\star}{\rho} \left(\sum_{t=1}^T \|w_t\|\right) 
% % \end{align} 
% % \end{theorem}
% % We also provide a lower bound (Theorem \ref{theorem:lowerbound}) which states that in this setting the $\sum_{t \leq T} \|w_t\|$ term is optimal, up to constant factors.
% % There are 3 points to emphasize: 
% % \begin{enumerate}
% %     \item The spectral filtering guarantee is extended to capture asymmetric systems, without requiring small or clustered phases of complex eigenvalues. 
% %     \item There is no explicit\footnote{The hidden dimension does not affect the runtime, number of parameters, or convexity of the optimization in any way. It only affects $Q_\star$.} dependence on the hidden dimension $d_h$.
% %     \item  Informally, (\ref{eqn:realizable}) implies that the instantaneous loss of Algorithm \ref{alg:clsf} is $\|\hat{y}_T - y_T\| \lesssim \|w\| + \frac{1}{\sqrt{T}}$, and the lower bound states that for every algorithm there is a problem instance for which $\|\hat{y}_T - y_T\| \gtrsim \|w\|$. Note that an algorithm which allows errors to accumulate can have $\|\hat{y}_T - y_T\| \approx T \|w\|$ in the worst case \evannote{construct explicit examples for sysid, etc?}. In other words, Algorithm \ref{alg:clsf} is performing error correction by filtering over past observations.
% % \end{enumerate}

% % % \begin{assumption}
% % % Suppose that $\|y_t\| \leq R$ for all $t$ and that $f$ is continuous, and let $\pi$ be an $f$-invariant measure on $\mathcal{X}$. Suppose further that $M^2 := \mathbb{E}_{x \sim \pi} |h(x)|^2 < \infty$ and that the Koopman operator $K: L^2(\mathcal{X}, \pi) \to L^2(\mathcal{X}, \pi)$ is normal and compact, with complex eigenvalues $\lambda_j \to 0$ and eigenfunctions $\varphi_j$ such that $S := \sum_j |\varphi_j(x_0)|^2 < \infty$.
% % % \end{assumption}

% \begin{theorem}[General]
% \label{theorem:main}
% For any $\rho > 0$ and $N \in \mathbb{N}$, 
%  let $Q_\star$ denote the minimum value to the Luenberger program\footnote{This is covered in Section \ref{sec:luenbergerprogram}, and we apply it with $A = \operatorname{diag}(\lambda_1, \lambda_2, \ldots)$ and $C_j = \langle h, \varphi_j\rangle_{L^2(\mathcal{X}, \pi)}$ for $\varphi_j$ the eigenfunctions of $K$ on $L^2$.} (see Definition \ref{def:luenbergerprogram}), and define $Q := \max\{Q_\star, M\}$.
% Under these conditions, when run with $k = \Theta(\log^2 T \log(Q^3R^2p))$ filters, a constraint set with radius $\Theta(Q^2 \sqrt{k})$, and learning rate $\eta = \Theta\left(\frac{1}{Q^2R^3p \cdot k^2 \cdot \sqrt{T}}\right)$, Algorithm \ref{alg:mainA} generates predictions $y_t'$ with a loss bound
% \begin{align*}
%     \sum_{t \leq T} \|y_{t}' - y_t\| &\leq \underbrace{\|y_0\| 
%   \cdot
%   \frac{1}{1 - |\lambda_{N+1}|}}_{\text{truncation error}}+ \underbrace{M S \cdot \frac{1}{\rho}}_{\text{observer system error}} \\
%  &\quad +  \underbrace{M  Q_\star\|y_0\| \cdot  \frac{1}{\rho  (1-|\lambda_{N+1}|)}}_{\text{input mismatch due to truncation}} \\
%  &\quad + \underbrace{O\left(Q^6 R^3 p\log^2(Q^3 R^2 p) \sqrt{T} \log^6(T)\right)}_{\text{spectral filtering error}},
% \end{align*}
% % $$\sum_{t \leq T} \|y_{t}' - y_t\| = \widetilde{O}(\sqrt{T})$$
% \end{theorem}

% % \section{OLD Algorithm and Main Result}

% % \eh{this section needs to be reorganized a bit as follows: algorithm for nonlinear, theorem for nonlinear, then algorithm for linear, theorem for linear. }

% % Our main result is that incorporating past $y$'s into a predictor via spectral filtering allows us to learn marginally stable asymmetric LDS's in a way that is robust to adversarial noises. More precisely, we show:


% % \begin{algorithm}[htb]
% % \caption{closed loop spectral filtering }
% % \label{alg:clsf}
% % \begin{algorithmic}[1]
% % \STATE \textbf{Input:} Horizon \(T\), number of filters \(h\), step size $\eta$, convex constraint set \(\K\).
% % \STATE Compute $\{(\sigma_j, {\phi}_j)\}_{j=1}^k$ the top eigenpairs of a $T$-dimensional Hilbert matrix.
% % \STATE Initialize $M_i^0  \in \reals^{d_{out} \times d_{in}}, N_i^0  \in \reals^{d_{out} \times d_{out}}$ for all $i\in\{0,\dots,h\}$.
% % \FOR{$t = 0, \ldots, T-1$}
% % \STATE Compute $$\widehat{y}_t(M^t, N^t) = y_{t-1} + M_0^t u_{t-1} + \sum_{i=1}^h \sigma_i^{1/4}M_i^t \langle \phi_{j} ,  u_{t-1:t-T} \rangle + N_0^t y_{t-1} +  \sum_{i=1}^h \sigma_i^{1/4}N_i^t \langle \phi_{j} ,  y_{t-1:t-T} \rangle $$
% % % \STATE Update $\bz_{t} = (y_t - \hat{y}_t) + \sum_{i=1}^h K_i^t \langle \phi_{j} ,  z_{t-1:t-T} \rangle$.
% % % \STATE Predict $\tilde{y}_t = \hat{y}_t + \tilde{z}_t$.
% % \STATE Compute loss $\ell_t(M^t, N^t) = \norm{\widehat{y}_t(M^t, N^t) - y_t}^2$
% % \STATE Update to $(M^{t+1}, N^{t+1})$ using OGD or ONS rule and project to constraint set $\K$. 
% % \ENDFOR
% % \end{algorithmic}
% % \end{algorithm}






% % \begin{algorithm}[h!]
% % \caption{Nonlinear Spectral Filtering}
% % \label{alg:mainA}
% % \begin{algorithmic}[1]
% % \STATE \textbf{Input:} Horizon \(T\), number of filters \(k\), step size $\eta$, convex constraint set \(\K\).
% % \STATE Compute $\{(\sigma_j, \boldsymbol{\phi}_j)\}_{j=1}^k$  top eigenpairs of the $T$-dimensional Hankel matrix whose $i, j$-th entry is $\frac{2}{(i+j)^3 - (i+j)}$.
% \STATE Initialize $M_i^0 \in \reals^{p\times p}$ for all $i\in\{1,\dots,k\}$.
% \FOR{$t = 0, \ldots, T-1$}
% \STATE Compute $\hat{\by}_t(M^t) = \sum_{i=1}^k M_i^t  \langle y_{t-1:t-T} ,  \phi_i\rangle $
% \STATE Compute loss $\ell_t(M^t) = \norm{\hat{\by}_t(M^t) - \by_t}^2$
% \STATE Update $M^{t+1} \leftarrow \Pi_\K\left[M^t - \eta\nabla_M \ell_t(M^t) \right] $
% \ENDFOR
% \STATE \Return $M^T$
% \end{algorithmic}
% \end{algorithm}


\section{Setting} 

\subsection{Notation}
We use $\mathbb{D}_r$ to denote the complex disk of radius $r$, i.e. $\mathbb{D}_r = \{\lambda \in \mathbb{C}:\; |\lambda| \leq r\}$.
We use the notation $v_{a:b}$ to denote stacking the sequence $v_a, v_{a+1}, \ldots, v_{b-1}, v_b$. We allow indices to be negative for notational convenience, and quantities with negative indices are understood as zero. We define $\log_+(x) = 1 + \max(0, \log(x))$ for positive $x$. Unless otherwise stated, the norm $\|\cdot\|$ refers to the Euclidean $\ell_2$ norm for vectors and the corresponding operator/spectral norm for matrices. $\|\cdot\|_F$ denotes the Frobenius norm of a matrix. For $A \in \R^{d \times d}$ complex diagonalizable we let $$\kappa_{\textrm{diag}}(A) = \norm{H} \cdot \norm{H^{-1}}$$
be the condition number of the best similarity transformation $H \in \mathbb{C}^{d\times d}$ for which $A = HDH^{-1}$ with $D$ diagonal, which is sometimes called the "spectral condition number" of $A$. If $A$ is not diagonalizable we set $\kappa_{\textrm{diag}}(A) \equiv +\infty$, and if $A$ is normal then $\kappa_{\textrm{diag}}(A) = 1$. 

% For any vector space $X$ over the field $\mathbb{F}$, we let $\mathcal{F}(X, \mathbb{F})$ denote the set of all functionals $f: X \rightarrow \mathbb{F}$. 

\subsection{Problem Setting}
% Consider a partially observed deterministic nonlinear dynamical system with state space $\mathcal{X} \subseteq \reals^d$ and observation space $\mathcal{Y} \subseteq \reals^p$ given by
% $$ x_{t+1} = f(x_t) ,\quad  y_{t+1} = h(x_{t+1}) , $$
% for a given dynamics function $f: \mathcal{X} \to \mathcal{X}$, observation function $h: \mathcal{X} \to \mathcal{Y}$, and $x_0 \in \mathcal{X}$. 
% The goal is to accurately predict $y_{t+1}$ given the sequence of observations $(y_0, \ldots, y_t)$, and without knowing the dynamics. 
% We suppose for simplicity that the observations are one-dimensional, but our results extend to arbitrary hidden and observation dimensions. 
We consider nonlinear dynamical systems of the form (\ref{eqn:NLDS}) and linear systems of the form (\ref{eqn:LDS}), both of which generate a sequence $y_1, \ldots, y_T$ of observations.
The prediction notion we use is borrowed from the setting of online convex optimization \cite{hazan2016introduction}, where the performance of an algorithm is measured by its regret. Regret is a robust performance measure, as it does not assume a specific generative or noise model. It is defined as the difference between the loss incurred by our predictor and that of the best predictor in hindsight belonging to a certain class, i.e. 
\begin{equation*} \label{eqn:regret}
\regret_T(\mA, \Pi) = \sum_{t=1}^T \ell(y_t,\hat{y}_t^\mA ) - \min_{\pi \in \Pi} \sum_{t=1}^T \ell(y_t,\hat{y}_t^\pi) . 
\end{equation*}
Here $\mA$ refers to our predictor, $\hat{y}_t^\mA$ is its prediction at time $t$, and $\ell$ is a loss function. Further, $\pi$ is a predictor from the class $\Pi$, for example all predictions that correspond to noiseless linear dynamical systems of the form $\hat{y}_{t+1}^\pi = \sum_{i=0}^t C A^i B u_{t-i} $. A stronger predictor class is that of all closed-loop linear dynamical predictors that also take into account the initial state and noise, even though they are not observed by an algorithm, i.e. 
$$\hat{y}_{t+1}^\pi = CA^{t+1}x_0 + \sum_{i=0}^{t-1} C A^i (B u_{t-i} + w_{t-i}).
% + \sum_{i=1}^{t-1} C A^i w_{t-i} . 
$$

Our main results are proved as loss bounds in realizable settings, i.e. when the data is generated from a ground truth dynamical system and noise sequence. Since the derivations involved go through the language of online convex optimization, these can be converted to agnostic regret guarantees, which we do in Appendix \ref{appendix:otherguarantees}. We define the linear and nonlinear dynamical data generators for the realizable case as follows: 

\begin{assumption}[Linear data generator]
\label{definition:linear_realizable}
    The input/output signal $(u_t, y_t)_{t=1}^T$ is generated by an LDS $(A, B, C, x_0)$ as per equations (\ref{eqn:LDS}) with arbitrary perturbations $w_1, \ldots, w_T$. We assume that:
    \begin{enumerate}
    % \item[(i)] The dynamics matrix $A$ is diagonalizable over $\mathbb{C}$ with all eigenvalues in the unit disk \evannote{do we need $A$ to be diagonalizable?} \elad{I don't think so, and also all matrices up to vanishing perturbation are diagonalizable over $C$, what's the assumption here?}. We also assume w.l.o.g.\footnote{For a general diagonalizable matrix $A \in \R^{d_{\mathrm{h}} \times d_{\mathrm{h}}}$ with eigenvalues in the unit disk, we can write $A = H\tilde{A}H^{-1}$ for $\tilde{A}$ diagonal. Defining $\tilde{B} := H^{-1}B$ and $\tilde{C} := CH$, we have $\|\tilde{A}\| \leq 1$ and that the LDS's $(A, B, C)$ and $(\tilde{A}, \tilde{B}, \tilde{C})$ are equivalent.} that $A$ is Lyapunov stable, i.e. $\|A\| \leq 1$. \elad{why is it w.l.o.g, and how is comment (3) related?} \evannote{if $A$ is ill-conditioned, then it could have all eigenvalues in the unit disk but $\|A\|> 1$. In such a setting, there is an equivalent LDS where $A$ is well-conditioned and the $H$ and $H^{-1}$ go into $B$ and $C$... i'm just saying here that we will consider that system w.l.o.g.}
    \item[(i)] $A$ is diagonalizable with all eigenvalues in the unit disk $\mathbb{D}_1$ and $\kappa_{\textrm{diag}}(A) = 1$\footnote{For an LDS ($A,B,C$) with a general diagonalizable $A=HDH^{-1}$, one may consider the equivalent LDS ($D,H^{-1}B,CH$), in which case the product $R_BR_C$ which appears in our regret bounds scales up by $\kappa_{\textrm{diag}}(A)$.}.
    \item[(ii)] The system is bounded, i.e. $\|B\|_F\leq R_B,$ $\|C\|_F \leq R_C$, and $\|x_0\| \leq R_{x_0}$.
    \item[(iii)] For all $t$, the inputs and outputs are bounded by $\|u_t\|,\|y_t\| \leq R$, and $\|w_t\| \leq W$. %\elad{I added bound on perturbations} \evannote{we actually dont require a bound on $w_t$, only on $y_t$}
    \item[(iv)] The system $(A, C)$ is observable\footnote{This means that the matrix $\mathcal{O} = [C,CA,CA^{2},\ldots, CA^{d_{\text{obs}}-1}] \in \R^{d_\text{obs} d_h \times d_h}$ has rank $d_h$. In the case of 1-dimensional observations, observability is equivalent to invertibility of $\mathcal{O}$, and in general it describes whether the observation space sees all directions of state space. If $A$ has an eigenvalue with  multiplicity $> d_{\textrm{obs}}$, then $(A, C)$ cannot be observable. See Section 10.1 of \cite{hazan2025introductiononlinecontrol} for more information. The observability assumption can be weakened by ignoring the modes of $A$ corresponding to eigenvalues that don't need to be moved by pole placement -- for simplicity, we will assume complete observability.}.  
    \end{enumerate}
    
\end{assumption}
% Note that we do not make explicit assumptions about the disturbances $w_t$, though the condition that $\|y_t\| \leq R$ requires them to be bounded in some sense. 

\begin{assumption}[Nonlinear data generator]
\label{assumption:nonlinear}
    The signal $(y_t)_{t=1}^T$ is generated by a nonlinear dynamical system $(f,h,x_0)$ as per equations (\ref{eqn:NLDS}). We assume that:
    \begin{itemize}
    \item[(i)] For all $t$, the states and signal are bounded by $\norm{x_t}, \norm{y_t} \leq R$.
    \item[(ii)] The dynamics $f$ and observation $h$ are $1$-Lipschitz.
    \item[(iii)] The discretized linear lifting $(A', C')$ defined in Lemma \ref{lemma:discretization} is an observable LDS.     % \item[(i)] There exists a \evannote{Hilbert?} space of functionals $(\mathcal{F}, \langle \cdot , \cdot \rangle_\mathcal{F})$ such that $h \in \mathcal{F}$, $M := \norm{h}_\mathcal{F} < \infty$, and $\mathcal{F} \circ f \subseteq \mathcal{F}$.
        % \item[(ii)] For all $t$, the signal is bounded by $\norm{y_t} \leq R$.
        % % \item[(iii)] $M := \Exp_{x \sim \pi}[|h(x)|^2] \leq \infty$.
        % \item[(iii)] The Koopman operator associated with $f$ over $\mathcal{F}$ is compact, with complex eigenvalues $\lambda_j \rightarrow 0$ and eigenfunctions $\varphi_j$ such that $S := \sum_{j=1}^\infty |\varphi_j(x_0)|^2 < \infty$.
    \end{itemize}
\end{assumption}
% The quantities $M$ and $R$ are determined by the observation structure $h$, while the underlying dynamical system $f$ and the choice of $\mathcal{F}$ determine the decay of $\lambda_j$ and $\varphi_j(x_0)$. We discuss the Koopman operator in more detail in the next subsection.





\section{Analysis}

In this section we perform our main analysis. We first detail our construction of the Luenberger program and $Q_\star$ and its implications for learning linear systems, culminating in a proof of Theorem \ref{theorem:main_linear}. We then describe a reduction from nonlinear systems to high-dimensional linear systems via discretization, from which Theorem \ref{theorem:main_nonlinear} follows.
% \evannote{lastly, we do symmetric case}

\subsection{Linear Dynamical Systems}

In this subsection, we prove Theorem~\ref{theorem:main_linear}. Given a linear dynamical system (LDS), we construct another LDS that incorporates past observations as feedback, has a desirable transition matrix, and generates the same trajectories as the original system. This construction relies on the \textbf{Luenberger observer} framework \cite{luenbergerobserver1971}, which we sketch next.

For a given LDS $(A, B, C, x_0)$, we can always choose a gain $L \in \R^{d_h \times d_{\textrm{out}}}$ and construct the observer system
$$\tilde{x}_{t+1} = A\tilde{x}_t + Bu_t + L(y_t - \tilde{y}_t), \quad \tilde{y}_t = C\tilde{x}_t,\quad \tilde{x}_0 = x_0.$$
% $$\tilde{x}_{t+1} = (A-LC)\tilde{x}_t + Bu_t + Ly_t, \quad \tilde{y}_t = C\tilde{x}_t,\quad \tilde{x}_0 = x_0.$$
This system has the same dynamics as the original one, but with an additional term that uses the observation error $y_t - \tilde{y}_t$ as linear feedback. Since $y_t = Cx_t$ and $\tilde{y}_t = C\tilde{x}_t$, if we define the state error $\xi_t := x_t - \tilde{x}_t$ then we get
\begin{align*}\xi_{t+1} &= Ax_t + Bu_t + w_t - A\tilde{x}_t - Bu_t - LC(x_t - \tilde{x}_t) \\&= (A-LC)\xi_t + w_t.\end{align*}
Rolling this out reveals that for all $t$, $$x_t = \tilde{x}_t + (A-LC)^t (x_0 - \tilde{x}_0) + \sum_{t=1}^t (A - LC)^{i-1} w_{t-i}.$$

It is a known fact in control theory that if $(A, C)$ is observable then we can choose $L$ to place the eigenvalues of $A-LC$ as desired, which we describe further in Appendix \ref{appendix:luenbergerprogram}. 
% In particular, in the absence of noise we can make $\xi_t$ decay exponentially and get that $\|y_t - \tilde{y}_t\| \lesssim e^{-t}$. 
This is a standard construction which is frequently used with system identification for prediction and control from partial observations \cite{poleplacementapplication}, and the Kalman filter is related to a special case in which the gain $L$ is chosen optimally under Gaussian noise assumptions \cite{kalman1960new}. Our approach will be to predict via spectral filtering, which learns improperly and competes in hindsight against all observer systems.

For certain systems $(A, C)$ it is harder to infer hidden state information from observations, which is reflected in terms of the size of $L$ and the complexity of the observer dynamics $A - LC$. Many Luenberger observer-based methods which require $L$ explicitly end up approaching this through heuristics or structural assumptions \cite{luenbergerassumptionsA,luenbergerassumptionsB} or by performing a numerically-difficult optimization over $L$ \cite{luenbergeroptimization}. However, since spectral filtering matches the performance of the \textit{best} observer system in hindsight, we will make use of a quantity that captures the complexity of the best observer system:

\begin{definition}[Luenberger program]
\label{def:luenbergerprogram}
    Let $A \in \R^{d_h \times d_h}$ and $C \in \R^{d_{\textrm{out}} \times d_h}$. For each spectral constraint set $\Sigma \subset \mathbb{C}$ we consider the optimization problem over $L \in \R^{d_h \times d_{\textrm{out}}}$ given by
    $$\min \quad \kappa_{\textrm{diag}}(A-LC) \quad \text{s.t.} \quad A-LC \text{ has all eigenvalues in $\Sigma$}.$$
    When the constraint set of this optimization is nonempty, denote the minimal value by $Q_\star = Q(A, C, \Sigma) < \infty$ and a minimizing gain by $L_\star$.
\end{definition}
We will be applying Algorithm \ref{alg:clsf_regressionmain} to improperly learn over the class of observer systems. As we will see in Lemma \ref{lemma:spectralfiltering}, spectral filtering does well for $\Sigma \subset [0, 1]$ and very well for $\Sigma \subset [0, 1-\rho] \cup \{1\}$\footnote{Negative eigenvalues can be handled by spectral filtering easily with a constant amount of extra parameters and computation, see the extension in Theorem 3.1 of \cite{agarwal2023spectral}. For ease of presentation, we will continue to state things for nonnegative real eigenvalues.}, and the regression terms are able to handle $\Sigma \subset \mathbb{D}_{1-\gamma}$, so we will be mostly interested in $\Sigma = [0, 1-\rho] \cup \{1\} \cup \mathbb{D}_{1-\gamma}$. For generality, we will state things for arbitrary $\Sigma$ and specify where needed.

In Appendix \ref{appendix:luenbergerprogram} we will say more about this optimization and apply known bounds on $Q_\star$, and in particular we show that for observable systems $Q_\star \leq (\frac{N}{\operatorname{vol}(\Sigma)})^{O(N)}$, where $N$ is the number of system eigenvalues outside the desired region $\Sigma$. This can be tight.

\begin{remark}
The objective function $Q_\star$ is related to the quantity $\mathcal{S}_2$ from \cite{robustpoleplacement}, which describes the stability of eigenvalue placement. For this reason, such minimizations are referred to as "robust pole placement", see also the survey \cite{luenbergeroptimization}. It is known that the eigenvalue placement problem is very ill-conditioned, and small perturbations to the data $(A, C)$ yield $Q_\star$-sized changes in $L$ and the placed eigenvalues (see e.g. Theorem 3.1 of \cite{Mehrmann1996}). If $Q_\star$ is large, then a pipeline based on eigenvalue placement combined with system identification from data can be very brittle.
% for the above reason, it is still one of the common approaches to predicting and controlling an LDS from partial observation data \evannote{cite someone?}. From this perspective, Algorithm \ref{alg:clsf} can be viewed as a convex relaxation which is able to compete against the best observer system/eigenvalue placement construction, but without suffering this numerical instability.
\end{remark}
In words, the value of $Q_\star$ describes how hard it is (measured in terms of conditioning of the closed-loop dynamics) to approximately recover the state using a desirable observer system\footnote{
% We can restrict to real diagonalizable observers in order to apply spectral filtering, and since we will use autoregressive components our algorithm will tolerate $\gamma$-stable eigenvalues with nonzero complex part too. While spectral filtering works for marginally stable systems, at times we may require $A-LC$ to be strongly stable in order for the effects of additive process noise to decay.
Other improper algorithms will require a definition of $Q_\star$ tailored to their structure, see Appendix \ref{appendix:otherguarantees}. To maintain generality we will often use the phrases "desirable spectral structure" or "undesirable eigenvalues", where the implied spectral constraint set $\Sigma$ will depend on the algorithm.}.
We can describe the observer system construction in terms of $Q_\star$:

\begin{definition}[Observer system]
\label{prop:observersystem}
     Consider the following LDS
     \[
        x_{t+1} = Ax_t + Bu_t + w_t \,,\ \quad y_t = Cx_t\,.
     \]
     Fix $\Sigma \subset \mathbb{D}_{1-\delta}$ for some $\delta \in [0, 1]$, let $Q_\star = Q(A, C, \Sigma)$ be the value of the corresponding Luenberger program from Definition \ref{def:luenbergerprogram}, and suppose that $Q_\star < \infty$. Then, there exists an observer system which takes as inputs $(u_0, \dots, u_t)$ and $(y_0, \dots, y_t)$ and produces the estimated observation $\tilde{y}_{t+1}$ via the LDS
     \[
     \tilde{x}_{t+1} = \tilde{A}\tilde{x}_t + Bu_t + Ly_t\,, \quad \tilde{y}_t = C\tilde{x}_t\,, \quad \tilde{x}_0 = x_0
     \]
    with the properties:
    \begin{itemize}
        \item[(i)] $\|L\| \leq \frac{2Q_\star}{ \sigma_{\textrm{min}}(C)}$, where the denominator is the minimal singular value.
        \item[(ii)] $\tilde{A} = HDH^{-1}$ where $D$ is a diagonal matrix with entries in $\Sigma$ and $\kappa_{\textrm{diag}}(\tilde{A}) \equiv \norm{H}\norm{H^{-1}} = Q_\star$.
        \item[(iii)] $\sum_{t=1}^T\|y_t - \tilde{y}_t\| \leq \frac{Q_\star \|C\|}{\delta}\left(\sum_{t=0}^{T-1} \norm{w_t}\right)$, and if $w_t \equiv 0$ then $y_t = \tilde{y}_t$ even if $\delta = 0$.
    \end{itemize}
\end{definition}
Let $L_\star = L(A, C, \Sigma) \in \R^{d_{h} \times d_{\textrm{out}}}$ be a solution of the Luenberger program.
% , and define $\tilde{A} := A - L_\star C$. 
Properties (i) and (ii) follow from the objective function and constraint set of the Luenberger program, along with Lemma \ref{lemma:ackermannbounds}. As mentioned previously, if we define the errors $\xi_t := x_t - \tilde{x}_t$ then $\xi_t = \sum_{i=1}^t (A-L_\star C)^{i-1} w_{t-i}$ since $\xi_0 = 0$; so, in the noiseless setting we exactly reproduce the states and observations. In the noisy case, since $y_t - \tilde{y}_t = C\xi_t$ and all eigenvalues of $\tilde{A} = A-L_\star C$ have magnitude $\leq 1-\delta$, \begin{align*}
    \sum_{t=1}^T \|\tilde{y}_t - y_t\| &\leq \norm{C}Q_\star\sum_{t=1}^T \sum_{i = 1}^t (1-\delta)^{i-1} \|w_{t-i}\| \leq \frac{\norm{C}Q_\star}{\delta}\cdot  \left(\sum_{t=0}^{T-1} \|w_t\|\right)
    \end{align*}

The observer system construction shows that there is an LDS over $(u, y)$ with desirable eigenstructure that nearly matches the original one, which may be rolled out without requiring $w$ (without noise, the observer system matches it exactly). We will combine this with spectral filtering's regret guarantee, stated below for the unsquared $\ell_2$ loss (with immediate extension to Lipschitz losses -- a similar result can be proven for the squared loss). In addition to the inclusion of autoregressive terms in the algorithm, there is also a sharpening of the analysis that may be of independent interest, see Remark \ref{remark:spectralfilteringproof}. In the application of the below lemma, it should be imagined that the inputs $u_t$ are actually the stacked $[u_t, y_{t-1}]$ and we are running the algorithm with $P^t,N^t \equiv 0$ in order to compete against an observer LDS.

\begin{lemma} 
\label{lemma:spectralfiltering}
    On any sequence $(u_t, y_t)_{t=1}^T$ such that $\norm{u_t}, \norm{y_t} \leq R$, Algorithm \ref{alg:clsf_regressionmain} run with $P^t_1 = 1$ and $P^t_j, N^t \equiv 0$ for $j > 1$, $h = \Theta (\log^2 T)$, $m = \Theta\left(\frac{1}{\gamma}\log\left(\frac{T}{\gamma}\right)\right)$, $\mathcal{K} = B_{D}$ with $D = \Theta((m+h)R_BR_C\kappa)$, and $\eta_t = \Theta\left(\frac{R_BR_C\kappa}{R\sqrt{t}\log T }\right)$, achieves regret
    \[
    \sum_{t=1}^T \norm{\hat{y}_t - y_t} - \min_{\pi \in \Pi}\sum_{t=1}^T\norm{{\hat{y}_t^\pi} - y_t} \leq O\left(\frac{\log^7\left(\frac{T}{\gamma}\right)}{\gamma^2} \cdot R_B R_C \kappa R \cdot \sqrt{T} + \frac{R + \log_+(R_C \kappa R_{x_0}/R)}{\rho}\right),
    \]
    where $\Pi=LDS(\rho, \gamma, \kappa, R_{B}, R_C, R_{x_0}, R)$ is the class of LDS predictors $\pi$ parameterized by $(A, B, C, x_0)$ where $A$ is diagonalizable with eigenvalues in $[0, 1-\rho] \cup \{1\} \cup \mathbb{D}_{1-\gamma}$ for $\rho \leq \gamma$, $\kappa_{diag}(A) \leq \kappa$, $\norm{B}_F \leq R_B$, $\norm{C}_F \leq R_C$, $\|x_0\| \leq R_{x_0}$, and $\|\hat{y}_t^\pi\| \leq R$. 
    If the spectral gap $\rho = 0$, then the second term can be replaced by $R_C \kappa R_{x_0} \cdot \left(\frac{1}{\gamma} + \log T\right)$.
\end{lemma}
\begin{proof}
See Appendix \ref{appendix:SFlemma}.
\end{proof}

% \evannote{i adjusted this statement to have autoregressive + SF, need to also adjust the proof in Appendix \ref{appendix:SFlemma}. can we also keep track of $\|B\|$ and $\|C\|$ separately, and make sure that instead of $\|C\|_F$ we only sum over $\log T$ many rows?} \eh{bound needs to be strengthed to depend on rows of $CH$ and columns of $HB$}

% \begin{lemma}[Theorem 1 of \cite{hazan2017learning}] 
% \label{lemma:spectralfiltering}
%     On any sequence $(u_t, y_t)_{t=1}^T$ such that $\norm{u_t}, \norm{y_t} \leq R$, Algorithm \ref{alg:clsf_regressionmain} run with $h = \Theta (\log^2 T)$, $m = \Theta\left(\frac{1}{\rho}\log\left(\frac{T}{\rho}\right)\right)$, $\mathcal{K} = B_{R_M}$ with $R_M = \Theta (\kappa R_\Theta^2 \sqrt{hm})$, and $\eta = \Theta\left(\kappa R_\Theta^2/ R \sqrt{Td_{out}}\log T\right)$, achieves regret
%     \[
%     \sum_{t=1}^T \norm{\hat{y}_t - y_t} - \min_{\pi \in LDS(\frac{1}{m}, \kappa, R_\Theta)}\sum_{t=1}^T\norm{{\hat{y}_t^\pi} - y_t} \leq O\left(\frac{\kappa R_\Theta^2 R}{\rho} \sqrt{Td_{out}}\log^3 T \log\left(\frac{T}{\rho}\right)\right)
%     \]
%     where $LDS(\rho, \kappa, R_\Theta)$ is the class of LDS predictors parameterized by $\Theta = (A, B, C, x_0)$ where $A$ is diagonalizable with eigenvalues in $[0, 1] \cup \mathbb{D}_{1-\rho}$, $\kappa_{diag}(A) \leq \kappa$, and  $\norm{B}_F, \norm{C}_F, \norm{x_0} \leq R_\Theta$.
% \end{lemma}


\begin{proof}[Proof of Theorem \ref{theorem:main_linear}]
    Suppose that the data is generated by a ground truth LDS $(A, B, C, x_0)$ with disturbances $w_t$ such that Assumption \ref{definition:linear_realizable} is satisfied. Observability implies that $Q_\star(A, C, \Sigma) < \infty$ by Lemma \ref{lemma:ackermann}, where we only move those eigenvalues outside of $\Sigma$.
    Let $(\tilde{A}, \tilde{B}, C, x_0)$ be the observer system constructed in Definition \ref{prop:observersystem}, which we may treat as an LDS over concatenated inputs $[u_t, y_{t-1}]$ with a block-diagonal input gain $\tilde{B} = \begin{bmatrix}
        B & 0 \\ 0 & L
    \end{bmatrix}$. For $d_\mathcal{Y}=1$ this will have the bounds $\kappa_{\operatorname{diag}}(\tilde{A}) \leq Q_\star$ and $R_{\tilde{B}} = O(R_B + \frac{Q_\star}{R_C})$.
    
    For (c) we use the given $\rho$ and $\gamma = \frac{\log(mT)}{m}$ (so that $\frac{1}{\gamma} \log(\frac{T}{\gamma}) \leq m$) and define $\Sigma = [0, 1-\rho] \cup \mathbb{D}_{1-\gamma} \subset \mathbb{D}_{1-\rho}$. This is in the comparator class $\Pi = \mathrm{LDS}\left(\rho, \gamma, Q_\star, R_B + \frac{Q_\star}{R_C}, R_C, R_{x_0}, R + \frac{W}{\rho}\right)$, and so the error bound of the observer system implies
    $$\min_{\pi \in \Pi}\sum_{t=1}^T \|\hat{y}_t^\pi - y_t\| \leq \frac{R_CQ_\star}{\rho} \left( \sum_{t=0}^{T-1} \|w_t\|\right).$$
    The result follows from directly applying the spectral filtering guarantee of Lemma \ref{lemma:spectralfiltering}.

    For (b) we take $\Sigma = [0, 1 - \rho] \cup \{1\} \cup \mathbb{D}_{1-\gamma}$ with $\rho = \frac{1}{\sqrt{T}}$. The addition of the $\{1\}$ eigenvalue does not change the comparator class, and so we may apply the result of (c). Result (a) is proved the same way, using the $\rho = 0$ case of Lemma \ref{lemma:spectralfiltering}.
\end{proof}

\subsection{Nonlinear Dynamical Systems}\label{sec:koopman-discrete}

% \section{Discrete Koopman Theorem} 


In this subsection, we will prove Theorem \ref{theorem:main_nonlinear}. Our main tool is a space-discretized Markov chain argument that achieves the same goal as the Koopman operator, which may be of independent interest. The discretization argument is sufficient for sequence prediction and avoids traditional compactness and spectral assumptions, replacing them by simpler boundedness and Lipschitz conditions on the dynamics.  

\begin{lemma}
\label{lemma:discretization}
    Let $y_1, \ldots, y_T$ be generated by a deterministic discrete-time dynamical system $(f, h, x_0)$ satisfying Assumption \ref{assumption:nonlinear}. Let $\mathcal{S}$ be an $\epsilon$-net of the ball of radius $R$ in $\mathcal{X}$. The cardinality of $\mathcal{S}$ satisfies 
    $$N = |\mathcal{S}| \leq \left(\frac{2R}{\epsilon}\right)^{d_\mathcal{X}}.$$
    Then, there exists an LDS of hidden dimension $N$ with states $z_t \in \R^N$ and an observation matrix $C' \in \R^{d_\mathcal{Y} \times N}$ with Frobenius norm $\|C'\|_F \leq R \sqrt{Nd_{\mathcal{Y}}}$, such that the output sequence $y_t' = C'z_t$ satisfies 
    $$\sum_{t \leq T }\norm{y_t' - y_t} \leq \frac{T^2}{2}\epsilon.$$
\end{lemma}
\begin{proof}
    The $\epsilon$-net condition guarantees that for every $x \in \mathcal{X} \cap B_R$ there exists $s \in \mathcal{S}$ such that $\|x - s\| \leq \epsilon$. It is standard that such an $\epsilon$-net can be chosen with $|\mathcal{S}| = N$, and we enumerate the elements as $s^{(1)}, \ldots, s^{(N)}$.

    Define the mapping $\pi: \mathcal{X} \to \mathcal{S}$ as the nearest neighbor projection, i.e. $\pi(x) = \argmin_{s \in \mathcal{S}} \|x - s\|$. 
We consider the deterministic transition operator $T: \mathcal{S} \to \mathcal{S}$ as $T(s) = \pi(f(s))$, yielding a Markov chain with states $s_t$ in a finite state space. If we let $z_t \in \R^N$ be the indicator vectors corresponding to this state and the evolution operator $A' \in \R^{N \times N}$ with $A'_{ij} = 1$ if $T(s^{(j)}) = s^{(i)}$ and $0$ else, then we get a linear evolution
$$z_{t+1} = A'z_t, \quad y_t' = Cz_t$$
with $C \in \R^{d_\mathcal{Y} \times N}$ sending the one-hot vector $z$ corresponding to $s^{(j)}$ to $h(s^{(j)})$. Since $\|y_t\| \leq R$, the matrix elements of $C'$ can be taken to be bounded by $R$.

With these definitions in place, let $x_0, \ldots, x_T$ be a trajectory of the original nonlinear system with observations $y_0, \ldots, y_T$. We let $s_0 = \pi(x_0)$ and $s_{t+1} = T(s_t)$ be the rollout of the finite state space Markov chain, and define $\xi_t = \|s_t - x_t\|$ the error against the true states. Then,
\begin{align*}
    \xi_{t+1} &= \|s_{t+1} - x_{t+1}\| \\
    &= \|\pi (f(s_t)) - f(x_t)\| \\ &\leq \|f(s_t) - f(x_t)\| + \|f(s_t) - \pi (f(s_t))\| \\
    &\leq \xi_t + \epsilon,
\end{align*}
where we used 1-Lipschitzness of $f$ and the $\epsilon$-net condition in the last line. Since $\xi_0 \leq \epsilon$ again by the $\epsilon$-net construction, we find that $\xi_t \leq t\epsilon$ and so $\|y_t' - y_t\| \leq t \epsilon$
by Lipschitzness of $h$.
\end{proof}

This essentially says that a bounded and marginally stable nonlinear dynamical system can be approximated via a $T^{O(d_\mathcal{X})}$-dimensional LDS over a horizon of length $T$. 
An alternate approach to approximating nonlinear systems via high-dimensional linear ones is to use a spectral truncation of the Koopman operator, as outlined in e.g. \cite{brunton2021modernkoopmantheorydynamical}. However, 
proving convergence of the finite-dimensional spectral truncations of Koopman operators can require strong assumptions or compactification tricks such as different input and output spaces \cite{IEEEkoopmandifferentspaces}, without which 
it becomes difficult to describe convergence of continuous and residual spectra. Furthermore, the mechanism of pole placement via observer systems becomes obscured in such settings.

We adopt the improper learning perspective that the only thing that matters is approximately correct observations $y_t' \approx y_t$, and that convergence of spectra is not strictly necessary for sequence prediction. The discretization and Markov chain outlined above achieves this, and the proof is clean and simple with fewer assumptions\footnote{Note, however, that any construction made via truncated Koopman operators will also fall into the comparator class of high-dimensional LDS, against which we have a regret guarantee from Theorem \ref{theorem:main_linear}. If there were a better high-dimensional LDS that approximates the sequence $y_1, \ldots, y_T$, then Algorithm \ref{alg:clsf_regressionmain} would match it without needing to know anything about it.}. 
% To demonstrate the flexibility of our improper learning approach, we will prove \evannote{add whatever symmetric guarantee we will make} using a more standard Koopman formulation later on in \evannote{ref wherever we do it}. 

\begin{remark}
\label{remark:perronfrobenius}
    This Markov chain is closely related to the Perron-Frobenius operator $P$, the adjoint to the Koopman operator. Informally, $P$ is a linear operator on spaces of measures on the state space $\mathcal{X}$ which sends a measure $\mu$ to its pushforward under the dynamics, i.e. $P\mu = f_\#\mu$. One could imagine an infinite-dimensional LDS of the form
    $$\mu_{t+1} = P\mu_t,\quad y_{t} = \int_{\mathcal{X}} h(x) d\mu_t(x), \quad \mu_0 = \delta_{x_0}$$
    which generates the same observations as the original system,
    where both the dynamics and observations are linear with respect to the evolving measure. The Markov chain in Lemma \ref{lemma:discretization} is constructing exactly this LDS, but on a discretized state space to ensure finite-dimensionality (i.e. probability measures are represented via probability vectors). This discretization argument circumvents any need for proofs of convergence of spectral truncations of the Koopman operator or its adjoint $P$.
\end{remark}
Importantly, the observation that the Markov chain constructed in Lemma \ref{lemma:discretization} is (dual to) the Koopman operator implies that the discretized LDS retains the spectral properties of the Koopman operator. Technically, since we use the Euclidean norm for the analysis of this discretized LDS we inherit the spectral properties of the Koopman operator on $L^2(\mathcal{X}, \pi)$ function spaces, where the choice of measure $\pi$ is equivalent to the choice of discretization grid. A very useful result of our improper learning formulation of the problem is that our analysis adapts in hindsight to the best choice of discretization grid, i.e. the Koopman operator on a function space that gives it the most desirable spectral properties. In this way, we sidestep the suboptimalities and difficulties of any explicit choices of function space.


In any case, we are left with a high-dimensional LDS which approximates the original system, and we will conclude by applying Theorem \ref{theorem:main_linear}. First, we define the relevant quantities which we use in the statement and proof of Theorem \ref{theorem:main_nonlinear}.

\begin{definition}[Nonlinear Luenberger program]
\label{def:nonlinearluenbergerprogram}
    Let $(f, h, x_0)$ be a nonlinear dynamical system as per (\ref{eqn:NLDS}) with $\|x_t\|, \|y_t\| \leq R$ and $f, h$ being $1$-Lipschitz. For any discretization $\epsilon > 0$ and spectral constraint set $\Sigma \subset \mathbb{C}$, we consider the optimization program over $L \in \R^{\left(\frac{2R}{\epsilon}\right)^{d_\mathcal{X}} \times d_\mathcal{Y}}$ given by
    $$\min \quad \kappa_{\textrm{diag}}(A'-LC') \quad \text{s.t.} \quad A'-LC' \text{ has all eigenvalues in $\Sigma$},$$
    where $A'$ and $C'$ are the Markov chain transition and observation matrices constructed in Lemma \ref{lemma:discretization}.
    When the constraint set of the optimization is nonempty, with an overloading of notation we denote the minimal value by $Q_\star = Q_\star(f, h, \epsilon, \Sigma) = Q_\star(A', C', \Sigma)  < \infty$ and a minimizing gain by $L_\star$.
\end{definition}

\begin{proof}[Proof of Theorem \ref{theorem:main_nonlinear}]
    
Let $(f, h, x_0)$ be a nonlinear dynamical system satisfying Assumption \ref{assumption:nonlinear}, and let 
    $$z_{t+1} = A'z_t, \quad y_t' = C'z_t$$
    be the high-dimensional LDS constructed in Lemma \ref{lemma:discretization} at some scale $\epsilon$, where $z_0$ is the indicator vector of $\pi(x_0)$. With $\epsilon  = \frac{1}{T^{3/2}}$ this satisfies Assumption \ref{definition:linear_realizable} with $R_{x_0} = 1$ (since $z_0$ is a one hot probability vector), $W = 0$, the same bound $R$ on the signal, $R_B=0$ (since there are no inputs), and $R_C \leq R\sqrt{N} \leq R \cdot \left(2RT^{3/2}\right)^{d_\mathcal{X}/2}$. Letting $\hat{y}_1, \ldots, \hat{y}_T$ be the predictions made by Algorithm \ref{alg:clsf_regressionmain}, by Theorem \ref{theorem:main_linear}(b) (or more precisely by a regret phrasing, see Theorem \ref{theorem:agnosticlinear}(b)),
    $$\sum_{t=1}^T \norm{\hat{y}_t - y_t} \leq \sum_{t=1}^T\norm{y_t' - y_t} + O\left(Q_\star^2 \log(Q_\star) \cdot R m^2 \log\left((2RT^{3/2})^{d_\mathcal{{X}}}\right) \log^7(mT) \cdot \sqrt{T}\right)$$
    The result follows from the discretization error bound of Lemma \ref{lemma:discretization}.
\end{proof}
    

% \newpage
    
%     The only remaining subtlety is that the observer system constructed from the LDS of Lemma \ref{lemma:discretization} expects to receive the discretized observations $y_t'$ as inputs, whereas Algorithm \ref{alg:clsf_regressionmain} receives the true $y_t$. This injects a state error of size $\|L\| \cdot \|y_t' - y_t\|$ at each step, and with a small enough discretization $\epsilon$ this is tolerable.
    
    
%     \evannote{be careful about Frobenius!} with $R_\Theta = R\sqrt{d_\mathcal{Y}}$ and $\|y_t'\| \leq 2R$ \evannote{what about $\sigma$? and observability?}. We may construct its observer system via Definition \ref{prop:observersystem}, which is an LDS over $y_t'$ given by
%     $$\tilde{x}_{t+1} = \tilde{A} \tilde{x}_t + L_\star y_t', \quad \tilde{y}_t = C'\tilde{x}_t, \quad \tilde{x}_0 = z_0$$
%     with observations $\tilde{y}_t = y_t' = C'\tilde{A}^j z_0 + \sum_{j=1}^tC'\tilde{A}^j L_\star y'_{t-j}$. Since $\kappa_{\textrm{diag}}(\tilde{A}) = Q_\star(A', C', 0, \frac{1}{m})$, we see that the modified LDS over $y_t$ given by
%     $$\bar{x}_{t+1} = \tilde{A} \bar{x}_t + L_\star y_t, \quad \bar{y}_t = C'\bar{x}_t, \quad \bar{x}_0 = z_0$$
%     will satisfy $$\|\bar{y}_t - y_t'\| \leq \|C'\|  \cdot \|L_\star\| \cdot \sum_{j=1}^t \|y_j' - y_j\| \leq \frac{\|C'\| \cdot (1 + Q_\star)}{2\sigma_{\text{min}}(C')} \cdot t^2 \epsilon$$
%     Now, by Theorem \ref{theorem:main_linear}, Algorithm \ref{alg:clsf_regressionmain} generates predictions $\hat{y}_1, \ldots, \hat{y}_T$ with 
%     $$\sum_{t=1}^T \|\hat{y}_t - y_t\| \leq \sum_{t=1}^T \|y_t - \bar{y}_t\| + O\left(Q_\star S^2 R \sqrt{d_{\mathcal{Y}}} \cdot \log^3T \sqrt{T}\right)$$
%     By Lemma \ref{lemma:discretization} and the earlier bound summed over $t$, the first term is $\leq \frac{R\sqrt{d_\mathcal{Y}}(1+Q_\star)}{\sigma_{\text{min}}(C')} \cdot (T^2 + T^3) \epsilon \leq R\sqrt{d_\mathcal{Y}} S \cdot \sqrt{T}$, which can be absorbed into the second term.


\ignore{

\newpage
% In this section we prove a version of the Koopman operator which is sufficient for our purposes. It builds upon a particularly simple understanding of dynamical systems through discrete time and space discretization. 


\begin{theorem}[Existence and Properties of an Approximate Finite Linear System Matching Observations]
Let $(\mathcal{X}, f)$ be a deterministic discrete-time dynamical system with state $x_t \in \mathcal{X}$ and dynamics $x_{t+1} = f(x_t)$. Let $y_t = h(x_t)$ be the observation at time $t$, where $h : \mathcal{X} \to \mathcal{Y}$.

Assume:
\begin{itemize}
    \item[\textbf{(A1)}] (\textbf{Boundedness}) There exists $R > 0$ such that for all $t$,
    \[
    \|x_t\| \leq R, \quad \|y_t\| \leq R.
    \]
    \item[\textbf{(A2)}] (\textbf{Marginal stability / Non-expansiveness}) The map $F : \mathcal{Y} \times \mathcal{X} \to \mathcal{Y} \times \mathcal{X}$ defined by
    \[
    F(y, x) := (h(f(x)), f(x))
    \]
    is non-expansive:
    \[
    \|F(a) - F(b)\| \leq \|a-b\| \quad \text{for all } a, b \in \mathcal{Y} \times \mathcal{X}.
    \]
\end{itemize}

Let $d_{\mathcal{Y}}$ and $d_{\mathcal{X}}$ denote the (finite) dimensions of $\mathcal{Y}$ and $\mathcal{X}$, and set $D = d_{\mathcal{Y}} + d_{\mathcal{X}}$. Let $S$ be an $\epsilon$-net of the ball of radius $R$ in $\mathbb{R}^D$. The cardinality of $S$ (and thus the dimension $N$ of the finite LDS) satisfies
\[
N = |S| \leq \left(\frac{3R}{\epsilon}\right)^D.
\]

Then there exists a finite linear dynamical system (LDS) of dimension $N$, with state $z_t \in \mathbb{R}^N$ and a linear projection $C \in \mathbb{R}^{d_{\mathcal{Y}} \times N}$ with operator norm $\|C\| = 1$, such that the output sequence $\hat{y}_t = C z_t$ satisfies
\[
\|\hat{y}_t - y_t\| \leq t\epsilon
\]
for every $t$.
\end{theorem}

\begin{proof}
Let $\mathcal{Z} = \mathcal{Y} \times \mathcal{X}$ and $D = d_{\mathcal{Y}} + d_{\mathcal{X}}$. By assumption (A1), all pairs $(y_t, x_t)$ lie in the closed Euclidean ball of radius $R$ in $\mathbb{R}^D$. 

Let $S$ be an $\epsilon$-net for this ball, meaning that for every $(y, x)$ with $\|y\| \leq R$, $\|x\| \leq R$, there exists $s \in S$ such that $\|(y, x) - s\| \leq \epsilon$. It is standard that such an $\epsilon$-net can be chosen with 
\[
|S| \leq \left(\frac{3R}{\epsilon}\right)^D.
\]
We enumerate the elements of $S$ as $s^{(1)}, \ldots, s^{(N)}$.

Define the mapping $\pi: \mathcal{Z} \to S$ as the nearest neighbor projection, i.e., $\pi(z) = \operatorname{argmin}_{s \in S} \|z - s\|$.

Now define the (deterministic) finite system: For $s_t \in S$, set
\[
s_{t+1} := \pi(F(s_t)),
\]
where $F(y, x) = (h(f(x)), f(x))$ as in the assumptions. This defines a deterministic transition operator $T_S: S \to S$.

Let the finite-dimensional state at time $t$ be the indicator vector $z_t \in \mathbb{R}^N$ corresponding to $s_t$: that is, $z_t$ is the standard basis vector $e_k$ if $s_t = s^{(k)}$.

Define the (matrix) evolution operator $A \in \mathbb{R}^{N \times N}$ by $A_{ij} = 1$ if $T_S(s^{(j)}) = s^{(i)}$, else $A_{ij} = 0$. Then $z_{t+1} = Az_t$; thus, the evolution is linear.

Define the output projection $C \in \mathbb{R}^{d_{\mathcal{Y}} \times N}$ by $C e_k = y^{(k)}$, where $s^{(k)} = (y^{(k)}, x^{(k)})$. For any $z_t$, the output $\hat{y}_t = C z_t$ recovers the $y$ component of $s_t$. $C$ simply picks out the observation coordinate, and thus has operator norm $1$ (with respect to the Euclidean norm).

Now, let $(y_0, x_0), (y_1, x_1), \ldots$ be a trajectory of the original nonlinear system, and set $s_0 = \pi(y_0, x_0)$. Define $s_{t+1} = T_S(s_t)$ and $z_{t+1} = Az_t$ as above. Let $e_t = \|s_t - (y_t, x_t)\|$. By construction, $e_0 \leq \epsilon$.

We claim $e_{t+1} \leq e_t + \epsilon$ for all $t$. Indeed,
\begin{align*}
e_{t+1} &= \|s_{t+1} - (y_{t+1}, x_{t+1})\| \\
&= \|\pi(F(s_t)) - F(y_t, x_t)\| \\
&\leq \|F(s_t) - F(y_t, x_t)\| + \|\pi(F(s_t)) - F(s_t)\| \\
&\leq \|s_t - (y_t, x_t)\| + \epsilon \qquad \text{(non-expansiveness of $F$, definition of $\epsilon$-net)} \\
&= e_t + \epsilon.
\end{align*}
By induction, $e_t \leq t \epsilon$ for all $t$.

The actual output at time $t$ is $\hat{y}_t = C z_t$, corresponding to the $y$-component of $s_t$. Thus,
\[
\|\hat{y}_t - y_t\| = \|P(s_t) - y_t\| \leq \|s_t - (y_t, x_t)\| = e_t \leq t\epsilon,
\]
where $P: \mathcal{Y} \times \mathcal{X} \to \mathcal{Y}$ is the projection to the first $d_{\mathcal{Y}}$ coordinates.

Therefore, the claimed finite LDS with linear output projection $C$ approximates the true output sequence up to $t\epsilon$ at every time $t$, with the stated system dimension and projection norm.

\end{proof}
}



\ignore{

\begin{theorem}[Finite-dimensional Koopman Operator for Marginally Stable Systems]
Let $(\mathcal{X}, f)$ be a deterministic discrete-time dynamical system with state $x_t \in \mathcal{X}$ and dynamics $x_{t+1} = f(x_t)$. Let $y_t = h(x_t)$ be the observation at time $t$, where $h : \mathcal{X} \to \mathcal{Y}$.

Assume:
\begin{itemize}
    \item[\textbf{(A1)}] (\textbf{Finite-memory observability}) There exists $H \in \mathbb{N}$ and a function $\phi : \mathcal{Y}^H \to \mathcal{X}$ such that for all $t$,
    \[
    x_t = \phi(y_t, y_{t-1}, \ldots, y_{t-H+1}).
    \]
    \item[\textbf{(A2)}] (\textbf{Boundedness}) There exists $R > 0$ such that for all $t$,
    \[
    \|x_t\| \leq R, \quad \|y_t\| \leq R.
    \]
    \item[\textbf{(A3)}] (\textbf{Marginal stability (Non-expansiveness)}) The \textbf{lifted transition map}
    \[
    T(y_1, \ldots, y_H) := \left( h\left(f\left(\phi(y_1, \ldots, y_H)\right)\right),\ y_1, \ldots, y_{H-1} \right)
    \]
    is \textbf{non-expansive}, i.e.,
    \[
    \| T(a) - T(b) \| \leq \| a - b \| \quad \text{for all}\quad a, b \in \mathcal{Y}^H.
    \]
\end{itemize}

Let $\epsilon > 0$. Let $S$ be an $\epsilon$-net of the ball $\mathcal{B}_R^H \subset \mathcal{Y}^H$ (i.e., for every $(y_1, ..., y_H)$ with $\|y_i\| \leq R$, there exists $s \in S$ with $\|(y_1, ..., y_H) - s\| \leq \epsilon$).

Define the projection $\pi: \mathcal{Y}^H \to S$ as mapping each history to its nearest point in $S$.

Consider the induced finite system on $S$ defined by:
\[
s_{t+1} = \pi\left( T(s_t) \right)
\]
where $s_t = (y_t, y_{t-1}, ..., y_{t-H+1}) \in S$.

Let $T_S : S \to S$ denote this deterministic transition map.

For any observable $g : S \to \mathbb{R}$, define the finite Koopman operator $K$ by
\[
(Kg)(s) = g(T_S(s)),\quad s \in S.
\]

Suppose $g : \mathcal{Y}^H \to \mathbb{R}$ is $L_g$-Lipschitz.

Then, for any initial true history $h_0 \in \mathcal{Y}^H$ and its projection $s_0 = \pi(h_0)$, after $t$ steps:
\[
\left| (K^t g)(s_0) - (K_{\mathrm{true}}^t g)(h_0) \right| \leq L_g\, t\, \epsilon
\]
where $K_{\mathrm{true}}$ is the true (infinite-dimensional) Koopman operator acting on $g$.

\end{theorem}

\begin{proof}
\textbf{Step 1: Well-defined finite system.}\\
By (A2), all $x_t, y_t$ and thus all $H$-histories live in a bounded set, so the $\epsilon$-net $S$ is finite. The projected system on $S$ is deterministic and finite.

\textbf{Step 2: Non-expansive error propagation.}\\
Let $h_0 \in \mathcal{Y}^H$ be the true initial history and $s_0 = \pi(h_0)$ its projection. Let $h_t$ be the true $t$-step history and $s_t$ the discretized $t$-step history. Let $e_t := \|h_t - s_t\|$. Initially, $e_0 \leq \epsilon$.

At each step,
\begin{align*}
e_{t+1} &= \| T(h_t) - \pi(T(s_t)) \| \\
&\leq \| T(h_t) - T(s_t) \| + \| T(s_t) - \pi(T(s_t)) \| \\
&\leq \| h_t - s_t \| + \epsilon \\
&= e_t + \epsilon
\end{align*}

\textbf{Step 3: Induction.}\\
By induction, $e_t \leq t \epsilon$.

\textbf{Step 4: Observable error.}\\
Since $g$ is $L_g$-Lipschitz,
\[
|g(s_t) - g(h_t)| \leq L_g\, e_t \leq L_g\, t\, \epsilon.
\]
Thus,
\[
\left| (K^t g)(s_0) - (K_{\mathrm{true}}^t g)(h_0) \right| \leq L_g\, t\, \epsilon.
\]
This completes the proof.
\end{proof}
}

\section{Discussion}
\label{sec:discussion}

As a review of our nonlinear guarantee, we combined the Luenberger observer with a global linearization to construct a high-dimensional LDS with desirable spectral structure that competes against the original system, and showed that spectral filtering has a regret guarantee against this competitor. For a fixed nonlinear system, we defined the Luenberger program with value $Q_\star(\Sigma)$, which is the best possible condition number such that the constructed observer system has poles contained in $\Sigma$, and derived regret guarantees directly scaling with $Q_\star$.

\paragraph{Interpretation of loss bounds.} Suppressing dependence on $R$ and $d_\mathcal{X}$ for exposition, our sharpest analysis yields a regret guarantee of the form
$$\operatorname{Regret}_T(\mA, \Pi) \leq \widetilde{O}\left(\mathrm{poly} \log(Q_\star) \cdot \sqrt{T}\right)$$
for $\Pi$ being a class of marginally stable bounded nonlinear predictors,
where for $\mA$ we run Algorithm \ref{alg:clsf_regressionmain} with $O(\mathrm{poly} \log (Q_\star, T))$ many parameters and the Azoury-Vovk-Warmuth forecaster as our optimizer (see Theorem \ref{theorem:logdependence}). Furthermore, $Q_\star$ is exponential in the number of undesirable eigenvalues of \textit{some} high-dimensional LDS which approximates the nonlinear system -- we outline one construction via a Markov chain on a uniform discretization grid, but any other approximating LDS is just as valid. We find that if there is \textit{any} approximating LDS with $o(\log T)$ many undesirable eigenvalues, then we get an efficient algorithm with regret that is sublinear in $T$. 

Phrased in the other direction, if $N_\star$ denotes the optimal number of undesirable eigenvalues, minimized over all high-dimensional approximating LDS's, then it can be learned using a linear method such as spectral filtering with $\operatorname{poly}(N_\star)$ parameters and $\operatorname{poly}(N_\star)$ samples. We arrive at the main takeaway of our learning results: {if a nonlinear system has some high-dimensional approximating LDS with a certain spectral structure, then it can be efficiently learned by an improper algorithm suited for that structure}. 
With the choice of Algorithm \ref{alg:clsf_regressionmain} as our improper learning method, we get the statement that "\textit{the complexity of learning nonlinear signals via Algorithm \ref{alg:clsf_regressionmain} scales with the minmal number of marginally stable complex eigenvalues over all linear approximations}".

\paragraph{Spectral properties of linearizations.}
    Note that the spectral properties of Koopman operators depend on the choice of function space. Though we do not make this precise, by nature of improper learning our analysis implies that if there exists {\em any} function space on which the (truncation of the) Koopman operator satisfies the above desiderata then spectral filtering will succeed. 
    
    For many physical systems $f$ which satisfy some detailed balance condition (e.g. Langevin dynamics/energy minimization, see Section \ref{subsection:langevin}) the Koopman operator can be made self-adjoint, which yields $Q_\star(f, h, \epsilon, [-1, 1]) = O(1)$ for all $\epsilon$ and $h$ which are observable. 
    % We explore this connection between reversibility/time-symmetry and small $Q_\star$ in Appendix \ref{app:timesymm} \evannote{remove this last sentence if we cannot prove it, or if the proof is clean we should put it in main paper}.
    In other words, for many physical systems we truly do expect the existence of a high-dimensional approximating LDS with few marginally stable complex eigenvalues (see Section \ref{sec:numerics} for numerics regarding this). Combining this with the previous conclusion, we get the statement that "\textit{Algorithm \ref{alg:clsf_regressionmain} can learn physical systems provably and efficiently with few parameters}".

\paragraph{A control-theoretic perspective on learnability.} 
High-dimensional LDS's are in a sense universal dynamical computers, since by e.g. Lemma \ref{lemma:discretization} any reasonable nonlinear system can be approximated to arbitrary accuracy within this class. As such, it is very natural to define "learnability" of a nonlinear signal in terms of learnability of \textit{some} high-dimensional linear approximation. 

Armed with these guiding principles, a central theme in our analysis is the use of control-theoretic tools -- Luenberger observer systems and global linearization via Koopman-like methods -- not as components to be explicitly designed or learned, but to construct a theoretical benchmark in the form of a high-dimensional LDS predictor that an algorithm should aspire to match.
By design, $Q_\star(f, h, \Sigma)$ serves as a fundamental measure of the difficulty of learning a system $f$ from its observations $h$ using an improper algorithm that competes against LDS's with spectrum in $\Sigma$. 

This yields a notion of learnability that is tailored to a given improper algorithm.
For example, spectral filtering naturally competes against real-diagonalizable LDS's ($\Sigma \subset [-1, 1]$, i.e. physical systems with detailed balance) and regression naturally competes against stable LDS's ($\Sigma \subset \mathbb{D}_{1-\gamma}$, i.e. systems with an attractive fixed point), and Algorithm \ref{alg:clsf_regressionmain} demonstrates that we can compose these building blocks to cover a larger spectral region. The ability of these algorithms to learn a nonlinear signal depends sharply on the corresponding $Q_\star$ values of the generating system via our regret bounds, and our results give a way to determine the natural algorithm for a nonlinear system based on its qualitative properties (such as reversibility or strong stability). 

\paragraph{Summary.}
To summarize in a few words, the Luenberger program provides a new lens through which to analyze dynamical systems: by characterizing their $Q_\star$ values, one can make principled a priori judgments about their amenability to learning from partial observations with improper methods. Algorithm \ref{alg:clsf_regressionmain} is a very good general choice for such a method, since it can efficiently and provably learn to predict systems that are either physical or strongly stable. This methodology is universal, since the algorithm can be run without needing to know anything about the system in advance. 


% \subsection{A Control-Theoretic Perspective on Learnability}

% To establish a rigorous foundation for learning from observations, we first re-examine the problem through the lens of classical control theory. The central idea is to use the Luenberger observer—a fundamental tool for state estimation—not as a component to be explicitly designed, but as a theoretical benchmark that defines the best possible linear predictor an algorithm should aspire to match.

% \paragraph{The Luenberger Observer as a Learning Benchmark.}
% Consider a general discrete-time Linear Dynamical System (LDS) with process noise $w_t$, 
% $$ x_{t+1} = Ax_t + Bu_t +w_t, \quad y_t = Cx_t ,$$
% where $x_t \in \R^{d_h}$ is the hidden state, $u_t$ is the input, and $y_t$ is the observation. The Luenberger observer is a companion system designed to estimate the hidden state $x_t$ using only the known inputs $u_t$ and measured outputs $y_t$. The observer's state, $\tilde{x}_t$, evolves according to:
% $$ \tilde{x}_{t+1} = A \tilde{x}_t + Bu_t + L (y_t - \tilde{y}_t) ,$$
% where  $\tilde{y}_t  = C \tilde{x}_t$ is the estimated output and $L$ is the observer gain matrix. The term $L(y_t - \tilde{y}_t)$ acts as a correction, driving the estimated state $\tilde{x}_t$  towards the true state $x_t$ by feeding back the observation error. 

% The dynamics of the state estimation error, $\xi_{t} = x_t - \tilde{x}_t$, reveal the mechanism of the observer. Subtracting the observer equation from the system equation yields:
% $$ \xi_{t+1} = (A-LC) \xi_t + w_t . $$
% This equation shows that the error dynamics are governed by the matrix $A-LC$. If the eigenvalues of this matrix are all within the unit circle then the system is stable, and the effect of any initial estimation error $\xi_0$ will decay to zero over time. In the presence of persistent noise $w_t$, a stable observer will produce a bounded tracking error. Consequently, a stable Luenberger observer, driven by the real inputs $u_t$ and outputs $y_t$, serves as a powerful predictor for the system's behavior. This provides our first key insight: the performance of the best possible linear predictor for a given LDS is determined by the properties of the best possible stable observer for that system.

% \paragraph{Robust Pole Placement and the Luenberger Program.} 
% The crucial task in designing an effective observer is the selection of the gain matrix $L$. A fundamental result in control theory is that if the pair $(A,C)$ is observable—meaning the hidden state $x_t$ 
% can be uniquely determined from a sequence of outputs—then the gain $L$ can be chosen to place the eigenvalues of $A - LC$ at any desired locations in the complex plane. This technique is known as pole placement, where "poles" are synonymous with the system's eigenvalues.

% However, simply placing the eigenvalues in a stable region is not sufficient for a robust learning benchmark. An observer whose stability is fragile—meaning tiny perturbations to the system matrices $A$ or $C$ could render it unstable—is a poor benchmark. This motivates the concept of robust pole placement, where the goal is not only to achieve stability but to do so in a well-conditioned manner. We formalize this objective with the \textit{Luenberger program}, an optimization problem that seeks the best observer with desirable spectral structure in terms of its conditioning.  

% \begin{definition}
% Let $A \in \R^{d_h \times d_h}$ and $ C \in \R^{d_{out}\times d_h}$. For each spectral constraint set $\Sigma \subset \mathbb{D}_1$,  we consider the optimization problem over $L \in \R^{d_{h} \times d_{out}}$  given by:
% $$\min \quad \kappa_{\textrm{diag}}(A-LC) \quad \text{s.t.} \quad A-LC \text{ has all eigenvalues in $\Sigma$}.$$
% When the constraint set is nonempty, we denote the minimal value by $Q_\star = Q(A,C,\Sigma) < \infty$.
% \end{definition}
% Here, $\kappa_{\operatorname{diag}}(M) = \|H\| \cdot \|H^{-1} \| $ is the spectral condition number of the matrix $M$, where $H$ is the matrix of eigenvectors that diagonalizes $M$; this number measures the sensitivity of the eigenvalues of
% $M$ to perturbations. A value close to 1 (its minimum for normal matrices) indicates that the eigenvalues are robust, while a large value indicates that they are sensitive and the problem is ill-conditioned. Well-chosen constraints can ensure that the observer is stable (and also with real eigenvalues, which will be crucial for our spectral filtering approach), but if the original system eigenvalues are significantly far away from $\Sigma$ then it can require large $\|L\|$ and $Q_\star$.


% \paragraph{Quantifying Learning Difficulty: The Dynamics Condition Number $Q_\star$.}
% Consider a learning algorithm which excels at learning systems with a certain spectral structure (such as stability for system identification).
% The optimal value of the Luenberger program, $Q_\star$, serves as a fundamental measure of the difficulty of learning a system from its observations using such an algorithm. It quantifies the conditioning of the best possible observer system with given spectral structure, which has important implications for any learning algorithm. There is much flexibility in {\it where} to place the poles, and while we advocate for some choices that give our theoretical results, optimal placement is far from solved and remains an important research direction. 


% % A large $Q_\star$ signifies that the pole placement problem is ill-conditioned. This has severe consequences for any learning pipeline that relies on explicit system identification. 

% %Such a pipeline would first estimate the matrices
% %$\hat{A}$ and $\hat{C}$ from data, and then compute an observer gain $L$ based on these estimates. However, when $Q_\star$ is large, even minuscule estimation errors can lead to enormous changes in the required gain $L$ and, more critically, in the resulting eigenvalues of $\hat{A}- L \hat{C}$. A learning algorithm might therefore produce an observer that it believes to be stable, but which is in fact unstable when applied to the true system.

% In general, the value of $Q_\star$ can be very large, as we formally prove in the appendix. This brittleness is not just a theoretical concern. As we demonstrate later, for a simple and common dynamical structure—the cyclical permutation matrix—the value of $Q_\star$ is provably exponential in the hidden dimension, i.e., $Q_\star = \Omega(2^{d_h})$. Such a system is perfectly observable in the classical sense, yet it is hard to learn robustly via standard observer design due to its extreme ill-conditioning.
% This establishes a quantitative and continuous link between the control-theoretic notion of observability and the machine-learning notion of learnability. 

% Classical observability is a binary property: a system is either observable or not. If it is, a gain $L$ exists to place the poles arbitrarily. However, this provides no information about the practical feasibility of finding and using such a gain. The condition number
% $Q_\star$ measures this property quantitatively. A system with a small $Q_\star$ can be considered "robustly observable," meaning a desirable (stable, perhaps with real eigenvalues) observer is easy to find and insensitive to small errors. A system with a large $Q_\star$ is "brittle," being theoretically observable but practically very difficult to learn. 

% As the main theorems of this paper will show, the regret of our learning algorithm scales directly with $Q_\star$. Thus, $Q_\star$ acts as a bridge: a high degree of robust observability (low $Q_\star$) translates directly to a high degree of learnability (low regret). This provides a new lens through which to analyze dynamical systems: by characterizing their $Q_\star$ values, one can make principled, a priori judgments about their amenability to learning from partial observations.

% \paragraph{Time-Symmetric Dynamical Systems}
% An important class of dynamical systems which admits practical learnability via our methods is that of time-symmetric nonlinear dynamical systems. Every physcial system derived from Newtonian mechanics, if we exclude friction and noise, is by its mathematical nature time-reversible. This admits a symmetric Koopman lifting as we prove in Appendix \ref{sec:symmetric-koopman}, and thus poles that naturally lie in the interval $[-1,1]$. The observer system we construct can therefore be identical to the true system, and as a consequence we can learn it efficiently using spectral filtering with $Q_\star$ that is a small constant. 

\section{Experiments}\label{sec:experiments}

In this section, we confirm our theoretical results with various experiments. First, we demonstrate that spectral filtering over observations allows for the learning of asymmetric LDS with process noise. Then, we investigate the algorithm's performance on several nonlinear sequence prediction tasks. Lastly, we study the numerical behavior of $Q_\star$ on some systems of interest.

\subsection{Linear Systems}


Recall Algorithm \ref{alg:clsf_regressionmain}, which performs spectral filtering over the observation history $y_t$ as well as the open-loop controls $u_t$. If we run the algorithm with $N^t \equiv 0$ and $m=1$ then we essentially recover the original vanilla spectral filtering algorithm of \cite{hazan2017learning}, which has a guarantee for symmetric LDS with no process noise. We will refer to this baseline as \verb|SF|, and we will denote the full Algorithm \ref{alg:clsf_regressionmain} by \verb|SF+obs|. For simplicity, we will only compare these two methods to highlight the contributions of this work; for comparisons between spectral filtering-based algorithms and other LDS learning methods, please see \cite{liu2024flash} or \cite{shah2025spectraldsprovabledistillationlinear}. 

Theorem \ref{theorem:main_linear} proves two main benefits of the additional filtering over observations: (1) robustness to adversarial process noise and (2) the ability to learn asymmetric LDS. To demonstrate these advancements, we consider signals $(y_t)_{t=1}^T$ generated by an LDS according to (\ref{eqn:LDS}) in two settings:
\begin{enumerate}
    \item[(1)]$A$ is a Gaussian random $128 \times 128$
    matrix, normalized for $\|A\| = 1$, and $w_t = 0.1 \cdot \sin(3\pi t) \cdot [1,\ldots, 1] \in \R^{128}$ are correlated sinusoidal disturbances.
    \item[(2)] $A$ is the $16 \times 16$ cyclical permutation matrix and $w_t = 0$.
\end{enumerate}
In both cases, $d_{in}=d_{out}=1$, we sample $B, C, x_0, (u_t)_{t=1}^T$ as i.i.d. standard Gaussians, we use the \verb|optax.contrib.cocob| optimizer for simplicity\footnote{Since the involved optimizations are convex, these results are robust to optimizer choice. By contrast, learning either of the above problems with another approach such as S4 or Mamba are nonconvex and very dependent on the optimizer used.}, and we run Algorithm \ref{alg:clsf_regressionmain} with $m = 1$ and $h=24$. The instantaneous losses are plotted in Figure \ref{fig:experiments_lds}, averaged over random seeds and smoothed.

\begin{figure}[!htb]
    \centering
    \includegraphics[width=0.45\linewidth]{figs/lds_online_1.png}
    \includegraphics[width=0.45\linewidth]{figs/lds_online_2.png}
    \caption{Instantaneous losses $\|\hat{y}_t - y_t\|^2$ plotted as a function of $t$ for experiments (1) and (2), averaged over 12 random seeds with a smoothing filter of length 100 and 1000, respectively. 
    % Predictions $\hat{y}_t$ are generated using Algorithm \ref{alg:clsf_regressionmain}, and we compare to the original spectral filtering baseline of \cite{hazan2017learning}.
    }
\label{fig:experiments_lds}
\end{figure}

As expected, filtering over observations improves the tolerance to asymmetry and correlated process noise, and on setup (1) the minimal loss achieved is a fixed determinstic nonzero value determined by the noise structure. On setup (2), we see that vanilla spectral filtering is unable to make any progress for the permutation system (which is the most asymmetric marginally stable LDS); by contrast, the simple addition of filtering over observations enables learning. Furthermore, in (2) the number of steps needed to achieve 0 loss is very large\footnote{Interestingly, for the permutation system there is an initial phase of rapid decrease, followed by a very slow second phase of refinement. Anecdotally, we find that for the first order optimizers we experimented with, the length of the initial phase appears independent of $d_h$ whereas the length of the second phase is governed by $Q_\star$. This indicates that perhaps other optimization algorithms would achieve faster convergence, which we leave to future work.}, as predicted by the value of $Q_\star$ (see Example \ref{example:permutationqstar}). Note that even though the systems in (1) are asymmetric, it is in a much tamer way that results in smaller $Q_\star$ and faster learning. 
To sum up, the benefits of filtering over observations appear exactly as predicted by Theorem \ref{theorem:main_linear}.




\subsection{Nonlinear Systems}
We now turn to autonomous nonlinear systems, which generate a signal $(y_t)_{t=1}^T$ according to (\ref{eqn:NLDS}). The task is still next-token prediction, in which we use $y_1, \ldots, y_t$ to form predictions $\hat{y}_{t+1}$. In Section \ref{sec:experiments_nonlinearalgs} we will detail the four algorithms we will run, after which we will present their empirical performance on simple nonlinear systems of interest.

\subsubsection{Methods}
\label{sec:experiments_nonlinearalgs}
Consider running Algorithm \ref{alg:clsf_regressionmain} with $J_t,M_t \equiv 0$, which corresponds to spectral filtering with regression over the observation history. As before, we run this with $k=24$ and $m=1$ (yielding $(k+m)\cdot d_\mathcal{Y}^2$ many parameters), and for optimizer we use either COCOB or ONS with tuned learning rate, depending on computational restrictions. For the following experimental results, we refer to this as \verb|SF|. As baselines, we consider:
\begin{itemize}
    \item Directly learning an observer LDS $(A, B, C, x_0)$ via online gradient descent on the next-token MSE loss, with predictions of the form
    $$\hat{y}_{t+1} = \sum_{j=0}^t CA^{j}By_{t-j} + CA^{t+1}x_0.$$
    We refer to this baseline as \verb|LDS|, where we fix $d_h=64$ to be the hidden dimension of the LDS. This predictor has $d_h^2 + 2d_h\cdot d_\mathcal{Y} + d_h$ parameters, and the involved optimization is nonconvex\footnote{In fact, the \verb|LDS| baseline tends to blow up unless trained with the Adam optimizer with tuned learning rate, which we use for this baseline throughout the nonlinear experiments.}.
    
    \item Performing an online variant of the eDMD algorithm \cite{eDMDrowley}, which fixes a dictionary $\{f_j\}_{j=1}^n$ of nonlinear observables $f_j: \mathcal{Y} \to \R$, defines the lifted states $z_t = [f_1(y_t), \ldots, f_n(y_t)] \in \R^n$, and attempts to fit a predictor $z_{t+1} = Az_t$, $\hat{y}_t = Cz_t$ for matrices $A \in \R^{n \times n}$, $C \in \R^{d_\mathcal{Y} \times n}$. We use the \verb|pykoopman| library \cite{pan2023pykoopmanpythonpackagedatadriven} implementation, which fits $A, C$ using least squares linear regression over the history of past observations. 
    For the choice of observables, we let $n = d_\mathcal{Y} + 20$, where $$f_j(y) = \begin{cases}
        y_j & j = 1, \ldots, d_{\mathcal{Y}} \\
        \mathrm{RBF}_j(y) & j = d_{\mathcal{Y}} + 1, \ldots, n
    \end{cases},$$ 
    where $\mathrm{RBF}_j$ refers to a radial basis function using thinplate splines with centers fitted to the data\footnote{This choice is described in Section 3 of \cite{eDMDrowley}, and is also used in the eDMD tutorials for the \verb|pykoopman| library, \href{https://pykoopman.readthedocs.io/en/master/tutorial_koopman_edmd_with_rbf.html#Koopman-regression-using-EDMD}{available online}. Furthermore, this dictionary seemed to be a good choice uniformly over all our tasks.}. In other words, our nonlinear dictionary concatenates the initial observations with 20 $\mathrm{RBF}$ observables. We refer to this baseline as \verb|eDMD|, and fit every $T/10$ steps for efficiency. This method has $n^2 + n\cdot d_\mathcal{Y}$ parameters.
    
    \item Lastly, we introduce a new method of eDMD combined with spectral filtering, which we call \verb|SFeDMD|. We fix the same set of observables as in the \verb|eDMD| baseline, but instead of fitting $A, C$ directly we apply the predictor from Algorithm \ref{alg:clsf_regressionmain} (with $k=24$, $m=0$) to learn the lifted dynamics. One could learn the parameters $N^t$ with either gradient descent or least squares linear regression: we choose the latter to have a more direct comparison against the \verb|eDMD| baseline. This method has $(k+m)\cdot n \cdot d_\mathcal{Y}$ parameters.
\end{itemize}

To summarize, the algorithms we run consist of \verb|SF| (which directly applies the observation-only version of Algorithm \ref{alg:clsf_regressionmain} using gradient descent), \verb|LDS| (which attempts to directly learn an observer LDS via gradient descent), \verb|eDMD| (which performs the classical eDMD algorithm via linear regression with a fixed dictionary of nonlinear observables), and \verb|SFeDMD| (which uses the same dictionary of nonlinear observables and learns the dynamics improperly via spectral filtering and linear regression). These four algorithms constitute both proper and improper learning methods on both direct observations and nonlinear liftings. Rather than using online optimization algorithms for all methods, we choose to use regression for both eDMD baselines in order to directly apply the \verb|pykoopman| library and for more direct comparison against eDMD experimental results of other works.



\subsubsection{Lorenz System}

The Lorenz system is a 3-dimensional deterministic nonlinear dynamical system with state $x = [x^{(1)}, x^{(2)}, x^{(3)}] \in \R^3$ evolving according to the coupled differential equations\footnote{We simulate this ODE using a simple explicit Euler discretization with $\Delta t = 0.01$ to generate our plots and datasets.}
$$\begin{cases}
    \dot{x}^{(1)} = \sigma(x^{(2)} - x^{(1)}), \\
    \dot{x}^{(2)} = x^{(1)}(\rho - x^{(3)}) - x^{(2)}, \\
    \dot{x}^{(3)} = x^{(1)} x^{(2)} - \beta x^{(3)}
    \end{cases}$$
    for fixed constants $\sigma, \rho, \beta$. This was first considered by Edward Lorenz in 1963 as a simplified model for atmospheric dynamics, and since then it has become
a standard dynamical system on which to test prediction algorithms. For the choices $\sigma = 10, \rho = 28, \beta = \frac{8}{3}$ (which we make from now on) the dynamics become quite complicated, with chaotic solutions and the existence of a strange attractor of fractal dimension. Two typical trajectories are plotted in Figure \ref{fig:typicallorenztrajectory}, demonstrating the characteristic two-lobed attractor shape as well as the sensitive dependence on initial condition.


\begin{figure}[!htb]
    \centering
    \includegraphics[width=0.45\linewidth]{figs/lorenztrajectory1.png}
    \includegraphics[width=0.45\linewidth]{figs/lorenztrajectory2.png}
    \caption{Two trajectories of the Lorenz system, run for $1,024$ steps starting from the initial conditions $[1, 1, 1]$ and $[1.1, 1, 0.9]$, respectively. Initial positions are marked with a red star. These two trajectories quickly diverge from each other despite their similar initial conditions, demonstrating the chaotic behavior.
    % Predictions $\hat{y}_t$ are generated using Algorithm \ref{alg:clsf_regressionmain}, and we compare to the original spectral filtering baseline of \cite{hazan2017learning}.
    }
\label{fig:typicallorenztrajectory}
\end{figure}

For each fixed initial condition $x_0$, the Lorenz system produces a sequence of states $x_1, \ldots, x_T$. We consider two prediction tasks: one in which the algorithm sees full observations $y_t = x_t$, and one in which the algorithm sees partial observations via $y_t = Cx_t$ for $C \in \R^{1 \times 3}$. We run the four algorithms from Section \ref{sec:experiments_nonlinearalgs} on both the fully and partially observed versions of the task, with initial conditions and the value of $C$ i.i.d. standard Gaussian. Losses are plotted in Figure \ref{fig:experiments_lorenz}, averaged over random seeds and smoothed.


\begin{figure}[!htb]
    \centering
    \includegraphics[width=0.45\linewidth]{figs/lorenz_losses.png}
    \includegraphics[width=0.45\linewidth]{figs/lorenzpo_losses.png}
    \caption{Instantaneous losses $\norm{\hat{y}_t - y_t}^2$ plotted on log scale as a function of $t$ for the fully and partially observed Lorenz system, respectively, averaged over 12 random seeds with a smoothing filter of length 1000.
    % Predictions $\hat{y}_t$ are generated using Algorithm \ref{alg:clsf_regressionmain}, and we compare to the original spectral filtering baseline of \cite{hazan2017learning}.
    }
\label{fig:experiments_lorenz}
\end{figure}

As we see, Algorithm \ref{alg:clsf_regressionmain} is able to learn successfully in either the partial or full observation setting when filtering over either the original or nonlinearly-augmented observations -- this corroborates the conclusion of Theorem \ref{theorem:main_nonlinear}. By contrast, attempting to learn an observer system directly suffers from instability and nonconvexity, making it unfeasible even on this fairly low-dimensional task. Furthermore, we see that spectral filtering over a sequence augmented by nonlinear observables is certainly the best way to go: the \verb|SFeDMD| method is able to outperform \verb|eDMD| with fewer parameters in the fully observed setting, despite the fact that they both make use of the same dictionary of observables. In addition, a filtering-based predictor is able to aggregate information across time in a way that can handle partial observations naturally, whereas \verb|eDMD| cannot\footnote{The classical eDMD approach to handle partial observation would be to use time-delayed observables to generate an injective representation of state via Takens' theorem, but (1) anecdotally this is easier said than done since the dynamics of a time-delay lifting are often much more complex and (2) this vastly increases the number of parameters and the complexity of gradient-based optimization.}. 

So far, we have investigated the performance of various methods on the task of next-token prediction; after all, this is the loss on which we are training and the regret metric that we use for our theory. While this is by far the most common objective in practice, predictors learned under this loss are often deployed in an autoregressive way\footnote{By autoregression, we mean that there is an initial context $y_1, \ldots, y_L$, which the model uses to produce a prediction $\hat{y}_{L+1}$. In the next iteration, the model uses the context $y_1, \ldots, y_L, \hat{y}_{L+1}$ to predict $\hat{y}_{L+2}$, which in turn gets appended to the context and the process is repeated. In this way, the predictor may be queried with contexts 
corrupted by its prior mistakes, which allows
errors to compound.}, such as for text generation in LLMs or rollouts of world models in MPC-type control and planning applications. While our theory does not capture this type of generalization (a predictor which gets very small next-token prediction loss may fail catastrophically in rollout), we find it interesting to inspect the autoregressive behaviors of our learned predictors. 

\begin{figure}[!htb]
    \centering
    \includegraphics[width=0.45\linewidth]{figs/lorenz_rollouts1.png}
    \includegraphics[width=0.45\linewidth]{figs/lorenz_rollouts2.png}
    \caption{Two sets of autoregressive trajectories of 
    length 512, plotted alongside the ground truth trajectory given by continuing to simulate the Lorenz ODE. The initial positions at which the rollouts start are marked with a red star.}
\label{fig:lorenz_autoregressive}
\end{figure}

In Figure \ref{fig:lorenz_autoregressive} we plot two autoregressive trajectories of the learned \verb|eDMD| and \verb|SFeDMD| predictors\footnote{The \verb|LDS| and \verb|SF| predictors consistently diverge during rollouts, which may be due to their online training via gradient descent. At each iteration of training the predictor is adapted to the current state and dynamics, and so it may forget about behaviors that it saw earlier in training. For this reason we suspect that training in batch on a dataset of trajectories may fix this.} started from contexts unseen during training. On the first one, both methods are able to follow the general trend of the dynamics during the whole rollout, though \verb|SFeDMD|'s predictions remain sharp for longer. On the second trajectory, the errors made by \verb|eDMD| combine with the chaotic nature of this system to produce a qualitatively different rollout, which \verb|SFeDMD| is able to avoid. Of course, the ability of a predictor to generalize from next-token training to autoregressive rollout is heavily dependent on the parameterization, and deep models or different nonlinear liftings likely produce different behaviors. Since autoregressive rollout is not the goal of this paper, we leave a more detailed investigation to future work.




\subsubsection{Double Pendulum}
Another nonlinear system of physical interest is the double pendulum (a pendulum with two joints), which also exhibits a chaotic behavior and can have complicated marginally stable dynamics. We use the implementation in the Brax \cite{freeman2021braxdifferentiablephysics} library, which is a \verb|jax|-based physics simulator -- this double pendulum is attached to a cart that can move horizontally on a rail, and it takes inputs which push the cart to the left or right. A fixed deterministic agent/policy operating in such an environment yields a deterministic nonlinear dynamical system which produces a sequence of states $x_0, \ldots, x_T$; for this environment, we define the state to be (the sines and cosines of) the angles and the angular velocities of both joints, together with the position and velocity of the cart to total $d_\mathcal{X} = 8$. We consider a fixed linear feedback controller of the form $u_t = Kx_t$ for a fixed $K \in \R^{1 \times 8}$ with i.i.d. standard Gaussian entries.
% \footnote{In truth, we only play nonzero controls every 5 steps, since otherwise the simulator has trouble keeping up and returns \verb|nan|.}
As before, we can construct the sequence prediction task with either full observations $y_t = x_t$ or partial observations $y_t = Cx_t$ for a fixed $C \in \R^{1 \times 8}$ with i.i.d. standard Gaussian entries. Losses for both settings are plotted in Figure \ref{fig:experiments_pendulum}, averaged over random seeds and smoothed.


\begin{figure}[!htb]
    \centering
    \includegraphics[width=0.45\linewidth]{figs/pendulum_losses.png}
    \includegraphics[width=0.45\linewidth]{figs/pendulumpo_losses.png}
    \caption{Instantaneous losses $\norm{\hat{y}_t - y_t}^2$ plotted on log scale as a function of $t$ for the fully and partially observed double pendulum, respectively, averaged over 12 random seeds with a smoothing filter of length 1000.
    % Predictions $\hat{y}_t$ are generated using Algorithm \ref{alg:clsf_regressionmain}, and we compare to the original spectral filtering baseline of \cite{hazan2017learning}.
    }
\label{fig:experiments_pendulum}
\end{figure}

This is a more difficult system than the Lorenz attractor (since a predictor needs to also learn the effects of $K$). 
Despite this, we see a similar story as in the Lorenz system: with full observations both \verb|eDMD| and \verb|SF| are able to predict well, and \verb|SFeDMD| is able to predict very well. Furthermore, \verb|SF| and \verb|SFeDMD| are able to succeed when given partial observations, which \verb|eDMD| struggles with. The \verb|LDS| baseline is brittle -- on some trajectories it achieves the same performance as \verb|SF|, but on others it can completely diverge even at reasonable learning rates.


\begin{figure}[H]
    \centering
    \includegraphics[width=0.45\linewidth]{figs/pendulum_preds.png}
    \includegraphics[width=0.45\linewidth]{figs/pendulumpo_preds.png}
    \caption{Predictions $\hat{y}_t$ of the four algorithms along with the ground truth $y_t$ plotted as a function of $t$ for the fully and partially observed double pendulum, respectively, on a single random seed. We zoom in on the last 100 steps of the trajectory. For the full observation case, since we cannot plot the full state we choose to plot the sine of the angle between the two arms of the pendulum, which is a representative coordinate for the task of next-state prediction. 
    }
\label{fig:pendulum_preds}
\end{figure}


To add a bit more clarity to the behaviors of these algorithms, we plot in Figure \ref{fig:pendulum_preds} the actual predictions $\hat{y}_t$ made during training at the end of a typical trajectory, along with the ground truth signal $y_t$. We see in the full observation setting that the pendulum is swinging irregularly in a way that the models can predict with reasonable accuracy (except for \verb|LDS|, which suffers from unstable optimization). In the partial observation setting the algorithms all do quite well, but we do get to see some differences between the algorithms -- \verb|eDMD| is unable to be very precise due to the partial observation setting and appears to be predicting $\hat{y}_{t+1} = y_t$, whereas \verb|SFeDMD| produces very precise predictions. 


\subsubsection{Langevin Dynamics}
\label{subsection:langevin}
As a final autonomous nonlinear system, we will consider stochastic gradient flow/Gibbs energy minimization/Langevin dynamics. For a clear and in-depth description of this dynamical system from a stochastic differential equation (SDE) perspective, see Section 1 of the textbook \cite{chewi2025logconcave}. Given a differentiable loss/energy/potential function $V: \R^{d_\mathcal{X}} \to \R$, one can define a dynamical system according to the SDE
$$dX_t = -\nabla V(X_t) dt + \sqrt{2}dW_t,\quad X_0 = x_0,$$
which produces a random sequence of states $X_1, \ldots, X_T$ (we use capital letters here to denote random variables). This is very common in the physical sciences and more recently in machine learning, and can be interpreted through many different perspectives:
\begin{itemize}
    \item A first perspective is to interpret this as a \textbf{gradient descent}: we can imagine a discretized version via\footnote{This simple integrator with $\eta = 0.01$ is how we will numerically simulate SDEs for our experiments.} $$X_{t+1} = X_t - \eta \nabla V(X_t) + \sqrt{2\eta}w_t,$$ where $\eta$ is a small step size and $w_t \sim \mathcal{N}(0, 1)$ is an independent noise term that can be viewed as a stochastic gradient. In this way, the states correspond to the iterates of a first-order optimization algorithm.
    \item Under very mild assumptions, the distribution of $X_T$ for very large $T$ will approach the probability distribution $\pi$ on $\R^{d_\mathcal{X}}$ with (unnormalized) density $e^{-V}$. A particular trajectory can be viewed as a way to transform $x_0$ into a sample from $\pi \propto e^{-V}$, which yields a \textbf{sampling} perspective.
    \item Langevin diffusions arise naturally in thermodynamics and statistical physics, in which case $V$ is the energy of a particle/configuration -- it is known that the equilibrium distribution (at least at temperature 1) will be the stationary distribution $\pi \propto e^{-V}$, which is sometimes also called the Gibbs distribution. The evolution of the physical system will follow these trajectories, yielding a \textbf{statistical physics} perspective.
\end{itemize}
In any case, it is clear that prediction of such systems is a very useful task. The main difficulty is this: it can take very long for the system to equilibriate, and the transient/nonequilibrium dynamics are often the most useful (predicting optimization trajectories becomes useless if we wait until convergence). For some intuition, the Koopman operator\footnote{Technically one would define the Koopman operator as a family $K_t$ sending $h(x) \to \mathbb{E}[h(X_t) \mid X_0 = x]$, and the derivative of $K_t$ is the generator of the SDE. The main story is the same -- approximating the nonlinear dynamics by infinite-dimensional linear ones.} may have multiple marginally stable modes, in which case there is slow convergence (i.e. no strong mixing), and in such cases prediction is difficult. Under extra assumptions (strong convexity of $V$, a Poincare inequality for $\pi$, etc.) we would have exponentially-fast convergence mixing, but the general case is most interesting.

A crucial property of such SDEs is that they satisfy a detailed balance condition, and so the Koopman operator is a self-adjoint operator on $L^2(\pi)$ (see Remark 1 of \cite{kostic2022learningdynamicalsystemskoopman}). In our language, this implies that \textbf{there exists a well-conditioned high-dimensional symmetric LDS approximating the dynamics} -- this is very good news for us, since it implies a small $Q_\star$ and so strong performance of Algorithm \ref{alg:clsf_regressionmain} due to Theorem \ref{theorem:main_nonlinear}. 

To make the prediction problem difficult, we consider a potential on $\R^{d_\mathcal{X}}$ with exponentially-many asymmetric wells along each coordinate and a nontrivial coupling between coordinates, given by
$$V(x) = \sum_{j=1}^{d_\mathcal{X}} \left(0.05 \cdot x_j^4 - x_j^2 + 0.1 \cdot x_j\right) - 0.2 \cdot \sum_{1 \leq i < j \leq d_{\mathcal{X}}}x_i^2 x_j^2.$$



\begin{figure}[!htb]
    \centering
    \includegraphics[width=0.35\linewidth]{figs/potential1d.png}
    \includegraphics[width=0.35\linewidth]{figs/potential2d.png}
    \caption{Densities of the stationary distribution $\pi$ corresponding to the chosen potential $V$ for $d_\mathcal{X}$ equal to 1 and 2, respectively. We see a number of asymmetric wells that is exponential in the dimension.
    }
\label{fig:langevinpotential}
\end{figure}

For convenience, we plot the density of the stationary distribution $\pi \propto e^{-V}$ in one and two dimensions in Figure \ref{fig:langevinpotential}. In Figure \ref{fig:langevin_losses} we plot the training losses for the $d_\mathcal{X} = 64$ instantiation of the problem, averaged over random seeds and smoothed -- we find that the \verb|LDS| method consistently fails in the same way, suggesting some obstacle for learning an observer system via gradient descent on this task. By contrast, the other methods do well, and in this setting we find that the addition of the nonlinear dictionary offers no help. 
The minimal loss value reached is consistent with the size of the noise at each step (which is typically of size 0.02), agreeing with the lower bound of Theorem \ref{theorem:lowerbound}. We anticipate more interesting behaviors for more complicated Langevin dynamics, experimentation on which we leave for future work.

\begin{figure}[H]
    \centering
    \includegraphics[width=0.45\linewidth]{figs/langevin_losses.png}

    \caption{Instantaneous losses $\norm{\hat{y}_t - y_t}^2$ plotted on log scale as a function of $t$ for the Langevin dynamics, averaged over 12 random seeds with a smoothing filter of length 200.
    % Predictions $\hat{y}_t$ are generated using Algorithm \ref{alg:clsf_regressionmain}, and we compare to the original spectral filtering baseline of \cite{hazan2017learning}.
    }
\label{fig:langevin_losses}
\end{figure}


 % \evannote{add experiment with the many-wells $V$ -- show both loss plots and pred plots}


\subsection{Numerical Tests}
\label{sec:numerics}

Thus far, we have demonstrated the behaviors of Algorithm \ref{alg:clsf_regressionmain} on various nonlinear systems and compared to some simple baselines; we saw that spectral filtering over observations seems to succeed without much issue, and without the very slow optimization behaviors that were present in the permutation system (recall Figure \ref{fig:experiments_lds}). There is the intuition developed in the previous section that physical systems ought to have (approximately) real Koopman spectrum and therefore should be amenable to learning via spectral filtering. However, it would be nice to probe this question numerically, which we will quickly do in this section.

A classical approach to estimating Koopman eigenvalues of a nonlinear system is to run the eDMD algorithm as we have been doing, in which one learns a linear dynamics $A$ on a lifted space. By inspecting the eigenvalues of $A$, we gain understanding of the structure of (a finite-dimensional approximation to) the Koopman operator\footnote{Technically, this spectral relationship is only proven under the assumption that the dictionary of nonlinear observables spans an invariant subspace of the Koopman operator. This assumption is required for most eDMD-based methods to work at all, and inspecting the eigenvalues of the eDMD model is commonplace in practice regardless. See Section 5.1 of \cite{brunton2021modernkoopmantheorydynamical}.}. We take the \verb|eDMD| models\footnote{One could attempt to do this with the observer systems learned directly via gradient descent in the \verb|LDS| method. These appear to be placed generically in the unit disk, though this result is inconsistent and brittle, varying significantly across seeds. By contrast, \verb|eDMD| learns the same system eigenvalues across seeds, which are the ones we report.} trained on the full observation versions of our three nonlinear systems and plot the learned eigenvalues in Figure \ref{fig:edmdeigvals}. As we can see, on these marginally stable physical systems the Koopman eigenvalues appear as expected: approximately real and near 1. 

\begin{figure}[ht] % <---
   \begin{subfigure}{0.32\textwidth}
       \includegraphics[width=\linewidth]{figs/lorenz_edmdeigvals.png}
   \end{subfigure}
\hfill % <--- 
   \begin{subfigure}{0.32\textwidth}
       \includegraphics[width=\linewidth]{figs/pendulum_edmdeigvals.png}
   \end{subfigure}
\hfill % <---
   \begin{subfigure}{0.32\textwidth}
       \includegraphics[width=\linewidth]{figs/langevin_edmdeigvals.png}
   \end{subfigure}

       \caption{Eigenvalues of the lifted linear dynamics learned by the eDMD algorithm on the Lorenz system, double pendulum, and Langevin dynamics, respectively. The $x$ and $y$ axes display real and complex components, respectively, and we draw the unit circle in red for the reader's convenience.
}
   \label{fig:edmdeigvals}
\end{figure}

% \begin{figure}[!htb]
%     \centering
%     \includegraphics[width=0.2\linewidth]{figs/lorenz_edmdeigvals.png}
%     \includegraphics[width=0.2\linewidth]{figs/pendulum_edmdeigvals.png}

% \includegraphics[width=0.2\linewidth]{figs/langevin_edmdeigvals.png}
%     \caption{Eigenvalues of the lifted linear dynamics learned by the eDMD algorithm on the Lorenz system, double pendulum, and Langevin dynamics, respectively. The $x$ and $y$ axes display real and complex components, respectively, and we draw the unit circle in red for the reader's convenience.
%     }
% \label{fig:edmdeigvals}
% \end{figure}

In addition, this numerical experiment allows us to compute an estimate for the spectral gaps: if we sort $1 = |\lambda_1| > |\lambda_2| \geq \ldots$, the spectral gap is given by $\rho = 1 - |\lambda_2|$. We give the empirical spectral gaps for these systems below:
$$\rho_{\text{Lorenz}} = 0.00168,\quad \rho_{\text{pendulum}} = 0.00477, \quad \rho_{\text{Langevin}} = 0.000694.$$
The mixing time of a system scales as $\frac{1}{\rho}$, and our regret bound of Theorem \ref{theorem:main_nonlinear} implies that after $\frac{1}{\rho^2}$ steps we get vanishing prediction error. For the above systems, the mixing time is around 1000 steps and the learning horizon is on the order of hundreds of thousands or millions of steps; however, we see from our experiments that Algorithm \ref{alg:clsf_regressionmain}'s learning begins immediately. In other words, \textit{spectral filtering appears able to predict transient dynamics}. This indicates that the dependence of our analysis on a $\frac{1}{\sqrt{T}}$ spectral gap is not sharp, and leaves room for theoretical improvement.


% \evannote{manually discretize state space to find the Markov chain matrix for Lorenz + pendulum. look at its spectrum and compare to that of eDMD koopman eigenvalues. maybe implement a way to find $Q_\star$ numerically? make a triple plot -- Markov chain eigenvalues, Koopman eigenvalues, and the spectral constraint set we are using}

% \evannote{implement Langevin prediction, test a grid sampled uniformly vs one with Langevin}
\section{Conclusions and Future Work}

Deviating from the extensive literature on learning in dynamical systems, we have introduced a new approach based on improper learning and spectral filtering. We demonstrated how we can use classical tools from control theory, such as the Koopman operator and the Luenberger program for an observer system, to prove the {\it existence} of a real-diagonalizable linear dynamical system that is equivalent to the given nonlinear one in terms of observations. We then used the power of spectral filtering over this real-diagonalizable system to learn it provably and efficiently via convex optimization. 
The new methodology is accompanied by experiments on both linear and nonlinear systems validating the predicted efficiency, robustness, and accuracy, and showing that the regret scaling follows the control-theoretic condition number $Q_\star$.
 
Unsurprisingly, augmenting the observations nonlinearly improves the performance of an otherwise linear predictor. 
We suspect that much of the two-part eDMD paradigm (fixing/learning a nonlinear dictionary of observables and learning the lifted dynamics) will benefit from the use of spectral filtering -- we propose no changes to the lifting part, but only to learn to predict lifted dynamics via spectral filtering instead of direct LDS parametrization.
Deep neural networks with alternating MLP and SSM layers are currently a state of the art for deep learning of dynamical systems, and they may be viewed as jointly learning a nonlinear lifting and (linear) dynamics in the lifted space\footnote{This is also a current paradigm for robotics control and planning tasks, in which one learns an embedding from e.g. pixel space to latent space together with a "world model" that tracks these latent dynamics.}. 
The success of the \verb|SFeDMD| algorithm over \verb|eDMD| strongly motivates the use of STU \cite{liu2024flash} over direct SSM layers in such arrangements, since we see that spectral filtering is able to make better use of a nonlinear lifting for the task of sequence prediction, with improvement in terms of accuracy, number of parameters, optimization robustness, and no decrease in efficiency\footnote{During inference/rollout, one can use fast online convolution methods such as \cite{agarwal2024futurefillfastgenerationconvolutional} or LDS distillation methods such as \cite{shah2025spectraldsprovabledistillationlinear} to achieve the same generation efficiency as directly parameterizing an LDS.}. We leave large-scale experimentation on more complex tasks to future work.

On the theoretical side, the most natural next step is to remove the need for a $\frac{1}{\sqrt{T}}$ spectral gap in the observer system (since we don't see its effect much in practice). The work is already done in terms of spectral filtering (we can use the $\rho = 0$ case of Lemma \ref{lemma:spectralfiltering}), and the only thing to be done is to construct a high-dimensional LDS which approximates the original nonlinear system for which both $R_C$ and $R_{x_0}$ can be small. We anticipate this to be achievable using the more standard Koopman operator methodology. In addition, one could adapt our analysis to tolerate adversarial disturbances to the nonlinear system by translating them to additive disturbances to the original LDS and then applying our Theorem \ref{theorem:main_linear}(c) -- with a Koopman eigenlifting, this would require control on the Lipschitz constants of Koopman eigenfunctions, which can become technically involved. We leave these extensions to future work. A more significant extension would be to allow for the presence of open-loop inputs in the nonlinear system -- both the Koopman formalism and our discretized Markov chain are not well-adapted to handle non-control-affine systems, and we consider this a valuable direction for future work.
% \input{A8_Acknowledgments}

\bibliography{main.bib}
\bibliographystyle{plain}

\appendix
\section{Proof of Lemma \ref{lemma:spectralfiltering}}   \label{appendix:SFlemma}

In this section we provide a proof of Lemma \ref{lemma:spectralfiltering}, which follows the proof techniques from \cite{hazan2017learning}.
\begin{proof}
We measure regret against $\mathrm{LDS}(\rho, \gamma, \kappa, R_{B}, R_C, R_{x_0})$, which we recall to be the class of LDS predictors with spectrum $\sigma(A) \subset [0, 1-\rho] \cup \{1\} \cup \mathbb{D}_{1-\gamma}$, spectral condition number $\kappa_{\mathrm{diag}}(A) \leq \kappa$, and the bounds $\|B\|_F \leq R_{B}$, $\|C\|_F \leq R_C$, and $\|x_0\| \leq R_{x_0}$. For notation, we let $\Theta = (A, B, C, x_0)$ denote the LDS predictor and $y_t^\Theta$ its predictions, and $A = HDH^{-1}$ with $D = \operatorname{diag}(\alpha_1, \ldots, \alpha_{d_h})$ the eigenvalues of $A$.

We first prove an approximation result, i.e. that there exist parameters $J, M, P, N$ that match the predictions of any given LDS.
As in Section 2.2 of \cite{hazan2017learning}, our predictor $\hat{y}_{t+1}$ will add a prediction of the derivative of the impulse response $y_{t+1}^\Theta - y_{t}^\Theta$ to $y_{t}$. The derivative is
\begin{align*}
    y_t^\Theta - y_{t-1}^{\Theta}  & = CBu_{t-1} + \sum_{i=1}^{t} C(A^{i} - A^{i-1})Bu_{t-i-1} + C(A^t - A^{t-1})x_0 \\
    & = CBu_{t-1} + \sum_{i=1}^{t} CH(D^{i} - D^{i-1})H^{-1}Bu_{t-i-1} + C(A^t - A^{t-1})x_0\,.
\end{align*}
We will order the eigenvalues so that there exists $N \in [d_h]$ such that $\alpha_j \in [0, 1]$ for all $j \leq N$ and $|\alpha_j| \leq 1 - \gamma$ for all $j > N$. Defining $\mu_\alpha = (\alpha - 1) [1, \alpha, \ldots, \alpha^{T-2}] \in \R^{T-1}$, we have
\begin{align*}
    y_t^\Theta - y_{t-1}^{\Theta} 
    & = CBu_{t-1} + \sum_{i=1}^{t} \sum_{j=1}^{d_h} (\alpha_j^i - \alpha_j^{i-1}) CHe_je_j^\top H^{-1}Bu_{t-i-1} + C(A^t - A^{t-1})x_0 \\
    &= CBu_{t-1} + \sum_{j=1}^{d_h} CHe_je_j^\top H^{-1}B \left(\sum_{i=1}^t (\alpha_j - 1) \alpha_j^{i-1} u_{t-i-1}\right) + C(A^t - A^{t-1}) x_0 \\ &= CBu_{t-1} + \sum_{j=1}^{d_h} CHe_je_j^\top H^{-1} B\mu_{\alpha_j}^\top u_{t-2:t-T} + C(A^t - A^{t-1}) x_0,
\end{align*}
where $u_{t-2:t-T} \in \R^{(T-1) \times d_{\textrm{in}}}$ and so $\mu_{\alpha_j}^\top u_{t-2:t-T} \in \R^{d_{\textrm{in}}}$.
We split the sum over hidden dimension into two terms: $j \leq N$ will be handled via spectral filtering and $j > N$ via regression. Letting $\{\sigma_i, \phi_i\}_{i=1}^{T-1}$ be the eigenvalues and eigenvectors of the Hankel matrix from Algorithm \ref{alg:clsf_regressionmain} (see Section 3 of \cite{hazan2017learning}), we have
\begin{align*}
    \sum_{j=1}^{N} CHe_je_j^\top H^{-1} \mu_{\alpha_j}^\top u_{t-2:t-T} &=  \sum_{j=1}^{N} CHe_je_j^\top H^{-1} B\mu_{\alpha_j}^\top  \left(\sum_{i=1}^{T-1} \phi_i\phi_i^\top\right) u_{t-2:t-T}
    \\ &= \sum_{i=1}^{T-1}\sum_{j=1}^{N} CHe_je_j^\top H^{-1} B(\mu_{\alpha_j}^\top \phi_i) (\phi_i^\top u_{t-2:t-T})
\end{align*}
Define the optimal parameters $J_1 = CB$, $J_i = CH \left(\sum_{j=N+1}^{d_h} (\alpha_j^i - \alpha_j^{i-1}) e_je_j^\top H^{-1}B\right)$ for $i >1$, $M_i = \sigma_i^{-1/4} CH\left(\sum_{j=1}^N \mu_{\alpha_j}^\top \phi_i e_je_j^\top \right) H^{-1}B$ for $i \in [T-1]$, and $P_1 = 1$, $P_i = 0$ for $i >1$, and $N \equiv 0$. 
% \evannote{@Elad this is the first wrong statement due to the $y_{t-1} \neq y_{t-1}^\Theta$ mismatch and factor of 2 in regret... I will ignore it from now on, but i think it is fixed by assuming a Lipschitz signal $y_t$}
Then, we have the approximation error\footnote{Technically, we are considering the class of LDS predictors which add the derivative of the impulse response to the ground truth $y_{t-1}$. Otherwise, we would suffer an additional $y_{t-1} - y_{t-1}^\Theta$ approximation error.}
$$y^\Theta_t - \hat{y}_t(J, M, P, N) = \sum_{i=m+1}^t J_i u_{t-i} + \sum_{i=h+1}^{T-1} \sigma_i^{1/4}M_i \langle \phi_i, u_{t-2:t-T}\rangle + C(A^t - A^{t-1})x_0.$$
We bound
\begin{align*}\|J_i\| &\leq \left(\sup_{\alpha \in \mathbb{D}_{1-\gamma}} |\alpha^i - \alpha^{i-1}|\right) \cdot \sum_{j=N+1}^{d_h} \|CHe_j\| \cdot \|e_j^\top H^{-1}B\| \\&
\leq 2 (1-\gamma)^{i-1} \cdot \|CH\|_F \cdot \|H^{-1}B\|_F \\
&\leq 2(1-\gamma)^{i-1} R_B R_C \kappa,
\end{align*}
where in the second line we applied Cauchy-Schwartz and in the third we used $\|S_1S_2\|_F \leq \min(\|S_1\|\|S_2\|_F,\|S_1\|_F\|S_2\|)$. Similarly, 
\begin{align*}
    \|M_i\| &\leq \sigma_i^{-1/4} \cdot \left(\sup_{\alpha \in [0, 1]} |\mu_{\alpha}^\top \phi_i|\right) \cdot \sum_{j=1}^N \|CHe_j\| \cdot \|e_j^\top H^{-1}B\|\\
    &\leq 6^{1/4} R_BR_C\kappa,
\end{align*}
where the second line uses Lemma E.4 of \cite{hazan2017learning}. Since $\|u_{t-i}\| \leq R$ and $|\langle \phi_i, u_{t-2}:t-T\rangle| \leq R \sqrt{T}$ by Cauchy-Schwartz, it remains to bound $\norm{C(A^t - A^{t-1})x_0}$. If $A$ has all eigenvalues 1 then this term is zero, and so suppose that the largest eigenvalue of $A$ not equal to $1$ has magnitude $1-\rho$. Then,
\begin{align*}
    \norm{C(A^t - A^{t-1})x_0} \leq \sum_{j=1}^{d_h} |\alpha_j^t - \alpha_j^{t-1}| \cdot \|CHe_j\| \cdot \|e_j^\top H^{-1} x_0\|.
\end{align*}
% As before, for $j > N$ we know $|\alpha_j^t - \alpha_j^{t-1}| \leq 2\cdot (1-\gamma)^{t-1}$, and for $j \leq N$ we can use that $|\alpha^t - \alpha^{t-1}| \leq \frac{1}{t}$ for $\alpha \in [0, 1]$. 
For $j$ such that $\alpha_j = 1$ the summand is 0, and for all other $j$ we know $|\alpha_j^t - \alpha_j^{t-1}| \leq 2 \cdot (1-\rho)^{t-1}$.
Furthermore, again by Cauchy-Schwartz we have $\sum_{j=1}^{d_h} \norm{CHe_j} \cdot \norm{e_j^\top H^{-1}x_0} \leq \norm{CH}_F \cdot \norm{H^{-1}x_0} \leq R_C \kappa R_{x_0}$. Combining these ingredients,
\begin{align*}
    \norm{y_t^\Theta - \hat{y}_t(J, M, P, N)} &\leq 
    2R_B R_C\kappa R \cdot \sum_{i=m+1}^T (1-\gamma)^{i-1} + 6^{1/4}R_B R_C\kappa R \sqrt{T} \cdot \sum_{i=h+1}^{T-1} \sigma_i^{1/4} + 2R_C\kappa R_{x_0} \cdot (1-\rho)^{t-1} 
\end{align*}
We note that $(1-\gamma)^n \leq e^{-\gamma n}$ and that $\sigma_i^{1/4} \leq K e^{-i/(2\log T)}$ for an absolute constant $K < 10^6$ by Lemma E.2 of \cite{hazan2017learning}. Since $\sum_{i > n}e^{-ai} \leq \frac{e^{-an}}{a}$, we get the approximation result
\begin{align*}
    \norm{y_t^\Theta - \hat{y}_t(J, M, P, N)} &\lesssim R_BR_C \kappa R \cdot \frac{e^{-\gamma (m-1)}}{\gamma} + R_BR_C \kappa R \log(T) \sqrt{T} \cdot e^{-\frac{h}{2 \log T}} + R_C \kappa R_{x_0} \cdot e^{-\rho (t-1)}
\end{align*}
However, we also know that $ \norm{y_t^\Theta - \hat{y}_t(J, M, P, N)} \leq 2R$, and so if we let $\tau = \max\left(\frac{\log( R_C \kappa R_{x_0} / 2R)}{\rho}, 1\right)$ then
$R_C \kappa R_{x_0} e^{-\rho \tau} \leq 2R$, and therefore
\begin{align*}
    \norm{y_t^\Theta - \hat{y}_t(J, M, P, N)} &\lesssim R_BR_C \kappa R \cdot \frac{e^{-\gamma (m-1)}}{\gamma} + R_BR_C \kappa R \log(T) \sqrt{T} \cdot e^{-\frac{h}{2 \log T}} + R \cdot \begin{cases}
        1 & t \leq \tau \\ 
        e^{-\rho(t-\tau)} & t > \tau
    \end{cases}
\end{align*}
Summing over all $t$ and using $\sum_{t > \tau}e^{-\rho (t-\tau)} \leq \frac{1}{\rho}$ gives
\begin{align*}
    \sum_{t=1}^T\norm{y_t^\Theta - \hat{y}_t(J, M, P, N)} &\lesssim R_BR_C \kappa R T\cdot \frac{e^{-\gamma (m-1)}}{\gamma} + R_BR_C \kappa R \log(T) T^{3/2} \cdot e^{-\frac{h}{2 \log T}} \\ &\quad + \frac{ R \log_+(2R_C \kappa R_{x_0} / R)}{\rho}\;.
\end{align*}
Taking $m = \Theta\left(\frac{1}{\gamma} \log\left(\frac{T}{\gamma}\right)\right)$ and $h = \Theta\left(\log^2T\right)$, we get for any LDS parameterized by $\Theta$ there exist parameters $J_\Theta, M_\Theta, P_\Theta, N_\Theta$ for which
$$\sum_{t=1}^T \norm{y_t^\Theta - \hat{y}_t(J, M, P, N)} \leq O\left(R_B R_C \kappa R \log(T) \sqrt{T} +\frac{ R\log_+(R_C \kappa R_{x_0} / R)}{\rho}\right).$$
If we don't want to assume a spectral gap (i.e. $\rho = 0$), then we can instead use that $\sup_{\alpha \in [0, 1]} |\alpha^t - \alpha^{t-1}| \leq \frac{1}{t-1}$ and $\sup_{\alpha \in \mathbb{D}_{1-\gamma}} |\alpha^t - \alpha^{t-1}| \leq 2(1-\gamma)^{t-1}$ to get
$$\sum_{t=1}^T\norm{C(A^t - A^{t-1})x_0} \leq \|C\| \kappa \cdot \left(\frac{1}{\gamma} + \log T\right).$$
Now, we know that the loss is a convex function of $J, M, P, N$, which means that we can perform optimization via OGD. By our earlier bounds, the diameter of the constraint set is $D = O((m+h) R_B R_C \kappa)$. The gradient of the unsquared $\ell_2$ loss w.r.t. the prediction error is of size 1, and the gradient of the prediction w.r.t. $J$ and $P$ are bounded by $mR$ and $R$, respectively. The gradient w.r.t. $M_i$ is bounded by $\sigma_i^{1/4} \cdot |\langle \phi_j, u_{t-2:t-T}\rangle| \lesssim R \log T$, where we applied Corollary E.6 of \cite{hazan2017learning} and Holder's inequality. In total, we get the gradient bound $G \leq O(mR + hR \log T) \leq O((m+h)R\log T)$. By the OGD regret guarantee (Theorem 3.1 of \cite{hazan2016introduction}), if $\hat{y}_t$ are the iterates of Algorithm \ref{alg:clsf_regressionmain} then
\begin{align*}\sum_{t=1}^T \norm{\hat{y}_t - y_t} - \min_{\Theta} \sum_{t = 1}^T \norm{\hat{y}_t(J_\Theta, M_\Theta, P_\Theta, N_\Theta) - y_t} &\leq O\left((m+h)^2 R_B R_C \kappa R \log(T) \sqrt{T}\right).\end{align*}
\end{proof}
\begin{remark}
    The above proof shows that, if one increases $m$ and $h$ by $\log( R_B R_C \kappa R)$, then the approximation error can be made independent of those constants. By these proof techniques, the only place where one must incur $\Theta(\kappa)$ regret is in the diameter of the constraint set. 
    This dependence can be improved to $\Theta(\log \kappa)$ for the square loss with more sophisiticated online learning algorithms, such as the Azoury-Vovk-Warmuth forecaster \cite{azoury2001relative}, and in certain cases Online Newton Step \cite{hazan2016introduction}. This is implemented in Theorem \ref{theorem:logdependence}.
\end{remark}

\begin{remark}
\label{remark:spectralfilteringproof}
This proof has two main differences from that of Theorem 1 of \cite{hazan2017learning}. Firstly, we prove things for an algorithm which combines regression terms in order to learn systems with spectrum in $[0, 1] \cup \mathbb{D}_{1-\gamma}$; spectral filtering handles the $[0, 1]$ part and regression with $\widetilde{O}\left(\frac{1}{\gamma}\right)$ parameters handles the $\mathbb{D}_{1-\gamma}$ part. The second (and more theoretically interesting) difference is that rather than handling the vanishing initial state using the envelope $\sup_{\alpha \in [0, 1]} |\alpha^t - \alpha^{t-1}| \leq \frac{1}{t}$ to get linear decay, we make use of a spectral gap condition to get exponential decay of the initial state effect. 
% Actually, we require this in our application because $\|C\|_F$ is large, but as a byproduct we get 
\end{remark}


% \newpage
%     We measure regret against a class of LDS which predicts $\widehat{y}_t$ by adding the prediction of the derivative of the impulse response  $y_t^\Theta - y_{t-1}^\Theta$ to $y_{t-1}$ (see Section 2.2 of \cite{hazan2017learning}). For an LDS parameterized by $\Theta = (A, B, C, x_0)$  where $A = HDH^{-1}$, this derivative is:
%     \begin{align*}
%     y_t^\Theta - y_{t-1}^{\Theta}  & = CBu_{t-1} + \sum_{i=1}^{t-1} C(A^{i} - A^{i-1})Bu_{t-i-1} + C(A^t - A^{t-1})x_0 \\
%     & = CBu_{t-1} + \sum_{i=1}^{t-1} CH(D^{i} - D^{i-1})H^{-1}Bu_{t-i-1} + CH(D^t - D^{t-1})H^{-1}x_0\,.
%     \end{align*}
%     We shall assume that there exists $N \in [d_h]$ such that for all $i \leq N$, $\alpha_i \in [0, 1]$ and for all $i \geq N$, $|\alpha_i| \leq 1-\rho$.
%     Defining $\mu_\alpha = (\alpha - 1)[1, \alpha, \dots, \alpha^{T-2}] \in \reals^{T-1}$, we have
%     \begin{align*}
%         y_t^\Theta - y_{t-1}^{\Theta} 
%         % & = CBu_{t-1} + \sum_{i=1}^{t-1} CH\left(\sum_{j=1}^{d_h} (\alpha_j^i - \alpha_j^{i-1}) e_j e_j^\top\right)H^{-1}Bu_{t-i-1} + CH\left(\sum_{j=1}^{d_h} (\alpha_j^i - \alpha_j^{i-1}) e_j e_j^\top\right)H^{-1}x_0 \\
%         & = CBu_{t-1} + \sum_{j=1}^{N} CHe_je_j^\top H^{-1}B u_{t-2:t-T}\mu_{\alpha_j} + \sum_{j=1}^{N} (\alpha_j^t - \alpha_j^{t-1})CHe_je_j^\top H^{-1}x_0 \\
%         & + \sum_{i=1}^{t-1} CH\left(\sum_{j=N+1}^{d_h}(\alpha_i^i - \alpha_j^{i-1}) e_je_j^\top\right) H^{-1}Bu_{t-i-1} + \sum_{j=N+1}^{d_h} (\alpha_j^t - \alpha_j^{t-1})CHe_je_j^\top H^{-1}x_0
%     \end{align*}
%     Let $\phi_i$ be as defined in Algorithm \ref{alg:clsf_regressionmain}. Then,
%     \begin{align*}
%         \sum_{j=1}^{d_h} CHe_je_j^\top H^{-1}B u_{t-2:t-T}\mu_{\alpha_j} & = \sum_{j=1}^{d_h} CHe_je_j^\top H^{-1}B u_{t-2:t-T}\left(\sum_{i=1}^{T-1}\phi_i\phi_i^\top\right)\mu_{\alpha_j} \\
%         & = \sum_{i=1}^{T-1} CH\left(\sum_{j=1}^{d_h}\phi_i^\top\mu_{\alpha_j}e_je_j^\top\right) H^{-1}B u_{t-2:t-T}\phi_i \,.
%     \end{align*}
%     Setting $J_0 = CB$, setting $J_i = CH\left(\sum_{j=N+1}^{d_h}(\alpha_i^i - \alpha_j^{i-1}) e_je_j^\top\right) H^{-1}B$ for $i \in [m]$ and setting  $M_i = \sigma_i^{-1/4} CH\left(\sum_{j=1}^{d_h}\phi_i^\top\mu_{\alpha_j}e_je_j^\top\right) H^{-1}B$, we get that
%     \begin{align*}
%         \norm{y_t^\Theta - y_t(M, J)} & = \norm{\sum_{i=h+1}^{t-1} CH\left(\sum_{j=1}^{N}\phi_i^\top\mu_{\alpha_j}e_je_j^\top\right) H^{-1}B u_{t-2:t-T}\phi_i} + \norm{\sum_{j=1}^{N} (\alpha_j^t - \alpha_j^{t-1})CHe_je_j^\top H^{-1}x_0} \\
%         & + \norm{\sum_{i=m+1}^{t-1} CH\left(\sum_{j=N+1}^{d_h}(\alpha_j^i - \alpha_j^{i-1})e_je_j^\top\right)H^{-1}Bu_{t-i-1}} + \norm{\sum_{j=N+1}^{d_h} (\alpha_j^t - \alpha_j^{t-1}) CHe_je_j^\top H^{-1}x_0}
%         \end{align*}
%         Next, we shall bound each of these terms. In order to bound the first term, we use the Lemma E.4 of \cite{hazan2017learning}. For some absolute constants $K_0, K_1, c$,
%         \begin{align*}
%             \norm{\sum_{i=h+1}^{t-1} CH\left(\sum_{j=1}^{N}\phi_i^\top\mu_{\alpha_j}e_je_j^\top\right) H^{-1}B u_{t-2:t-T}\phi_i} & \leq \norm{C}_F\norm{H}\norm{H^{-1}}\norm{B}_F R \sqrt{T} \sum_{i=h+1}^{t-1} |\phi_i^\top \mu_{\alpha_j}| \\
%             & \leq \kappa(\Theta)R_\Theta^2R\sqrt{T}\sum_{i=h+1}^{t-1}K_0c^{-i/\log T} \\
%             & \leq K_1\kappa(\Theta)R_\Theta^2R\sqrt{T}\log T c^{-h/\log T} \,.
%         \end{align*}
%         The next term is bounded using the fact that $\max_{\alpha \in [0, 1]} |\alpha^t - \alpha^{t-1}| \leq \frac{1}{t}$,
%         \begin{align*}
%         \norm{\sum_{j=1}^{N} (\alpha_j^t - \alpha_j^{t-1})CHe_je_j^\top H^{-1}x_0} \leq \norm{C}_F\norm{H}\norm{H^{-1}}\norm{x_0}\max_{j \in [N]} |\alpha_j^t - \alpha_j^{t-1}| \leq  \frac{1}{t}\kappa(\Theta)R_\Theta^2
%     \end{align*}
%     The third is bounded by using the fact that $|\alpha_j| \leq 1-\rho$ for $j \geq N$.
%     \begin{align*}
%         \norm{\sum_{i=m+1}^{t-1} CH\left(\sum_{j=N+1}^{d_h}(\alpha_j^i - \alpha_j^{i-1})e_je_j^\top\right)H^{-1}Bu_{t-i-1}} & \leq \sum_{i=m+1}^{t-1} |\alpha_j^i - \alpha_j^{i-1}|\norm{C}_F\norm{H}\norm{H^{-1}}\norm{B}\norm{u_{t-i-1}} \\
%         & \leq 2\kappa(\Theta)R_\Theta^2R \sum_{i=m+1}^\infty (1-\rho)^{i-1} \leq \frac{2\kappa(\Theta)R_\Theta^2R}{\rho} (1-\rho)^m
%     \end{align*}
%     The fourth term is bounded similarly as follows:
%     \begin{align*}
%         \norm{\sum_{j=N+1}^{d_h} (\alpha_j^t - \alpha_j^{t-1}) CHe_je_j^\top H^{-1}x_0} & \leq \norm{C}_F\norm{H}\norm{H^{-1}}\norm{x_0}\max_{j \in [N+1, d_h]}|\alpha_j^t - \alpha_j^{t-1}| \\
%         & \leq  2\kappa(\Theta)R_\Theta^2 (1-\rho)^{t-1}
%     \end{align*}
%     Summing this up over all $t \in [T]$,
%     \begin{align*}
%         \sum_{t=1}^T \norm{y_t^\Theta - y_t(M, J)} & \leq K_1\kappa(\Theta) R_\Theta^2  R T^{3/2} \log T c^{-h/\log T} + \kappa(\Theta) R_\Theta^2 \log T \\
%         & \quad + \frac{2\kappa(\Theta)R_\Theta^2R T}{\rho} e^{-\rho m} + \frac{2\kappa(\Theta)R_\Theta^2}{\rho}\\
%         & \leq K_2\kappa(\Theta)R_\Theta^2R\sqrt{T}\log T + \frac{2\kappa(\Theta)R_\Theta^2}{\rho}
%     \end{align*}
%     for the choice of $h = C_1\log^2 T$ and $m = \frac{C_2}{\rho} \log\left(\frac{T}{\rho}\right)$ for some absolute constants $C_1$, $C_2$ and $K_2$. 
%     This proves the approximation result. Observe that $y_t(M, J)$ is linear in $M$ and $J$ and hence it is learnable by Online Gradient Descent. The diameter of the set is given by:
%     \begin{align*}
%         \norm{J_0} \leq R_\Theta^2, \quad \norm{J_i} \leq 2\kappa(\Theta)R_\Theta^2, \quad \norm{M_i}_F \leq R_\Theta^2 \norm{H}\norm{H^{-1}} \max_{l \in [d_h]} \sigma_l^{-1/4} \phi_l^\top \mu_{\alpha_j}
%     \end{align*}
%     Using lemma E.4 of \cite{hazan2017learning}, we observe that $\sigma_l^{-1/4} \phi_l^\top \mu_{\alpha_j}$ are upper bounded by an absolute constant. Thus, $\norm{M, J} \leq C_3\kappa(\Theta)R_\Theta^2  \sqrt{mh}$ for some absolute constant $C_3$. The gradient of the $\ell_2$ loss w.r.t $y_t(M, J)$ is $1$. Thus, the gradient w.r.t $M, J$ is bounded as follows:
%     \begin{align*}
%         \norm{\nabla_{M, J} y_t(M, J)} & \leq \norm{U_{t-1:t-m}} + \norm{\sum_{i=1}^h \sigma_i^{1/4}U_{t-2:t-T} \phi_i} \leq K_3R \log T \sqrt{mhd_{out}} 
%     \end{align*}
%     for some absolute constant $K_3$.
%     Using Theorem 3.1 of \cite{hazan2016introduction}, for $\eta = C_3\kappa(\Theta)R_\Theta^2 / R \sqrt{Td_{out}} \log T$ for some absolute constant $C_3$, we get that:
%     \begin{align*}
%         \sum_{t=1}^T \norm{\widehat{y}_t - y_t} - \min_{M, J \in B(C_3\kappa(\Theta)R_\Theta^2\sqrt{mh})}\sum_{t=1}^T\norm{{y_t(M, J)} - y_t} \leq K_4\kappa(\Theta)R_\Theta^2 R mh \sqrt{Td_{out}}\log T
%     \end{align*}
%     for some constant $K_4$. Using the approximation result, and substituting the value of $h$, we get that:
%     \begin{align*}
%         \sum_{t=1}^T \norm{\widehat{y}_t - y_t} - \min_{\Theta \in LDS(\Theta)}\sum_{t=1}^T\norm{{y_t(\Theta)} - y_t} \leq \frac{C_4\kappa(\Theta) R_\Theta^2 R}{\rho} \sqrt{Td_{out}}\log^3 T \log\left(\frac{T}{\rho}\right)
%     \end{align*}



\section{Lower Bounds} \label{sec:lower-bounds}

In this section we give lower bounds for learning sequences generated by linear and nonlinear dynamical systems, which complement our upper bounds and show them to be tight in certain respects. This does not mean they cannot be improved: to the contrary, this investigation motivates further study of the leading order constants and parameters of our bounds.

\subsection{LDS with Noise}

We have shown that spectral filtering with observation feedback is able to learn asymmetric marginally stable LDS's under adversarial disturbances. Our regret guarantee scales with two quantities: the optimal observer complexity $Q_\star$ and the disturbance sizes $ \sum_{t=0}^{T-1} \|w_t\|$. In this section, we investigate the the leading order terms of this result via lower bounds for sequence prediction in terms of these quantities. 
In particular, we prove that any algorithm which is not able to see $w_t$ when predicting $\hat{y}_{t+1}$ suffers a linear cost in terms of the disturbance sizes and ambient dimension up to constant factors:

\begin{theorem}
\label{theorem:lowerbound}
Let $\mA$ be any algorithm that predicts $\hat{y}_{t+1}$ using $u_1, \ldots, u_t$, $y_1, \ldots, y_t$, and $w_1, \ldots, w_{t-1}$. Then, there exists a problem instance $(A, B, C, x_0)$ and sequences $u_1, \ldots, u_T$ and $w_1, \ldots, w_T$ satisfying Assumption \ref{definition:linear_realizable} for which
$$\sum_{t = 1}^T \|\hat{y}_t^\mA - y_t\| \geq \Omega\left( d  \sum_{t=0}^{T-1} \| C w_t\| \right) $$
\end{theorem}
\begin{proof}
Consider the following linear dynamical system. We let $A$ be the permutation matrix over $d$ elements, $B$ is zero, and $C$ is the first standard basis vector, as given by
$$ A_{d}^{\mathrm{perm}} = 
 \begin{bmatrix}
 0 & 0 & \cdots & 0 & 1 \\
 1 & 0 & \cdots & 0 & 0 \\
 0 & 1 & \cdots & 0 & 0 \\
 \vdots & \vdots & \ddots & \vdots & \vdots \\
 0 & 0 & \cdots & 0 & 0 \\
 0 & 0 & \cdots & 1 & 0 \\
 \end{bmatrix} \ , \ B = 0 \ , \ 
 C =  
 \begin{bmatrix}
 1 \\
 0 \\
 \vdots \\
 0 \\
 \end{bmatrix} . 
$$
Notice that the observation at time $t$ is given by
$$ y_t = w_{1:t}(t \mod d) , $$
that is, the sum of all noises at a rotating coordinate according to the permutation.  Suppose $t \mod d = 1$ for simplicity. 

Consider any prediction algorithm $\mA$. We create a noise sequence as follows: in the $d$ iterations before the iteration $t$ such that $t \mod d = 1$, the noises $w_\tau(1)$ are going to be either $-y_{t-d}-1$ or $-y_{t-d} + 1$, chosen uniformly at random, and equal for these $d$ iterations. The following hold:
\begin{itemize}
\item Regardless of the algorithm $\mA$ prediction, it's loss is on expectation can predict w.l.o.g. $y_{t-d} \in \{ -d,d \} $, since the observation belongs to this set.  We have that $ \|\hat{y}_t^\mA - y_t\|  $ is zero w.p. $\frac{1}{2}$ and $2d$ w.p. $\frac{1}{2}$.

\item The expected loss of the algorithm is thus, over $T$ iterations,
$$ \sum_{t=1}^T \|\hat{y}_t^\mA - y_t\| \geq {d}  T = {d} \sum_t \|C w_t\| $$

\item The optimal predictor in hindsight that has access to $w_{t-1}$, which has the correct sign of the upcoming prediction, and thus has loss of zero.   

\end{itemize}


\end{proof}


\subsection{Computational Lower Bounds}

The preceding lower bound applies to linear dynamical systems with adversarial (or even stochastic) noise. We next give a computational lower bound sketch, based on cryptographic hardness assumptions,  for deterministic nonlinear dynamical systems.  

\begin{theorem}
% [Computational hardness for sequence prediction up to polynomial horizon]
\label{thm:computational-lb}
Let $\mathrm{PRG}:\{0,1\}^d \to \{0,1\}^T$ be a cryptographically secure pseudorandom generator with stretch $T=\mathrm{poly}(d)$. 
Consider the deterministic dynamical system
\[
\mathcal{X} = \{0,1\}^d \times \{0,1,\ldots,T-1\},\quad
F(s,t) = (s,(t+1) \bmod T),\quad
g(s,t) = 2\,\mathrm{PRG}(s)_t - 1 \ \in \ \{-1,+1\},
\]
and observations $y_t = g(x_t)$ from hidden initial state $(s,0)$.

Then, for every probabilistic polynomial-time predictor $\widehat{f}$, every polynomial $p(\cdot)$, and all sufficiently large $d$,
\[
\forall\, t\in\{0,\ldots,T-1\}:\qquad
\mathbb{E}\left[\,\big| \widehat{f}(y_0,\ldots,y_t) - y_{t+1} \big|\,\right] 
\ \ge\ 1 - \tfrac{1}{p(d)}.
\]
That is, the expected $\ell_1$ loss per step is $1 - o(1)$, matching the loss of a random $\pm 1$ guess, for all $t<T$.

Moreover, there exists an algorithm running in $O(2^d \cdot \mathrm{poly}(T))$ time that, given $(y_0,\ldots,y_{T-1})$, recovers the seed $s$ by exhaustive search and then predicts all future outputs exactly.
\end{theorem}

\begin{proof}[Proof sketch]
Suppose for contradiction that there exists a PPT predictor $\widehat{f}$, an index $t<T$, and a polynomial $q(\cdot)$ such that 
\[
\mathbb{E}\left[\,\big| \widehat{f}(y_0,\ldots,y_t) - y_{t+1} \big|\,\right] 
\ \le\ 1 - \tfrac{2}{q(d)}.
\]
Since $y_{t+1}\in\{-1,+1\}$, this means $\widehat{f}$ predicts the sign of $y_{t+1}$ with probability at least $\tfrac12 + \tfrac{1}{q(d)}$. 
We now construct a distinguisher $\mathcal{D}$ for PRG security: 
Given a string $z\in\{-1,+1\}^T$, feed $(z_0,\ldots,z_t)$ to $\widehat{f}$, output “PRG” if the sign of $\widehat{f}$ matches $z_{t+1}$, and “uniform” otherwise. 
If $z$ is $\mathrm{PRG}(s)$ for random $s$, this succeeds with probability $\tfrac12+\tfrac{1}{q(d)}$; if $z$ is uniform, the success probability is exactly $\tfrac12$. 
This yields a polynomial-time distinguisher with advantage $1/q(d)$, contradicting PRG security.

The exhaustive search recovery algorithm enumerates all $2^d$ seeds $s'$, computes $\mathrm{PRG}(s')$, and finds the unique one matching the observed $(y_0,\ldots,y_{T-1})$. This takes $O(2^d \cdot \mathrm{poly}(T))$ time and recovers the seed exactly, enabling perfect prediction thereafter.
\end{proof}


\section{Luenberger Program}
\label{appendix:luenbergerprogram}


In this section, we investigate the Luenberger program of Definition \ref{def:luenbergerprogram}, restated here for convenience: for given $A \in \R^{d_h \times d_h}$, $C \in \R^{d_{\textrm{obs}}\times d_h}$, and $\Sigma \subset \mathbb{D}_1$, we consider the optimization
\begin{equation}
\label{eqn:luenbergerprogram}
\min_{L \in \R^{d_h \times d_{\textrm{obs}}}} \quad \kappa_{\textrm{diag}}(A-LC) \quad \text{s.t.} \quad A-LC \text{ has all eigenvalues in $\Sigma$},\end{equation}
    and we denote the minimal value by $Q_\star = Q(A, C, \Sigma)$ and a minimizing matrix, when it exists, by $L_\star$. First, we summarize some known formulas for the important quantities:
\begin{lemma}
\label{lemma:ackermann}
    Suppose that $(A, C)$ is observable. Then, for any monic, degree-$d_h$ polynomial $p$ with distinct roots, there is a unique $L \in \R^{d_h \times d_{\textrm{obs}}}$ such that $A - LC$ is diagonalizable with characteristic polynomial $p$. In particular, if $(A, C)$ is observable then $Q_\star(A, C, \Sigma) < \infty$ for all $\Sigma$ containing more than $d_h$ points which is symmetric about the real line. 
    % Furthermore, if $d_{\textrm{obs}} = 1$ then we have that
    % $$A - LC = H\operatorname{diag}(\lambda_1, \ldots, \lambda_{d_h}) H^{-1}$$
    % for $\lambda_j$ the roots of $p$, where
    % $$H = \left[(A-\lambda_1 \operatorname{Id})^{-1}C, \ldots, (A-\lambda_{d_h} \operatorname{Id})^{-1}C\right].$$
    % In this case we also have
    % Ackermann's formula
    % $$L = p(A) \cdot \mathcal{O}^{-1} \cdot e_n = H \cdot [1,\ldots, 1].$$
\end{lemma}
\begin{proof}
This is a standard result in eigenvalue placement, and for $d_{\text{obs}}=1$ there is an explicit formula known as Ackermann's formula (see e.g. Corollary 2.6 of \cite{Mehrmann1996}).
    % The formulas for $H$ and $L$ in terms of $H$ come from Corollary 2.5 of \href{https://etna.ricam.oeaw.ac.at/vol.4.1996/pp89-105.dir/pp89-105.pdf}{this paper}. Ackermann's formula is stated as Corollary 2.6.
\end{proof}
The above shows that observability implies $Q_\star < \infty$. Also, the bijection between $p$ and $L$ informs us that we may instead consider the optimization (\ref{eqn:luenbergerprogram}) as being parameterized by the desired poles. 
% Furthermore, we expect $Q_\star$ to depend on the condition number of $\mathcal{O}$ and the distances between the eigenvalues of $A$ and the optimal choice of poles.
As a crude but general upper bound:

\begin{lemma}
\label{lemma:ackermannbounds}
Suppose that $(A, C)$ is observable and $\dim(\Sigma) \geq 1$, let $\operatorname{vol}(\Sigma)$ denote the appropriate notion of volume, and let $N$ be the number of eigenvalues of $A$ not in $\Sigma$. Then, we have the bounds
$$Q_\star(A, C, \Sigma) \leq \left(1+\frac{2N^8}{\operatorname{vol}(\Sigma)}\right)^{2N-2} = \left(\frac{N}{\operatorname{vol}(\Sigma)}\right)^{O(N)},$$
$$\|L_\star\| \leq \frac{2Q_\star}{\sigma_{\textrm{min}}(C)}.$$
\end{lemma}
\begin{proof}
Let $\hat{A} = A - LC$ for notation. Firstly, since eigenvalue placement respects blocks of $A$ we can w.l.o.g. assume that $A$ is $N \times N$.
By Theorem 1 of 
% \href{https://www.sciencedirect.com/science/article/pii/S0024379597800298}{this paper}
\cite{conditionnumbereigengap}, we have
$$\kappa_{\textrm{diag}}(\hat{A}) \leq \left(1 + \frac{\alpha}{\delta}\right)^{2N-2}$$
for $\delta$ the minimum eigengap of $\hat{A}$ and $\alpha \leq N^{3/4} \cdot \mu_2(\hat{A})$, where we use the $\mu_j$ notation of 
% \href{https://www.sciencedirect.com/science/article/pii/0024379587902539}{this paper}
\cite{nonnormality}. By (C0), (C3), and (C5) of \cite{nonnormality} we know $\mu_2^2 \leq 4\sqrt{N} \|\hat{A}\|^2$ and therefore $\alpha \leq 2N\|\hat{A}\|$. If we take the rows of $H$ to be unit norm (i.e. normalizing the eigenvectors) then by Theorem 2 of 
% \href{https://www.emis.de/journals/GM/vol11nr1_2/gheorghiu/gheorghiu.pdf}{this paper} 
\cite{Gheorghiu2003}
we have that $$\|\hat{A}\| \leq 1 + N^{9/4} \rho(A) \gamma(A) \leq 1 + N^{23/4} \leq N^6,$$
where we also used $\ell_1,\ell_2$ norm inequalities.
Combining everything (and observing that we can choose to place the eigenvalues with an eigengap of $\frac{\operatorname{vol}(\Sigma)}{N}$) gives the bound on $\kappa_{\textrm{diag}}(\hat{A})$. Since it is independent of $p$, this implies the bound on $Q_\star$.
The bound on $\|L\|$ in terms of the condition number of $H$ is equation (22) of \cite{gainsizebound}
% \href{https://sci-hub.ru/10.1080/0020718508961188}{this paper} 
and the fact that $Q_\star \geq 1$.
\end{proof}
The above bound is very general, but is exponential in the hidden dimension. 
Next, we look into a representative example of an asymmetric system to demonstrate the finer properties of this optimization:



\begin{example}[Permutation system]
\label{example:permutationqstar}
    Let $A$ be a cyclical $n \times n$ permutation matrix, i.e.
\[
A = \begin{bmatrix}
0 & 1 & 0 & \cdots & 0 \\
0 & 0 & 1 & \cdots & 0 \\
\vdots & \vdots & \ddots & \ddots & \vdots \\
0 & 0 & \cdots & 0 & 1 \\
1 & 0 & \cdots & 0 & 0
\end{bmatrix},
\]
$C = e_1$ be the projection on the first coordinate, and $p$ be a monic degree-$n$ polynomial (with coefficients $p_j$ and roots $\lambda_j$).
Then, the gain $L \in \R^{n \times 1}$ and eigenvector matrix $H$ 
satisfy $$\|L\|^2 = \sum_{j=0}^{n} |p_j|^2, \quad \kappa(H) = \kappa\left(\operatorname{diag}\left(\frac{1}{1 - \vec{\lambda}^n}\right) \cdot V(\vec{\lambda})\right)$$
for $V(\vec{\lambda})$ the Vandermonde matrix with knots $\lambda_j$. In particular, $\|L\|$ can be made to be small but $\kappa(H)$ is always exponential if $\Sigma$ does not include the roots of unity, i.e. $$Q_\star = \Omega(2^n).$$ 
\end{example}
\begin{proof}
We will use the explicit formulas for $H$ and $L$ from Corollaries 2.5 and 2.6 of \cite{Mehrmann1996}.
    We know that $c^\top A^k e_j = \langle c, A^k e_j\rangle = \langle c, e_{j+k \mod n}\rangle = \delta_{j+k\mod n = 1}$, and so 
    $$CA = e_{2}, \quad \ldots, \quad CA^{n-1} =e_{n}$$
    This means that the observability matrix is 
    $$\mathcal{O} = \begin{bmatrix}
        C \\ CA \\ \vdots \\ CA^{n-1}
    \end{bmatrix} = \begin{bmatrix}
        e_1 \\ e_2 \\ \vdots \\ e_n
    \end{bmatrix} = \mathrm{Id}$$
In particular, since we know that $A$ is unitarily diagonalized as $A = FDF^{-1}$ for $F$ the DFT matrix and $D = \operatorname{diag}(1, \omega, \omega^2, \ldots, \omega^{n-1})$ with $\omega = e^{2\pi i / n}$, we find 
$$L = p(A) \cdot e_n = F \operatorname{diag}(1, p(\omega), p(\omega^2), \ldots, p(\omega^{n-1})) F^{-1} \cdot e_n$$
% This tells us that we may write $e_n$ in the DFT basis as
% $$e_n = \sum_{k=0}^{n-1} \frac{1}{\sqrt{n}} \omega^{(n-1)k} v^{(k)} \implies p(A) \cdot e_n = \sum_{k \leq n} \frac{\omega^{(n-1)k} p(\omega^{k-1})}{\sqrt{n}} v^{(k)}$$
% \evannote{the above line is wrong, but everything below is right and was checked in colab}
Writing this in the standard basis, we find that the $j^{th}$ coordinate of $L$ for $1 \leq j \leq n$ is
$$\sum_{k =0}^{n-1} \frac{\omega^{jk} \cdot p(\omega^k)}{n} $$
and so
\begin{align*}\|L\|^2 &= \frac{1}{n^2}\sum_{j = 1}^n \sum_{k_1 = 0}^{n-1} \sum_{k_2 = 0}^{n-1} {\omega^{j(k_1 + k_2)} \cdot p(\omega^{k_1}) \cdot p(\omega^{k_2})} 
\\ &= \frac{1}{n^2} \sum_{k_1 = 0}^{n-1} \sum_{k_2 = 0}^{n-1} {p(\omega^{k_1}) \cdot p(\omega^{k_2})} \cdot \sum_{j = 1}^n \omega^{j(k_1 + k_2)}
\end{align*}
We note that $\sum_{j=1}^n \omega^{q \cdot j}$ is equal to 0 unless $q$ is a multiple of $n$, in which case it equals $n$. Therefore, we are left with
\begin{align*}\|L\|^2 &= \frac{1}{n} \sum_{k_1 + k_2 = n} p(\omega^{k_1}) \cdot p({\omega^{k_2}}) \\&= \frac{1}{n}\sum_{k=0}^{n-1} p(\omega^k) \cdot p(\omega^{n-k})
\\&= \frac{1}{n} \sum_{k=0}^{n-1}|p(\omega^k)|^2,
\end{align*}
where we used that $p(\omega^{-k}) = \overline{p(\omega^k)}$ for a polynomial $p$ with real coefficients. Lastly, noting that
$$p(\omega^k) = \sum_{j=0}^n p_j \omega^{kj}$$
% \href{https://en.wikipedia.org/wiki/Discrete_Fourier_transform#Definition}
{we see that} the sequence $p(\omega^k)$ is none other than the DFT of the sequence $p_j$ of coefficients of $p$. By 
% \href{https://en.wikipedia.org/wiki/Discrete_Fourier_transform#The_Plancherel_theorem_and_Parseval's_theorem}
{Parseval's identity}, we get the stated result about $L$.

For $H$, we once again write $A = FDF^{-1}$, and so that $j^{th}$ eigenvector $h_j$ of $A - LC$ equals
$$h_j = (A - \lambda_j \operatorname{Id})^{-1}C = F \operatorname{diag}\left(\frac{1}{1 - \lambda_j}, \frac{1}{\omega - \lambda_j}, \ldots, \frac{1}{\omega^{n-1} - \lambda_j}\right)F^{-1}e_1$$
% By unitarity of $F$ and the fact that $F^{-1}e_1 = \frac{1}{\sqrt{n}}[1, \ldots, 1]$, this means that 
% $$\langle u_i, u_j\rangle = \frac{1}{n}\sum_{k=0}^{n-1} \frac{1}{\overline{(\omega^k - \lambda_i)}(\omega^k - \lambda_j)}$$
Noting that $F^{-1}e_1 = \frac{1}{\sqrt{n}}[1, \ldots, 1]$, we see that if $V'$ is the Cauchy matrix with entries
$$V_{ij}' = \frac{1}{\omega^{j-1} - \lambda_i} \quad (1 \leq i,j\leq n)$$
then $H = \frac{1}{\sqrt{n}}FV$. Therefore, the condition number of $H$ equals that of $V'$. 
By equation (7) of 
% \href{https://academicworks.cuny.edu/cgi/viewcontent.cgi?article=1385&context=gc_cs_tr}{this paper}
\cite{cauchyvandermondematrices}, we find that $V'$ has the same condition number as $(D^n - \operatorname{Id})^{-1}V$, where $V$ is the Vandermonde matrix with knots $\lambda_j$ and $D$ is the diagonal matrix with entries $\lambda_j$. We have arrived at the fact that
$$Q_\star = \min_{\vec{\lambda} \subset \Sigma} \kappa\left(\operatorname{diag}\left(\frac{1}{1 - \vec{\lambda}^n}\right) \cdot V(\vec{\lambda})\right).$$
The only Vandermonde matrices with subexponential condition number are those with knots close to the roots of unity \cite{pan2015badvandermondematrices}, in which case the row rescalings blow up.
As such, it can be seen that for any admissible choice of $\lambda_j$'s, the condition number of this row-rescaled Vandermonde matrix will be exponential in the dimension.
\end{proof}



% The question of conditioning of the eigenvalue placement problem was already studied in \href{https://sci-hub.ru/10.1137/s0895479896302382}{this paper}. 
% Combining some of their bounds yields:

% \begin{lemma}[Theorems 2.1 and 2.2 of \href{https://sci-hub.ru/10.1137/s0895479896302382}{this paper}]
% Suppose that $A$ is diagonalizable with eigenvalues $\lambda_j$, $d_\text{obs} = 1$, and $(A, C)$ is observable, and let $p$ be a monic degree-$d_h$ polynomial with distinct roots $\gamma_j$ such that $\lambda_i \neq \gamma_j$ for all $i, j$. Then, \evannote{something that bounds $\kappa$ by $\|L\|$ or otherwise by $p(A)$ and $\mathcal{O}^{-1}$}
% \end{lemma}
% \begin{proof}
% We assume w.l.o.g. that $A$ is diagonal. 
%     Let $e = [1, \ldots, 1] \in \R^{d_h}$ be the all 1's vector. Denoting by $L_c$ and $L_e$ the eigenvalue placement gains such that $A - L_cc^\top$ and $A - L_ee^\top$ have characteristic polynomial $p$, and let $\kappa_c^F = \kappa_{\textrm{diag}}(A - L_cc^\top)$ and $\kappa_e^F = \kappa_{\textrm{diag}}(A - L_ee^\top)$ be the diagonalizing condition numbers w.r.t. the Frobenius norm. By Theorem 2.2 of \evannote{cite the paper}, we know
%     $$\frac{\|L_e\|}{\|c\|} \leq \|L_c\| \leq \norm{\frac{1}{c}} \cdot \|L_e\| \quad \text{and} \quad \frac{\kappa_e^F}{\sqrt{d_h} \cdot \|c\| \cdot \norm{\frac{1}{c}}} \leq \kappa_c^F \leq \sqrt{d_h} \cdot \|c\| \cdot \norm{\frac{1}{c}} \cdot \kappa_e^F,$$
%     where $\frac{1}{c}$ denotes the vector with entries $\frac{1}{c_j}$ (note that since $A$ is diagonal and $(A, C)$ is observable, each coordinate of $c$ must be nonzero). By Theorem 2.1 of \evannote{cite the paper}, 
%     $$\|L_e\|^2 = \sum_{i=1}^{d_h} \frac{\prod_{k=1}^{d_h} |\lambda_i - \gamma_k|^2}{\prod_{k=1,k\neq i}^{d_h} |\lambda_i - \lambda_k|^2}$$
%     and Theorem 2.3 says that
%     $$\|$$
% \end{proof}
% \elad{looks like we can place all eigenvalues, equally spaced, in the interval [1/4,1/2], and that will give something like $d^d \sim 2^{d \log d}$, which is probably the best we can do in general. For special cases we can do much better perhaps, but not sure it is worth very detailed study at this point.}


\section{Other Guarantees}
\label{appendix:otherguarantees}
The presented framework is quite general: using the Markov chain discretization together with observer systems, we can construct comparator LDS's with any specified spectral structure (with a cost measured by $Q_\star$), and we combine this with a regret guarantee against the chosen class of comparator LDS's. 

Before deriving other forms of loss and regret bounds, we stop to note that since we prove things for the unsquared $\ell_2$ loss, everything extends immediately to all Lipschitz loss sequences (since we use OCO, the losses may be time-varying). Furthermore, since we assume boundedness of the outputs, we can also use the square loss with an $O(R)$ Lipschitz constant. 

To start, our proof technique yields for free the agnostic versions of our main results. 
% \evannote{do we require a lipschitz assumption on the signal?} 
It is nice to view our results in this way, as it defines exactly which types of computation are able to be efficiently implemented via spectral filtering -- on arbitrary data, if there exists a computer of the classes defined below which fits the data well with small $Q_\star$, then spectral filtering will succeed.

\begin{theorem}[Nonlinear, agnostic]
    Let $y_1, \ldots, y_T$ be an arbitrary sequence bounded by $R$. Fix the number of algorithm parameters $h = \Theta(\log^2 T)$ and $m \geq 1$, and for notation define $\gamma = \frac{\log(mT)}{m}$.

    Let $\Pi_\text{NL}\left( q\right)$ denote the class of nonlinear dynamical systems $(f, h, x_0)$ such that $\|x_t\| \leq R$ and $Q_\star\left(f, h, \frac{1}{T^{3/2}}, \left[0, 1 - \frac{1}{\sqrt{T}}\right] \cup \{1\} \cup \mathbb{D}_{1-\gamma}\right) \leq q$. Then, running Algorithm \ref{alg:clsf_regressionmain} with $J^t, M^t \equiv 0$, $\mathcal{K} = R_D$ for $D = \Theta((m+h)q^2)$, and  $\eta_t = \Theta\left(\frac{q^2}{R\sqrt{t}\log T}\right)$ produces a sequence of predictions $\hat{y}_1, \ldots, \hat{y}_T$ for which
    $$\sum_{t=1}^T \|\hat{y}_t - y_t\| - \min_{\pi \in \Pi_{\text{NL}}(q)} \sum_{t=1}^T \|\hat{y}_t^\pi - y_t\| \leq O\left(q^2 \log(q) \cdot d_\mathcal{X} R \log(RT) m^2 \log^7(mT)  \cdot \sqrt{T}\right)  . $$
\end{theorem}
\begin{proof}
    The proof proceeds in the same way as that of Theorem \ref{theorem:main_nonlinear}: we use Lemma \ref{lemma:discretization} to construct a high-dimensional LDS, which we have a regret guarantee against via Theorem \ref{theorem:agnosticlinear}.
\end{proof}

\begin{theorem}[Linear, agnostic]
    \label{theorem:agnosticlinear}
    Let $u_1, \ldots, u_T$ and $y_1, \ldots, y_T$ be arbitrary sequences bounded by $R$.
    Fix the number of algorithm parameters $h = \Theta(\log^2T)$ and $m \geq 1$, and for notation define $\gamma = \frac{\log(mT)}{m}$
    and $Q = R_BR_Cq + q^2$. Set $\mathcal{K} = R_D$ for $D = \Theta\left((m+h) Q\right)$ and $\eta_t = \Theta\left(\frac{Q}{R \sqrt{t}\log T }\right)$, and let $\hat{y}_1, \ldots, \hat{y}_T$ be the sequence of predictions that Algorithm \ref{alg:clsf_regressionmain} produces. 

Fix the constants $R_B, R_C, R_{x_0}, W$, and 
    let $\Pi_{\text{L}}\left(w, q, \Sigma\right)$ denote the class of LDS predictors (with arbitrary disturbances bounded by $\sum_{t=1}^{T-1}\|w_t\| \leq w$) which satisfy Assumption \ref{definition:linear_realizable} with the specified constants and for which $Q_\star(A, C, \Sigma) \leq q$. Then, we have the following three regret guarantees:

\begin{enumerate}
    \item[(a).] With $w = 0$,
    \begin{align*}
        \sum_{t=1}^T \norm{\hat{y}_t - y_t} - \min_{\pi \in \Pi_{\mathrm{L}}\left(0, q, [0, 1] \cup \mathbb{D}_{1-\gamma}\right)} \sum_{t=1}^T \|\hat{y}_t^\pi - y_t\| \leq O\left(Q R m^2 \log^7(mT) \cdot \sqrt{T} + R_C q R_{x_0 } \cdot (m + \log T)\right).
    \end{align*}
    \item[(b).] With $w = 0$,
    \begin{align*}
        \sum_{t=1}^T \norm{\hat{y}_t - y_t} - \min_{\pi \in \Pi_{\mathrm{L}}\left(0, q, \left[0, 1 - \frac{1}{\sqrt{T}}\right] \cup \{1\} \cup \mathbb{D}_{1-\gamma}\right)} \sum_{t=1}^T \|\hat{y}_t^\pi - y_t\| \leq O\left(Q R m^2 \log_+(R_Cq R_{x_0} / R)\log^7(mT) \cdot \sqrt{T}\right).
    \end{align*}
    \item[(c).] For any $\rho \in (0, \gamma]$, with $R' = R + \frac{W}{\rho}$,
    \begin{align*}
        \sum_{t=1}^T \norm{\hat{y}_t - y_t} - \min_{\pi \in \Pi_{\mathrm{L}}\left(w, q, \left[0, 1 - \rho\right]\cup \mathbb{D}_{1-\gamma}\right)} \sum_{t=1}^T \|\hat{y}_t^\pi - y_t\|  \leq \; & O\left(Q \cdot R' m^2 \log^7(mT) \cdot \sqrt{T} \right) \\ + &O\left(\frac{R' \log_+(R_Cq R_{x_0} / R') + R_C q w}{\rho}\right).
    \end{align*}
\end{enumerate}
\end{theorem}
\begin{proof}
    By the observer system construction, the approximation result that constitutes the first part of the proof of Lemma \ref{lemma:spectralfiltering} extends immediately to this agnostic case. Then, since OGD's guarantee is given in the regret formulation of online convex optimization, we get the stated results.
\end{proof}

The above statements describe the performance of Algorithm \ref{alg:clsf_regressionmain}, which is a combination of spectral filtering and regression. However, our techniques yield a sharp analysis for other types of algorithms. For simplicity, we consider two additional algorithms: (1) regression by itself and (2) spectral filtering combined with the Chebyshev sequence preconditioning of \cite{marsden2025universalsequencepreconditioning}. We will state the corresponding results for the nonlinear case with minimal proofs -- the reader can fill in the details, as well as any statements for the linear case.


\begin{algorithm}[htb]
\caption{Regression}
\label{alg:appendixregression}
\begin{algorithmic}[1]
\STATE \textbf{Input:} Horizon \(T\), number of autoregressive components $m$, step sizes $\eta_t$, convex constraint set \(\K\), prediction bound $R$.
\STATE Initialize $P_j^0\in \R^{d_{out} \times d_{out}}$ for $j \in \{1, \ldots, m\}$. 
\FOR{$t = 0, \ldots, T-1$}
\STATE Compute \begin{align*}\hat{y}_t(\Theta^t) &=  \sum_{j=1}^{m}P_j^ty_{t-j} , \end{align*}
and if necessary project $\hat{y}_t(P^t)$ to the ball of radius $R$ in $\mathcal{Y}$.
\STATE Compute loss $\ell_t(P^t) = \norm{\hat{y}_t(P^t) - y_t}$.
\STATE Update $P^{t+1} \leftarrow \Pi_\K\left[P^t - \eta_t  \nabla_P \ell_t(P^t)\right]$. 
\hfill \textcolor[HTML]{408080}{\# can be replaced with other optimizer}
\ENDFOR
\end{algorithmic}
\end{algorithm}

\begin{theorem}
    Let $y_1, \ldots, y_T$ be an arbitrary sequence bounded by $R$. Let $\Pi_{\mathrm{NL}}(q, \gamma)$ be the class of nonlinear dynamical systems $(f, h, x_0)$ such that $\|x_t\| \leq R$ and $Q_\star(f, h, \frac{1}{T^{3/2}}, \mathbb{D}_{1-\gamma}) \leq q$. Then, running Algorithm \ref{alg:appendixregression} with $m \geq \frac{1}{\gamma} \log\left(\frac{T}{\gamma}\right)$, $\mathcal{K} = R_D$ for $D = \Theta(mq^2)$, and $\eta_t = \Theta\left(\frac{q^2}{R\sqrt{t}\log T}\right)$ produces a sequence of predictions $\hat{y}_1, \ldots, \hat{y}_T$ for which
    $$\sum_{t=1}^T \|\hat{y}_t - y_t\| - \min_{\pi \in \Pi_{\text{NL}}(q, \gamma)} \sum_{t=1}^T \|\hat{y}_t^\pi - y_t\| \leq O\left(q^2 \log(q) \cdot d_\mathcal{X} R \log(RT) m^2 \log^7(mT)  \cdot \sqrt{T}\right)  . $$
\end{theorem}
\begin{proof}[Proof sketch]
    As before, we use Lemma \ref{lemma:discretization} to perform a global linearization via discretization. We then use eigenvalue placement to place all the eigenvalues of an observer system into $[0, 1-\gamma]$, and the dynamics matrix of this observer system will have spectral condition number $q$. The regression part of the bound from Lemma \ref{lemma:spectralfiltering} together with OGD's regret guarantee lets us conclude.
\end{proof}
\begin{remark}
\label{remark:regression}
    Consider a system with $N$ marginally stable modes, such as the permutation system of size $N$ from Example \ref{example:permutationqstar}. For such a system, $Q_\star$ is exponential in $N$ when we take the spectral constraint set $\Sigma = \mathbb{D}_{1-\gamma}$. The regret bound tells us that we require $T$ to also be exponential in $N$, which means we must have $m$ linear in $N$. This recovers what we would expect from the Cayley-Hamilton theorem, which implies that any LDS with hidden dimension $d_h$ should be learnable with $d_h$ autoregressive terms. This also exemplifies a main benefit that spectral filtering offers: it allows for $Q_\star$ to be much smaller when there are real and marginally stable eigenvalues, which in turn allows for learnability using fewer parameters.
\end{remark}



\begin{algorithm}[htb]
\caption{Observation Spectral Filtering + Chebyshev}
\label{alg:clsf_chebyshev}
\begin{algorithmic}[1]
\STATE \textbf{Input:} Horizon \(T\), number of filters \(h\), step sizes $\eta_t$, convex constraint set \(\K\), prediction bound $R$.
\STATE Compute $\{(\sigma_j, {\phi}_j)\}_{j=1}^k$ the top eigenpairs of the $(T-1)$-dimensional Hankel matrix of \cite{hazan2017learning}.

\STATE Initialize $N_i^0\in \R^{d_{out} \times d_{out}}$ for $i \in \{1, \ldots, h\}$. 
\FOR{$t = 0, \ldots, T-1$}
\STATE Compute $$\hat{y}_t(N^t) = -\sum_{j=1}^{\log T}c_jy_{t-j} +  \sum_{i=1}^h \sigma_i^{1/4}N_i^t \langle \phi_{j} ,  y_{t-1:t-T} \rangle, $$
and if necessary project $\hat{y}_t(N^t)$ to the ball of radius $R$ in $\mathcal{Y}$.
\STATE Compute loss $\ell_t(N^t) = \norm{\hat{y}_t(N^t) - y_t}$.
\STATE Update $N^{t+1} \leftarrow \Pi_\K\left[N^t - \eta_t  \nabla_N \ell_t(N^t)\right]$. 
\hfill \textcolor[HTML]{408080}{\# can be replaced with other optimizer}
\ENDFOR
\end{algorithmic}
\end{algorithm}

    
\begin{theorem}
    Let $y_1, \ldots, y_T$ be an arbitrary sequence bounded by $R$. Let $\Pi_{\mathrm{NL}}(q)$ be the class of nonlinear dynamical systems $(f, h, x_0)$ such that $\|x_t\| \leq R$ and $Q_\star\left( f, h, \frac{1}{T^{3/2}}, \Sigma\right) \leq q$ for $$\Sigma := \left\{\lambda \in \mathbb{D}_1 : \quad |\arg(\lambda)| \leq \frac{1}{64 \log^2(T)} \right\}.$$
    Then, running Algorithm \ref{alg:clsf_chebyshev} with $h = \Theta(\log^2(T))$, $\mathcal{K} = R_D$ for $D = \Theta(q^2)$, and $\eta_t = \Theta\left(\frac{q^2}{R\sqrt{t}\log T}\right)$ produces a sequence of predictions $\hat{y}_1, \ldots, \hat{y}_T$ for which
    $$\sum_{t=1}^T \|\hat{y}_t - y_t\| - \min_{\pi \in \Pi_{\text{NL}}(q)} \sum_{t=1}^T \|\hat{y}_t^\pi - y_t\| \leq \widetilde{O}\left(q^2 \cdot d_\mathcal{X} R \cdot T^{10/13}\right) . $$
\end{theorem}
\begin{proof}[Proof sketch]
    Form the observer system with a choice of $L$ such that all eigenvalues of $A - LC$ have complex component less than $\delta = \frac{1}{64 \log^2(T)}$. Apply Theorem C.2 of \cite{marsden2025universalsequencepreconditioning}, where the same trick as in Lemma \ref{lemma:spectralfiltering} may be used to improve the dependence on $R_C$.
\end{proof}

Lastly, we recall from the proof of Lemma \ref{lemma:spectralfiltering} that the approximation error can be made independent of $Q_\star$ with a cost of $\log Q_\star$ extra parameters, though the diameter of the constraint set will still depenend linearly on $Q_\star$. More sophisticated online algorithms can have logarithmic dependence on the diameter, yielding the following result. As before, we state this only for the nonlinear realizable case, and the linear case (as well as the agnostic cases) are proved along the same lines.

\begin{theorem}
    \label{theorem:logdependence}
    Let $y_1, \ldots, y_T$ be the observations of a nonlinear dynamical system satisfying Assumption \ref{assumption:nonlinear}, and let $\gamma \in (0, 1)$. 
    Let $Q_\star = Q_\star\left(f,h,\frac{1}{T^{3/2}}, \left[0, 1 - \frac{1}{\sqrt{T}}\right] \cup \{1\} \cup \mathbb{D}_{1-\gamma}\right)$ as specified in Definition \ref{def:nonlinearluenbergerprogram}. Fix the number of algorithm parameters $h = \Theta\left(\log^2(Q_\star T)\right)$ and $m = \Theta\left(\frac{1}{\gamma} \log\left(\frac{Q_\star T}{\gamma}\right)\right)$, and adjust Algorithm \ref{alg:clsf_regressionmain} (with $J^t,M^t \equiv 0$) to run unconstrained with the Azoury-Vovk-Warmuth forecaster for the square loss. Then, the predictions $\hat{y}_1, \ldots, \hat{y}_T$ satisfy
    $$\sum_{t=1}^T \|\hat{y}_t - y_t\|^2 \leq O\left(\log^2(Q_\star) \cdot d_\mathcal{X}R^2 \log(RT) \log(T)\cdot \sqrt{T}\right).$$
    
\end{theorem}
\begin{proof}[Proof sketch]
With this choice of parameters, after constructing the discretized LDS according to Lemma \ref{lemma:discretization}, the approximation bound from the proof of Lemma \ref{lemma:spectralfiltering} says that there exists a choice of parameters for which the summed squared losses are bounded by $O\left(\log^2(Q_\star) \cdot d_\mathcal{X} R^2 \log(RT) \sqrt{T}\right)$. The regret bound from equation (3.12) of \cite{azoury2001relative} adds $O\left(R^2 \log((m+h)Q_\star T)\right)$ from the optimization.
\end{proof}
The above result tells roughly the same story as our main Theorem \ref{theorem:main_nonlinear}: as long as only a small portion of the Koopman eigenvalues are undesirable, spectral filtering has a strong regret guarantee. This statement of the result trades poor dependence on $Q_\star$ in the regret bound for an increase in the number of parameters.
\section{Multi-dimensional Observations} \label{sec:higher-d-y}
In this section, we discuss what changes when allowing higher $d_\mathcal{Y}$. To begin, in such settings we can have $(A, C)$ being observable even if $A$ has an eigenvalue with multiplicity $d_\mathcal{Y}$, and so eigenvalue placement will allow us to move undesirable eigenvalues with multiplicity. In addition, while Ackermann's formula as stated no longer holds, the bound 
$$\norm{L_\star} \leq \frac{2Q_\star}{\sigma_{\text{min}}(C)}$$
of Lemma \ref{lemma:ackermannbounds}
continues to hold. The regret guarantee of Lemma \ref{lemma:spectralfiltering} depends on $\|L_\star\|_F \cdot \|C\|_F$ in a sharp way, and we use that $\|L_\star\|_F = \|L_\star\|$ and $\|C\|_F = \|C\| = \sigma_{\text{min}}(C)$ in the 1-dimensional observation case in order to simplify several expressions. While in most places our bounds will simply introduce a $\sqrt{d_\mathcal{Y}}$ factor gotten from converting between operator and Frobenius norms, we will also see a $\frac{\|C\|}{\sigma_{\text{min}}(C)}$ factor in the regret bounds of Theorems \ref{theorem:main_nonlinear} and \ref{theorem:main_linear}.

To summarize, increasing the observation dimension (1) strengthens eigenvalue placement, (2) adds $\sqrt{d_\mathcal{Y}}$ factors to many of the norms, and (3) creates a difference between the max and min singular values of $C$, which will show up in the $\|L_\star\|_F \cdot \|C\|_F$ terms in the regret bounds. It is difficult to get control on the conditioning of $C$ in our Markov chain construction, but morally everything stated remains true.

\section{Theoretical Subtleties}
In this section, we discuss a few nuances in our theory. Some appear paradoxical at first glance, but with more care they yield useful information and intuition about these proof techniques.

\paragraph{Strong stability.}
To begin, since we can construct the observer system with whichever poles we want, one could ask "why not just make $A-LC$ to be $\gamma$-stable for $\gamma = \frac{1}{2}$ and use $\widetilde{\Theta}\left(\frac{1}{\gamma}\right)$ regression terms?". The answer, as explored a bit in Remark \ref{remark:regression}, is that the number of autoregressive terms needed also scales with the spectral condition number of the resulting observer system, i.e. one would actually need $\widetilde{\Theta}\left(\frac{\log Q_\star}{\gamma}\right)$ many regression terms. There is a price to pay for more extreme eigenvalue placement given by $Q_\star$, and the cost of a fully stable observer system is seen by the number of regression terms needed to learn it. To drive this point home, if we consider a symmetric system with $N$ many marginally stable eigenvalues, regression would require $N$ parameters, whereas spectral filtering can get the job done independently of $N$.

\paragraph{Use of observations.}
Next, since we use pole placement for two purposes (to change the system eigenstructure and to construct an LDS over observations against which to apply spectral filtering), one could ask "what value of $L$ should be taken if the system already has desirable spectral structure?". For example, if we start with a symmetric autonomous LDS $x_{t+1} = Ax_t$ and $y_t = Cx_t$, there is no immediate input-output sequence-to-sequence mapping against which spectral filtering should hope to compete against. The classical way to enforce real-diagonalizability and marginal stability of $A-LC$ in this case would be to choose $L = 0$, but this does not fix the problem that there is no natural $[1, \alpha, \alpha^2, \ldots, \alpha^T]$ structure to an input-output mapping. The solution is to recognize that spectral filtering competes against \textit{all} observer systems -- for the choice of $L = 0$ we pay the full price of the initial state, and if $A$ is marginally stable we cannot hope for any decay. However, one can choose $L$ nonzero but such that $A- LC$ remains desirable (for example, if $A$ has an eigenvalue strictly less than 1 we can move it and only it, whereas if $A$ has all eigenvalues $1$ then we can take $L$ proportional to $C^\top$). In this way, we construct an observer LDS with a nontrivial input-output mapping which retains the desirable spectral structure, and spectral filtering can efficiently learn to represent this. 

% \paragraph{Spectral gap assumption.}




% \newpage
% \input{older-analysis}
% \input{older-stuff}
% \newpage
% \input{A7-symmetrickoopman}
% \newpage
% %\evannote{placeholder}
% \newpage
% \input{A5-old-koopman}
\ignore{

\section{Time-symmetric dynamical systems}\label{app:timesymm}
\begin{theorem}[Symmetric lifting for time-symmetric dynamics]\label{thm:symm-lift}
Let $(\mathcal{X},f)$ be a deterministic, discrete-time, time-symmetric system $x_{t+1}=f(x_t)$ with observation $y_t=h(x_t)$.
Assume:
\begin{itemize}
\item[(A1)] (\emph{Boundedness}) $\|x_t\|\le R$ and $\|y_t\|\le R$ for all $t$.
\item[(A2)] (\emph{Time symmetry}) There is an involution $R:\mathcal{X}\to\mathcal{X}$ with $f^{-1}=R\circ f\circ R$.
\item[(A3)] (\emph{Lipschitz}) $f$ and $h$ are $1$-Lipschitz (w.r.t.\ a common norm).\footnote{Or equip the product space $\mathcal{Y}\times\mathcal{X}$ with the max norm; see the proof.}
\end{itemize}
Then for any horizon $T$ and any $\varepsilon>0$, there exist:
\begin{itemize}
\item a finite-dimensional lifted state $z_t\in\mathbb{R}^N$ and a \emph{symmetric} matrix $\mathsf{S}\in\mathbb{R}^{N\times N}$ with spectrum in $[-1,1]$,
\item a linear readout $C:\mathbb{R}^N\to\mathcal{Y}$ and an initialization $z_{-1},z_0$,
\end{itemize}
such that the \emph{symmetric two-step LDS}
\[
z_{t+1}=2\,\mathsf{S}\,z_t-z_{t-1},\qquad \hat y_t=C z_t,
\]
approximates the true outputs with
\[
\sup_{0\le t\le T}\,\|\hat y_t-y_t\|\ \le\ (T+1)\,\varepsilon.
\]
Equivalently, writing $\xi_t=\begin{bmatrix}z_t\\ z_{t-1}\end{bmatrix}$ yields a first-order LDS
$\xi_{t+1}=\mathsf{A}\,\xi_t$ with
$\mathsf{A}=\begin{bmatrix}2\mathsf{S}&-I\\ I&0\end{bmatrix}$ and output $\hat y_t=[\,C\ \ 0\,]\xi_t$.
\end{theorem}

\begin{proof}
\textbf{Step 1: Discretize the joint space with a 1-Lipschitz update.}
Let $\mathcal{Z}:=\mathcal{Y}\times\mathcal{X}$ with norm $\|(y,x)\|_{\max}:=\max\{\|y\|,\|x\|\}$.
Define $F:\mathcal{Z}\to\mathcal{Z}$ by $F(y,x)=(h(f(x)),\,f(x))$.
Since $f,h$ are $1$-Lipschitz, for all $(y,x),(y',x')$,
\[
\|F(y,x)-F(y',x')\|_{\max}=\max\{\|h(f(x))-h(f(x'))\|,\ \|f(x)-f(x')\|\}\le \|x-x'\|\le \|(y,x)-(y',x')\|_{\max}.
\]
Thus $F$ is $1$-Lipschitz on $(\mathcal{Z},\|\cdot\|_{\max})$.
By (A1), $(y_t,x_t)$ stays in the radius-$R$ ball in $\mathcal{Z}$.
Fix $\varepsilon>0$, let $S\subset\mathcal{Z}$ be an $\varepsilon$-net of that ball, and let $\pi:\mathcal{Z}\to S$ be nearest neighbor.
Define the \emph{finite} deterministic map $T_S:S\to S$ by $T_S(s):=\pi(F(s))$.

Write $s_t:=T_S^t(s_0)$ where $s_0:=\pi(y_0,x_0)$, and denote $s^{(j)}=(y^{(j)},x^{(j)})\in S$.
Let $U\in\{0,1\}^{|S|\times|S|}$ be the adjacency matrix of $T_S$: $U_{ij}=1$ iff $T_S(s^{(j)})=s^{(i)}$.
Define the discretized outputs $y'_t:=\mathrm{proj}_{\mathcal{Y}}(s_t)$ and the linear map $C':\mathbb{R}^{|S|}\to\mathcal{Y}$ by $C'e_j=y^{(j)}$.

By $1$-Lipschitzness of $F$ and the net property, the standard induction gives
\[
e_t:=\|s_t-(y_t,x_t)\|_{\max}\ \le\ (t+1)\varepsilon,\qquad\Rightarrow\qquad \|y'_t-y_t\|\le e_t\le (t+1)\varepsilon\quad (0\le t\le T).
\tag{$\star$}
\]

\textbf{Step 2: Lift $T_S$ to a permutation that preserves outputs.}
For $v\in S$ set $\mathrm{Pre}(v):=\{u\in S: T_S(u)=v\}$, $k(v):=|\mathrm{Pre}(v)|$, and $k^+(v):=\max\{1,k(v)\}$.
Fix for each $w\in S$ a bijection $b_w:\{1,\dots,k(w)\}\to \mathrm{Pre}(w)$.
For each $v\in S$ let $j_v\in\{1,\dots,k(T_S(v))\}$ be the index with $b_{T_S(v)}(j_v)=v$.

Define the lifted state set $\tilde S:=\{(v,i): v\in S,\ i\in[k^+(v)]\}$ and the map $\Pi:\tilde S\to\tilde S$ by
\[
\Pi(v,i)\ :=\ (T_S(v),\,j_v).
\]
\emph{Claim:} $\Pi$ is a bijection. Injectivity: if $\Pi(v,i)=\Pi(v',i')$, then $T_S(v)=T_S(v')=w$ and $j_v=j_{v'}$, hence $b_w(j_v)=v=v'$. Surjectivity: given $(w,j)$, set $v:=b_w(j)$ and take any valid $i$; then $\Pi(v,i)=(w,j)$.

Define linear maps $R:\mathbb{R}^{|\tilde S|}\to\mathbb{R}^{|S|}$ and $E:\mathbb{R}^{|S|}\to\mathbb{R}^{|\tilde S|}$ by
\[
(R\tilde z)_v=\sum_{i=1}^{k^+(v)}\tilde z_{(v,i)},\qquad
E e_v=\begin{cases}\frac{1}{k(v)}\sum_{i=1}^{k(v)} e_{(v,i)},&k(v)\ge1,\\[3pt] e_{(v,1)},&k(v)=0.\end{cases}
\]
Then $RE=I$ and, for each basis vector $e_v$,
\[
R\,\Pi\,E\,e_v\ =\ R\,e_{(T_S(v),j_v)}\ =\ e_{T_S(v)}\ =\ Ue_v,
\]
so $R\Pi E=U$. Let $\tilde C:=C'R$ and initialize $\tilde z_0:=E z_0$ (where $z_0$ is the one-hot vector of $s_0$).
Then for all $t$, $\tilde z_t:=\Pi^t\tilde z_0$ satisfies
\[
\tilde y_t\ :=\ \tilde C\,\tilde z_t\ =\ C'R\,\Pi^t E z_0\ =\ C' U^t z_0\ =\ y'_t.
\tag{$\diamond$}
\]

\textbf{Step 3: Make the linear dynamics symmetric via the cosine operator.}
Since $\Pi$ is a permutation matrix, $\Pi^{-1}=\Pi^\top$.
Define the \emph{symmetric} cosine operator
\[
\mathsf{S}\ :=\ \tfrac12\big(\Pi+\Pi^\top\big),
\]
which is symmetric with spectrum in $[-1,1]$.
The lifted states obey the exact two-step identity
\[
\tilde z_{t+1}+ \tilde z_{t-1}=(\Pi+\Pi^{-1})\tilde z_t=2\,\mathsf{S}\,\tilde z_t
\quad\Longleftrightarrow\quad
\tilde z_{t+1}=2\,\mathsf{S}\,\tilde z_t-\tilde z_{t-1}.
\]
With $C:=\tilde C$ and $z_t:=\tilde z_t$, define $\hat y_t:=C z_t$.
By $(\diamond)$ and $(\star)$,
\[
\sup_{0\le t\le T}\|\hat y_t-y_t\|\ =\ \sup_{0\le t\le T}\|y'_t-y_t\|\ \le\ (T+1)\,\varepsilon.
\]
Finally, stacking $\xi_t=\begin{bmatrix}z_t\\ z_{t-1}\end{bmatrix}$ gives the equivalent first-order LDS
$\xi_{t+1}=\begin{bmatrix}2\mathsf{S}&-I\\ I&0\end{bmatrix}\xi_t$, with readout $[\,C\ \ 0\,]$.
\end{proof}

\begin{remark}[On the role of time symmetry]
Assumption (A2) motivates the use of a symmetric comparator: for invertible time-symmetric dynamics the Koopman
operator is unitary on an invariant $L^2$ space, and its cosine $\tfrac12(U+U^{-1})$ is self-adjoint with spectrum in $[-1,1]$.
Our finite construction mirrors this exactly via the permutation $\Pi$ and the symmetric $\mathsf{S}=\tfrac12(\Pi+\Pi^\top)$.
\end{remark}

}






%%%%%%%%%%%%%%%%%%%%%%%%%%%%%%%%%%%%%%%%%%%%%%%%%%%%%%%%%%%%%%5
%%%%%%%%%%%%%%%%%%%%%%%%%%%%%%%%%%%%%%%%%%%%%%%%%%%%%%%%%%%%%%5
%% original attempt
%%%%%%%%%%%%%%%%%%%%%%%%%%%%%%%%%%%%%%%%%%%%%%%%%%%%%%%%%%%%%%5
%%%%%%%%%%%%%%%%%%%%%%%%%%%%%%%%%%%%%%%%%%%%%%%%%%%%%%%%%%%%%%5


\ignore{
\section{Time-symmetric dynamical systems} \label{sec:symmetric-koopman}
\begin{theorem}
Let $(\mathcal{X}, f)$ be a deterministic, discrete-time, time-symmetric dynamical system: $x_{t+1} = f(x_t)$, with observation $y_t = h(x_t)$, where $h: \mathcal{X} \to \mathcal{Y}$. Assume:
\begin{itemize}
    \item[\textbf{(A1)}] (\textbf{Boundedness}) All trajectories remain in a ball of radius $R$: $\|x_t\| \leq R, \|y_t\| \leq R$ for all $t$.
    \item[\textbf{(A2)}] (\textbf{Time symmetry}) There exists an involutive map $R: \mathcal{X} \to \mathcal{X}$ such that $f^{-1} = R \circ f \circ R$.
    \item[\textbf{(A3)}] (\textbf{Lipschitz regularity}) $f$ and $h$ are Lipschitz.
\end{itemize}

Let $\epsilon > 0$. Then for any time horizon $T$, there exists a symmetric linear operator $\tilde{K}_f$ and a linear projection $P$ to the observation space, such that the sequence of outputs $\hat{y}_t = P z_t$ (where $z_{t+1} = \tilde{K}_f z_t$ and $z_0$ corresponds to the initial state) satisfies
\[
\sup_{0 \leq t \leq T} \|\hat{y}_t - y_t\| \leq T\epsilon,
\]
where $y_t = h(x_t)$ are the true system outputs.
\end{theorem}

\begin{proof}
Let $\mathcal{Z} = \mathcal{Y} \times \mathcal{X}$, and consider all $(y, x)$ with $\|y\| \leq R$ and $\|x\| \leq R$. By (A1), the true trajectory $(y_t, x_t)$ always remains in this set.

\textbf{Step 1: Discretization.}  
Choose an $\epsilon$-net $S$ of $\mathcal{Z}$ (with respect to the Euclidean norm), so that for every $(y, x) \in \mathcal{Z}$, there exists $s \in S$ with $\|(y, x) - s\| \leq \epsilon$.  
Let $\pi: \mathcal{Z} \to S$ denote the nearest-neighbor projection.

\textbf{Step 2: Discrete Dynamics Operator.}  
Define the map
\[
F: \mathcal{Z} \to \mathcal{Z}, \qquad F(y, x) := (h(f(x)), f(x)).
\]
For $s = (y, x) \in S$, define $T_S(s) := \pi(F(y, x))$ (i.e., apply $f$ and $h$, then project to the nearest element in $S$).

Let $K_f$ be the $|S| \times |S|$ matrix with $(K_f)_{ij} = 1$ if $T_S(s^{(j)}) = s^{(i)}$, and $0$ otherwise.

\textbf{Step 3: Symmetrization.}  
Define the symmetric operator $\tilde{K}_f = \frac{1}{2}(K_f + K_f^\top)$. By construction, $\tilde{K}_f$ is symmetric.

\textbf{Step 4: Evolution and Output.}  
Let $z_0 \in \mathbb{R}^{|S|}$ be the indicator vector for $s_0 = \pi(y_0, x_0)$.  
Iteratively compute $z_{t+1} = \tilde{K}_f z_t$.

Define the observation output $\hat{y}_t = P z_t$, where $P: \mathbb{R}^{|S|} \to \mathcal{Y}$ is the linear map such that $P e_j = y^{(j)}$ if $s^{(j)} = (y^{(j)}, x^{(j)})$ (i.e., the $y$-component of each $s^{(j)}$).

\textbf{Step 5: Error Propagation.}  
Let $(y_t, x_t)$ be the true system trajectory, and $s_t$ the representative discretized state at time $t$.

Define the error $e_t := \|s_t - (y_t, x_t)\|$.  
By construction, $e_0 \leq \epsilon$.

Suppose $e_t \leq t\epsilon$. Then
\begin{align*}
e_{t+1} &= \| s_{t+1} - (y_{t+1}, x_{t+1}) \| \\
&= \| \pi(F(s_t)) - F(y_t, x_t) \| \\
&\leq \| F(s_t) - F(y_t, x_t) \| + \| \pi(F(s_t)) - F(s_t) \| \\
&\leq L_F \| s_t - (y_t, x_t) \| + \epsilon
\end{align*}
where $L_F$ is the Lipschitz constant of $F$.

\textbf{Now, where is the symmetry used?}

Because the system is \emph{time-symmetric}, the transition operator $K_f$ and its transpose $K_f^\top$ correspond to forward and backward evolution, and are close in operator norm for sufficiently fine discretization.  
Thus, the action of their average $\tilde{K}_f$ preserves the $\epsilon$-close approximation to the true trajectory:  
- The symmetrization does not increase the per-step error compared to $K_f$ alone, because $(K_f^\top z_t)$ can be interpreted as the time-reversed step, and their average is within $O(\epsilon)$ of the true transition for reversible systems.

Thus, the total error at each step can be bounded by $e_{t+1} \leq e_t + \epsilon$, just as in the purely forward case.

By induction, $e_t \leq t \epsilon$.

\textbf{Step 6: Output Error.}  
The projection $P$ is $1$-Lipschitz (since it is coordinate projection), so
\[
\|\hat{y}_t - y_t\| = \|P(s_t) - y_t\| \leq \|s_t - (y_t, x_t)\| = e_t \leq t\epsilon.
\]
So
\[
\sup_{0 \leq t \leq T} \|\hat{y}_t - y_t\| \leq T\epsilon.
\]

\textbf{Symmetry:}  
$\tilde{K}_f$ is symmetric by construction, and the argument above shows its use: time symmetry ensures the average of forward and backward transition operators remains an accurate approximation (at the scale of the discretization) for observable evolution.

\end{proof}
}

\section{Nonlinear Systems with Stochastic Noise}
 \label{sec:nl-w-noise}
 We have discussed deterministic systems of the form $x_{t+1} = f(x_t)$ with $x_0$ fixed. This yields a Koopman operator sending $h(x) \mapsto h(f(x))$ and and a Markov chain which tracks the deterministic transitions. Both of these notions can be extended seamlessly to the case of stochastic systems of the form
 $$x_{t+1} = f(x_t) + w_t, \quad y_t = h(x_t), \quad x_0 \sim \mu_{x_0}$$
 for an initial condition drawn from $\mu_{x_0}$, where
 $w_t \sim \mu_w$ are i.i.d. stochastic noises. For example, applying the Koopman operator $n$ times sends the function $h(x) \mapsto \mathbb{E}[f(x_n) \mid x_0 = x]$, and the Markov chain simply gets non-integer entries that record the stochastic transition probabilities. As long as we assume that $\|x_t\|, \|y_t\| \leq R$ almost surely, the discretization argument still holds exactly as before, where now $y'_t = C'z_t$ approximates $\mathbb{E}[y_t]$. Note that we need not assume anything about the distributions $\mu_w, \mu_{x_0}$ themselves -- boundedness of states and observations is enough. The above reasoning is sufficient to prove:
 \begin{corollary}
     Let $x_{t+1} = f(x_t) + w_t$ be a stochastic dynamical system with observation function $h$ producing a sequence $y_1, \ldots, y_T$ of observations, where $w_t \sim \mu_w$ i.i.d. and $x_0 \sim \mu_{x_0}$. Assume that for all $t$, $\|x_t\|, \|y_t\| \leq R$ almost surely, $f, h$ are 1-Lipschitz, and the discretized linear lifting $(A', C')$ from Lemma \ref{lemma:discretization} with stochastic transitions is an observable LDS. 

     Fix the number of algorithm parameters $h = \Theta(\log^2 T)$ and $m \geq 1$, and for notation define $\gamma = \frac{\log(mT)}{m}$. Let $Q_\star = Q_\star\left(A', C', \left[0, 1 - \frac{1}{\sqrt{T}}\right] \cup \{1\} \cup \mathbb{D}_{1-\gamma}\right)$ as specified in Definition \ref{def:luenbergerprogram}. Then, running Algorithm \ref{alg:clsf_regressionmain} with $J^t, M^t \equiv 0$, $\mathcal{K} = R_D$ for $D = \Theta((m+h)Q_\star^2)$, and $\eta_t = \Theta\left(\frac{Q_\star^2}{R\sqrt{t} \log T}\right)$ produces a sequence of predictions $\hat{y}_1, \ldots, \hat{y}_T$ for which
     $$\mathbb{E}\sum_{t=1}^T \|\hat{y}_t - \mathbb{E} y_t\| \leq O\left(Q_\star^2 \log(Q_\star) \cdot d_\mathcal{X} R \log(RT) m^2 \log^7(mT)  \cdot \sqrt{T}\right). $$
 \end{corollary}
 \begin{proof}
     The proof goes the same way as that of Theorem \ref{theorem:main_nonlinear}. We form the discretized Markov chain $A', C'$, where now the transition matrix is given by the transition probabilities from one grid point to another after quantization. More precisely, we have
     $$A'_{ij}\;:= \mathbb{P}_{w \sim \mu_w }\left(\pi\left(f(s^{(j)})+w\right)=s^{(i)}\right).$$
     The LDS is then
     $$z_{t+1} = A'z_t, \quad y'_t = C'z_t, \quad z_0 = \pi_\# \mu_{x_0} ,$$
     where $\pi$ is the projection to the grid (i.e. $z_0$ is the probability vector for the location of the initial state). Since each $z_t$ is the law of the grid state of the discretized system, we see that $y_t' = \mathbb{E}[(h \circ \underbrace{ T \circ \ldots \circ T}_{\text{$t$ times}} \circ \pi)(x_0)]$ in the language of Lemma \ref{lemma:discretization}. Applying the same Lipschitz argument allows us to bound the difference between this discretized rollout and the true trajectory almost surely along the whole path, which implies the same bound on the size of the difference between their expectations. From here on we have again reduced to competing against the predictions of a deterministic LDS. Note that by the Lipschitzness argument, since $f(x_t) + w_t$ is a.s. bounded by $R$ by assumption, we will have $f(s_t) + w_t$ bounded by $R + t\epsilon \leq R + \frac{1}{\sqrt{T}}$ always. This means that under the assumptions we imposed, we can use an $\epsilon$-net of the ball of size $O(R)$ regardless of how large $w_t$ is. 

     The final subtlety is that our predictions $\hat{y}_{t+1}$ are in fact random, since we filter over $y_1, \ldots, y_t$ instead of the deterministic $y_t'$ which is needed to compete against the Markov chain LDS. This introduces an additional error of size $\|L\| \cdot \|y_t - y_t'\|$, which gets upper bounded by $Q_\star \cdot t\epsilon$ almost surely. This gets absorbed into our other terms.
 \end{proof}

\ignore{
 \newpage

\section{Nonlinear dynamical systems with noise} 
% \label{sec:nl-w-noise}

\begin{theorem}[Discretization with bounded noise]
Let $x_{t+1}=f(x_t)+w_t$ with $f,h$ $1$-Lipschitz. Suppose $\|x_t\|\!,\|y_t\|\le R$ for all $t$ and
the noise is a.s.\ bounded: $\|w_t\|\le W$. Fix horizon $T$ and $\varepsilon>0$.

Let $R':=R+2W+T\varepsilon$, and let $S$ be an $\varepsilon$-net of the ball of radius $R'$.
Let $\pi:X\to S$ be nearest-neighbor. For each $s^{(j)}\in S$ define
\[
P_{ij}\;:=\;\Pr\!\big(\,\pi\big(f(s^{(j)})+w\big)=s^{(i)}\,\big),\qquad
z_{t+1}=Pz_t,\quad \hat y_t=Cz_t,\ \ C e_j=h\big(s^{(j)}\big).
\]
Define the \emph{coupled} discretized process
\[
s_0=\pi(x_0),\qquad s_{t+1}=\pi\big(f(s_t)+w_t\big),\qquad y'_t:=h(s_t).
\]
Then for all $t\le T$:
\begin{align}
\|\,\hat y_t-\mathbb{E}[y_t]\,\| &\le t\varepsilon, \label{eq:exp-gap}\\
\|\,y'_t-y_t\,\| &\le t\varepsilon \quad\text{(pathwise; hence also }\mathbb{E}\|y'_t-y_t\|\le t\varepsilon\text{)}. \label{eq:path-gap}
\end{align}
\end{theorem}

\begin{proof}
Let $e_t:=\|s_t-x_t\|$. Note $e_0\le\varepsilon$ by construction.
Since $x_{t+1}=f(x_t)+w_t$ and $s_{t+1}=\pi(f(s_t)+w_t)$ with the \emph{same} $w_t$,
\[
e_{t+1}\;=\;\big\|\pi\!\big(f(s_t)+w_t\big)-\big(f(x_t)+w_t\big)\big\|
\;\le\;\underbrace{\big\|\pi\!\big(f(s_t)+w_t\big)-\big(f(s_t)+w_t\big)\big\|}_{\le\varepsilon}\;+\;\underbrace{\|f(s_t)-f(x_t)\|}_{\le e_t}
\;\le\;e_t+\varepsilon.
\]
Thus $e_t\le t\varepsilon$ and, by $1$-Lipschitzness of $h$, $\|y'_t-y_t\|\le e_t\le t\varepsilon$, proving \eqref{eq:path-gap}.
Moreover,
\[
\|\hat y_t-\mathbb{E}[y_t]\|
=\big\|\mathbb{E}[\,h(s_t)\,]-\mathbb{E}[\,h(x_t)\,]\big\|
\le \mathbb{E}\|h(s_t)-h(x_t)\|
\le \mathbb{E}e_t
\le t\varepsilon,
\]
which gives \eqref{eq:exp-gap}.

Finally, the $\varepsilon$-net radius $R'$ guarantees coverage: by induction,
$\|s_t\|\le \|x_t\|+e_t \le R+ t\varepsilon\le R'$ and
$\|f(s_t)+w_t\|\le \|f(x_t)\|+\|f(s_t)-f(x_t)\|+\|w_t\|\le (R+W)+e_t+W\le R+2W+T\varepsilon=R'$,
so the projection error $\le\varepsilon$ is valid at every step.
\end{proof}


\begin{remark}[On $R'$ depending on $\varepsilon$]
We first fix the discretization scale $\varepsilon$, then set $R'(\varepsilon):=R+2W+T\varepsilon$ and net the
ball $B_{R'(\varepsilon)}$. The induction uses the hypothesis $e_t\le t\varepsilon$ to ensure
$\|f(s_t)+w_t\|\le R+2W+T\varepsilon$, so the projection error is $\le\varepsilon$ and $e_{t+1}\le e_t+\varepsilon$.
\end{remark}


\paragraph{Implication for the main results.}
Let $(P,C)$ be the Markov-chain LDS constructed above. Then
\[
\underbrace{\sum_{t=1}^T\| \hat y_t - y_t\|}_{\text{algorithm on real data}}
\;\le\;
\underbrace{\sum_{t=1}^T\| \hat y_t - \hat y^{\,\star}_t\|}_{\text{regret vs.\ best comparator in }\mathrm{LDS}(P,C)}
\;+\;\underbrace{\sum_{t=1}^T\| \hat y^{\,\star}_t - \mathbb{E}[y_t]\|}_{\text{comparator mismatch}}
\;+\;\underbrace{\sum_{t=1}^T\| \mathbb{E}[y_t]-y_t\|}_{\text{noise \& coupling}},
\]
where $\hat y^{\,\star}$ is the best comparator sequence in your linear class.
The second bracket is controlled by \eqref{eq:exp-gap}, and the third by \eqref{eq:path-gap} in expectation (or pathwise a.s.\ under bounded noise).
Thus, Theorem \ref{theorem:main_linear} (c) applies to the comparator over $(P,C)$, and the discretization contributes only the additive
$O(T\varepsilon)$ terms (choose $\varepsilon=T^{-3/2}$ as in Lemma \ref{lemma:discretization}).%



\ignore{

\begin{theorem}[Discrete Stochastic Koopman Operator Approximates Observations (Discretization Error Only)]
Let $(\mathcal{X}, f)$ be a discrete-time nonlinear dynamical system with additive noise: $x_{t+1} = f(x_t) + w_t$, where $x_t \in \mathcal{X} \subset \mathbb{R}^d$, $w_t$ are i.i.d. random vectors (arbitrary distribution), and $f: \mathcal{X} \to \mathcal{X}$ is $1$-Lipschitz. Let $y_t = h(x_t)$, where $h: \mathcal{X} \to \mathcal{Y} \subset \mathbb{R}^m$ is $1$-Lipschitz.

Assume all states and outputs remain in a ball of radius $R$ for all $t$.

Fix $\epsilon > 0$ and time horizon $T$. Let $S$ be an $\epsilon$-net of the ball of radius $R$ in $\mathbb{R}^d$, with $|S| = N$. Let $\pi: \mathcal{X} \to S$ be the nearest neighbor projection. For each $s^{(j)} \in S$, define transition probabilities
\[
P_{ij} = \mathbb{P}\left( \pi(f(s^{(j)}) + w) = s^{(i)} \right),
\]
where $w$ is distributed as $w_t$. Let $P$ be the $N \times N$ transition matrix.

Let $z_0 \in \mathbb{R}^N$ be the indicator vector of $s_0 = \pi(x_0)$. Define $z_{t+1} = P z_t$. Define $C \in \mathbb{R}^{m \times N}$ by $C e_j = h(s^{(j)})$.

Let $\hat{y}_t = C z_t$ be the expected output from the Markov chain at time $t$.

Then for any $t \leq T$, the expected error satisfies
\[
\mathbb{E} \left[ \|\hat{y}_t - \mathbb{E}[y_t] \| \right] \leq t\epsilon,
\]
where the expectation on the right is over the noise in the true system.
\end{theorem}

\begin{proof}
Let $x_0$ be the initial state, and $s_0 = \pi(x_0)$. The initial error is $e_0 \leq \epsilon$.

At each time $t$, let $x_t$ be the state of the true system, and $s_t$ be the random variable for the Markov chain (with law $z_t$). Consider the expected error $e_t := \mathbb{E}[ \| s_t - x_t \| ]$.

For one step,
\[
x_{t+1} = f(x_t) + w_t, \qquad s_{t+1} = \pi( f(s_t) + w_t )
\]
(with the same noise used for both).

By triangle inequality and $1$-Lipschitz property,
\begin{align*}
e_{t+1} &= \mathbb{E}[ \| s_{t+1} - x_{t+1} \| ] \\
&= \mathbb{E}[ \| \pi(f(s_t) + w_t) - (f(x_t) + w_t) \| ] \\
&\leq \mathbb{E}[ \| \pi(f(s_t) + w_t) - (f(s_t) + w_t) \| ] + \mathbb{E}[ \| f(s_t) + w_t - f(x_t) - w_t \| ] \\
&\leq \epsilon + \mathbb{E}[ \| f(s_t) - f(x_t) \| ] \\
&\leq \epsilon + \mathbb{E}[ \| s_t - x_t \| ] \\
&= \epsilon + e_t.
\end{align*}
So, $e_{t+1} \leq e_t + \epsilon$, $e_0 \leq \epsilon$, so $e_t \leq t\epsilon$.

Now, since $h$ is $1$-Lipschitz, the expected output error is
\[
\mathbb{E}[ \| \hat{y}_t - \mathbb{E}[y_t] \| ] \leq e_t \leq t\epsilon.
\]
\end{proof}

}}

\end{document}

